% La familia : evolución y conflictos intergeneracionales — Espagnol 1ère A — Cours signé OnBuch+
% Programme officiel MINESEC. Compiler avec : tectonic EA01-familia-evolucion-conflictos.tex
\newcommand{\DOCMATIERE}{Espagnol}
\newcommand{\DOCNIVEAU}{1ère A}
\newcommand{\DOCMODULE}{Módulo 1 : Vie familiale et intégration sociale}
\newcommand{\DOCLECON}{EA01}
\newcommand{\DOCTITRE}{La familia : evolución y conflictos intergeneracionales}
% OnBuch+ — préambule « cours signés » ENRICHI (compilé avec tectonic / XeLaTeX)
% Système visuel « Pop » d'OnBuch+ : fond crème (jamais blanc), bordures encre
% épaisses, ombres « dures » décalées, titres Archivo Black en capitales.
% Version MATHÉMATIQUES : graphes (pgfplots/tikz), boîtes pédagogiques
% (définition, propriété, méthode, attention, le savais-tu, exercice résolu,
% exercice, TP, bilan), page de garde et signature OnBuch+ ancrée en bas.
%
% Macros attendues dans main.tex AVANT \input{preamble} :
%   \DOCMATIERE  \DOCNIVEAU  \DOCMODULE  \DOCLECON (numéro)  \DOCTITRE
\documentclass[11pt]{article}

\usepackage{fontspec}
\usepackage[spanish,french]{babel} % français = langue principale ; l'espagnol via \esp{..} / espagnol
\usepackage[a4paper,margin=2cm,top=3.4cm,bottom=2.8cm,headheight=1.4cm,footskip=1.3cm]{geometry}
\usepackage{amsmath,amssymb,amsfonts}
\usepackage{mathtools} % \xmapsto, \coloneqq... souvent utilisés par les modèles
\usepackage{cancel} % simplifications barrées (bilans de forces)
% \abs{z} (module, valeur absolue) et \norm{u} (norme), utilisés par les modèles
\providecommand{\abs}{}\let\abs\relax\DeclarePairedDelimiter{\abs}{\lvert}{\rvert}
\providecommand{\norm}{}\let\norm\relax\DeclarePairedDelimiter{\norm}{\lVert}{\rVert}
\usepackage[table]{xcolor} % [table] : fournit aussi \rowcolors (alternance de lignes)
\usepackage{pagecolor}
\usepackage{tikz}
\usetikzlibrary{arrows.meta,positioning,calc,shapes.geometric,shapes.misc,shapes.arrows,decorations.pathmorphing,decorations.pathreplacing,decorations.markings,patterns,shadows,mindmap,trees,fit,backgrounds,matrix,intersections,angles,quotes}
\usepackage{pgfplots}
\usepackage[version=4]{mhchem} % équations nucléaires \ce{^{235}_{92}U}
\usepackage[european,siunitx]{circuitikz} % schémas électriques
\pgfplotsset{compat=1.18}
\usepgfplotslibrary{fillbetween,groupplots} % aires entre courbes (calcul intégral)
\usepackage{enumitem}
\usepackage{fancyhdr}
% [most] charge toutes les bibliothèques tcolorbox (listings, minted, raster,
% documentation...) et gonfle inutilement la mémoire TeX sur un document long
% (~300 boîtes) : on ne charge que ce qui est réellement utilisé.
\usepackage[skins,breakable]{tcolorbox}
\usepackage{graphicx}
\usepackage{colortbl}
\usepackage{booktabs}
\usepackage{tabularx}
\usepackage{diagbox} % case d'en-tête diagonale des tableaux à double entrée
% Colonnes extensibles centrée / gauche / droite (C, L, R), souvent utilisées par les modèles
\newcolumntype{C}{>{\centering\arraybackslash}X}
\newcolumntype{L}{>{\raggedright\arraybackslash}X}
\newcolumntype{R}{>{\raggedleft\arraybackslash}X}
\usepackage{array}
\usepackage{multicol}
\usepackage{siunitx}
% parse-numbers=false : accepte aussi \SI{2,5 \times 10^{-3}}{...}
\sisetup{locale=FR,output-decimal-marker={,},group-minimum-digits=4,parse-numbers=false}
\DeclareSIUnit\ohm{\ensuremath{\symup{\Omega}}} % U+2126 absent de la police
\DeclareSIUnit\division{div} % réglages d'oscilloscope (ms/div, V/div)
\DeclareSIUnit\tr{tr} % tours (vitesse de rotation, ex. \SI{1500}{\tr\per\minute})
\DeclareSIUnit\molar{M} % concentration molaire (ex. \SI{3}{\molar}, \SI{1}{\micro\molar})
\DeclareSIUnit\an{an} % année (le modèle confond parfois avec \year, un compteur TeX) — ex. \SI{5}{\kilo\meter\per\an}
\DeclareSIUnit\FCFA{FCFA} % franc CFA (ex. \SI{500}{\FCFA\per\kilo\gram})
\DeclareSIUnit\degreeBrix{\textdegree Brix} % teneur en sucre d'un sirop/jus (réfractométrie)
\DeclareSIUnit\month{mois} % le modèle utilise parfois \month, un compteur TeX, comme unité de durée
\usepackage{qrcode}
% \degree : commande fréquemment utilisée par les modèles pour un angle ou une
% température en dehors de \SI{}{} (non fournie par les paquets chargés).
\providecommand{\degree}{^{\circ}}
% \qed (fin de démonstration), utilisé par les modèles sans amsthm
\providecommand{\qed}{\hfill$\square$}
% \barre : double barre des valeurs interdites dans un tableau de variations (tkz-tab)
\providecommand{\barre}{\ensuremath{\|}}
% environnement proof (amsthm non chargé) utilisé par les modèles
\ifdefined\proof\else\newenvironment{proof}[1][Démonstration]{\par\noindent\textit{#1.}\ }{\qed\par}\fi
\usepackage{lastpage}
\usepackage{hyperref}
\hypersetup{hidelinks}

% ============================================================
% Palette « Pop » OnBuch+ — mêmes hex que la classe P (plus_ui.dart)
% ============================================================
\definecolor{poporange}{HTML}{F59321}
\definecolor{popdark}{HTML}{E07A0C}
\definecolor{popcream}{HTML}{FFF6E8}
\definecolor{popink}{HTML}{1C1714}
\definecolor{poppurple}{HTML}{7A5AE0}
\definecolor{popgreen}{HTML}{2E9E6A}
\definecolor{poppink}{HTML}{E0457B}
\definecolor{popblue}{HTML}{2F7DE1}
\definecolor{popgold}{HTML}{C98A12}
\definecolor{popmuted}{HTML}{8A7F76}
\definecolor{popline}{HTML}{EADFD2}
\definecolor{popgray}{HTML}{8A7F76}
\colorlet{popgrey}{popgray}
% Alias tolérants (le modèle nomme parfois les couleurs différemment)
\colorlet{popwhite}{white}
\colorlet{popblack}{popink}
\colorlet{popred}{poppink}
\colorlet{popyellow}{popgold}
% Teintes claires (fonds de tableaux/encadrés) — le modèle nomme parfois
% les couleurs avec un suffixe L (ex. popblueL) sans les avoir définies.
\colorlet{poporangeL}{poporange!20}
\colorlet{popdarkL}{popdark!20}
\colorlet{poppurpleL}{poppurple!20}
\colorlet{popgreenL}{popgreen!20}
\colorlet{poppinkL}{poppink!20}
\colorlet{popblueL}{popblue!20}
\colorlet{popgoldL}{popgold!20}
\colorlet{popredL}{popred!20}
\colorlet{popgrayL}{popgray!20}
% Teintes claires pour fonds de tableaux / zones
\colorlet{popblueL}{popblue!10!white}
\colorlet{poporangeL}{poporange!14!white}
\colorlet{popgreenL}{popgreen!12!white}
\colorlet{poppurpleL}{poppurple!12!white}
\colorlet{poppinkL}{poppink!12!white}
\colorlet{popgoldL}{popgold!14!white}

\pagecolor{popcream}

% ============================================================
% Typographie
% ============================================================
\defaultfontfeatures{Ligatures=TeX}
\setmainfont{PlusJakartaSans-Medium.ttf}[Path=../../../fonts/,BoldFont=PlusJakartaSans-Bold.ttf]
% Maths en sans-serif (Fira Math) avec chiffres et lettres droites de Plus
% Jakarta, pour que les formules se fondent dans le texte.
\usepackage[math-style=ISO,bold-style=ISO]{unicode-math}
\setmathfont{FiraMath-Regular.otf}[Path=../../../fonts/]
\setmathfont{PlusJakartaSans-Medium.ttf}[Path=../../../fonts/,range={up/num,up/{Latin,latin},\mathup}]
\usepackage{icomma} % virgule décimale sans espace en mode math
% Caractères mathématiques Unicode absents des polices de texte (Plus Jakarta,
% Archivo Black) : rendus par leur équivalent mathématique, en texte comme en
% formule — sinon ils s'affichent en carré vide (ex. « Arithmétique dans ℤ »).
\usepackage{newunicodechar}
\newunicodechar{ℝ}{\ensuremath{\mathbb{R}}}
\newunicodechar{ℤ}{\ensuremath{\mathbb{Z}}}
\newunicodechar{ℂ}{\ensuremath{\mathbb{C}}}
\newunicodechar{ℕ}{\ensuremath{\mathbb{N}}}
\newunicodechar{ℚ}{\ensuremath{\mathbb{Q}}}
\newunicodechar{𝔻}{\ensuremath{\mathbb{D}}}
\newunicodechar{α}{\ensuremath{\alpha}}
\newunicodechar{β}{\ensuremath{\beta}}
\newunicodechar{γ}{\ensuremath{\gamma}}
\newunicodechar{δ}{\ensuremath{\delta}}
\newunicodechar{ε}{\ensuremath{\varepsilon}}
\newunicodechar{θ}{\ensuremath{\theta}}
\newunicodechar{λ}{\ensuremath{\lambda}}
\newunicodechar{μ}{\ensuremath{\mu}}
\newunicodechar{σ}{\ensuremath{\sigma}}
\newunicodechar{τ}{\ensuremath{\tau}}
\newunicodechar{φ}{\ensuremath{\varphi}}
\newunicodechar{ψ}{\ensuremath{\psi}}
\newunicodechar{ω}{\ensuremath{\omega}}
\newunicodechar{Σ}{\ensuremath{\Sigma}}
% majuscules produites par \MakeUppercase dans les titres (α-aminés -> Α)
\newunicodechar{Α}{\ensuremath{\alpha}}
\newunicodechar{Β}{\ensuremath{\beta}}
\newunicodechar{Γ}{\ensuremath{\Gamma}}
\newunicodechar{Δ}{\ensuremath{\Delta}}
\newunicodechar{Ε}{\ensuremath{\varepsilon}}
\newunicodechar{Θ}{\ensuremath{\Theta}}
\newunicodechar{Λ}{\ensuremath{\Lambda}}
\newunicodechar{Μ}{\ensuremath{\mu}}
\newunicodechar{Τ}{\ensuremath{\tau}}
\newunicodechar{Φ}{\ensuremath{\Phi}}
\newunicodechar{Ψ}{\ensuremath{\Psi}}
\newunicodechar{Ω}{\ensuremath{\Omega}}
\newunicodechar{↦}{\ensuremath{\mapsto}}
\newunicodechar{∈}{\ensuremath{\in}}
\newunicodechar{∉}{\ensuremath{\notin}}
\newunicodechar{∧}{\ensuremath{\wedge}}
\newunicodechar{⇔}{\ensuremath{\Leftrightarrow}}
\newunicodechar{⟺}{\ensuremath{\Longleftrightarrow}}
\newunicodechar{⇒}{\ensuremath{\Rightarrow}}
\newunicodechar{⟹}{\ensuremath{\Longrightarrow}}
\newunicodechar{∘}{\ensuremath{\circ}}
\newunicodechar{∪}{\ensuremath{\cup}}
\newunicodechar{∩}{\ensuremath{\cap}}
\newunicodechar{≡}{\ensuremath{\equiv}}
\newunicodechar{⊂}{\ensuremath{\subset}}
\newunicodechar{⊥}{\ensuremath{\perp}}
\newunicodechar{∃}{\ensuremath{\exists}}
\newunicodechar{∀}{\ensuremath{\forall}}
\newunicodechar{∛}{\ensuremath{\sqrt[3]{\,}}}
\newunicodechar{∜}{\ensuremath{\sqrt[4]{\,}}}
\newunicodechar{⃗}{\ensuremath{{}^{\to}}}
\newunicodechar{ⁿ}{\ifmmode^{n}\else\textsuperscript{\ensuremath{n}}\fi}
\newunicodechar{ˣ}{\ifmmode^{x}\else\textsuperscript{\ensuremath{x}}\fi}
\newunicodechar{ᵇ}{\ifmmode^{b}\else\textsuperscript{\ensuremath{b}}\fi}
\newunicodechar{ʳ}{\ifmmode^{r}\else\textsuperscript{\ensuremath{r}}\fi}
\newunicodechar{ᵘ}{\ifmmode^{u}\else\textsuperscript{\ensuremath{u}}\fi}
\newunicodechar{ʸ}{\ifmmode^{y}\else\textsuperscript{\ensuremath{y}}\fi}
\newunicodechar{ᵃ}{\ifmmode^{a}\else\textsuperscript{\ensuremath{a}}\fi}
\newunicodechar{ᵉ}{\ifmmode^{e}\else\textsuperscript{\ensuremath{e}}\fi}
\newunicodechar{ᵏ}{\ifmmode^{k}\else\textsuperscript{\ensuremath{k}}\fi}
\newunicodechar{ᵖ}{\ifmmode^{p}\else\textsuperscript{\ensuremath{p}}\fi}
\newunicodechar{ᴺ}{\ifmmode^{N}\else\textsuperscript{\ensuremath{N}}\fi}
\newunicodechar{ᶜ}{\ifmmode^{c}\else\textsuperscript{\ensuremath{c}}\fi}
\newunicodechar{ʰ}{\ifmmode^{h}\else\textsuperscript{\ensuremath{h}}\fi}
\newunicodechar{ᵠ}{\ifmmode^{q}\else\textsuperscript{\ensuremath{q}}\fi}
\newunicodechar{ₙ}{\ifmmode_{n}\else\textsubscript{\ensuremath{n}}\fi}
\newunicodechar{ₐ}{\ifmmode_{a}\else\textsubscript{\ensuremath{a}}\fi}
\newunicodechar{ᵢ}{\ifmmode_{i}\else\textsubscript{\ensuremath{i}}\fi}
\newunicodechar{ᵣ}{\ifmmode_{r}\else\textsubscript{\ensuremath{r}}\fi}
\newunicodechar{ₖ}{\ifmmode_{k}\else\textsubscript{\ensuremath{k}}\fi}
\newunicodechar{ᵦ}{\ifmmode_{\beta}\else\textsubscript{\ensuremath{\beta}}\fi}
\newunicodechar{ₓ}{\ifmmode_{x}\else\textsubscript{\ensuremath{x}}\fi}
\newfontfamily\popdisplay{ArchivoBlack-Regular.ttf}[Path=../../../fonts/]
\newfontfamily\poplogo{SpaceGrotesk-Bold.ttf}[Path=../../../fonts/]
\newfontfamily\popsemi{PlusJakartaSans-SemiBold.ttf}[Path=../../../fonts/]
\newfontfamily\popxbold{PlusJakartaSans-ExtraBold.ttf}[Path=../../../fonts/]
\newfontfamily\popxboldls{PlusJakartaSans-ExtraBold.ttf}[Path=../../../fonts/,LetterSpace=3.0]

\setlist[enumerate]{leftmargin=1.7em,itemsep=5pt,topsep=4pt}
\setlist[itemize]{leftmargin=1.7em,itemsep=4pt,topsep=4pt,label={\color{poporange}\textbullet}}
\setlength{\parindent}{0pt}
\setlength{\parskip}{5pt}
\renewcommand{\arraystretch}{1.25}

% pgfplots : style Pop par défaut
\pgfplotsset{
  every axis/.append style={axis line style={popink,line width=1pt},tick style={popink},
    grid style={popline,line width=0.5pt},label style={font=\small},tick label style={font=\footnotesize}},
  popcurve/.style={poporange,line width=1.8pt,no markers,smooth},
}
% Même style disponible dans un \draw TikZ simple (hors pgfplots)
\tikzset{popcurve/.style={poporange,line width=1.8pt}}
\tikzset{popgrid/.style={popline,very thin}}
\colorlet{popcurve}{poporange} % le nom du style est parfois utilisé comme couleur

% ============================================================
% Pill / squiggle
% ============================================================
\newcommand{\poppill}[3]{%
  \tikz[baseline=-0.55ex]\node[draw=popink,line width=1.2pt,rounded corners=2.6pt,%
    fill=#2,inner xsep=6.5pt,inner ysep=3pt]{%
    \popxboldls\fontsize{7}{7}\selectfont\color{#3}\MakeUppercase{#1}};%
}
\newcommand{\popsquiggle}[1]{%
  \tikz[baseline=-0.4ex]{
    \draw[#1,line width=2.6pt,line cap=round]
      plot [smooth] coordinates {(0,0.09) (0.34,0.01) (0.68,0.21) (1.02,0.01) (1.36,0.10)};
  }%
}
% Mot-clé surligné (marqueur orange sous le texte)
\newcommand{\cle}[1]{{\popxbold\color{popdark}#1}}

% ============================================================
% En-tête / pied
% ============================================================
\pagestyle{fancy}
\fancyhf{}
\renewcommand{\headrulewidth}{0pt}
\renewcommand{\footrulewidth}{0pt}
\lhead{\raisebox{-0.15ex}{\poplogo\fontsize{13}{13}\selectfont\color{popink}OnBuch\color{poporange}+}}
\rhead{\poppill{\DOCMATIERE\ \textbullet\ \DOCNIVEAU\ \textbullet\ Leçon \DOCLECON}{white}{popink}}
\lfoot{\popxbold\fontsize{7.6}{7.6}\selectfont\color{popmuted}\MakeUppercase{Cours signé OnBuch+}\ \textbullet\ {\popsemi\DOCTITRE}}
\rfoot{\tikz[baseline=-0.6ex]\node[rounded corners=8.5pt,fill=popink,minimum height=17pt,minimum width=17pt,inner xsep=5pt,inner sep=0]{\popxbold\fontsize{8}{8}\selectfont\color{popcream}\thepage\,/\,\pageref*{LastPage}};}

% ============================================================
% Titres
% ============================================================
\newcommand{\chaptitre}[2]{%
  \begin{tcolorbox}[enhanced,colback=poporange,colframe=popink,boxrule=2.6pt,arc=11pt,
      drop shadow=popink,left=15pt,right=15pt,top=12pt,bottom=11pt,before skip=3pt,after skip=18pt]
    \poppill{#2}{popink}{popcream}\par\vspace{9pt}
    {\raggedright\hyphenpenalty=10000\exhyphenpenalty=10000\popdisplay\fontsize{22}{24}\selectfont\color{popcream}\MakeUppercase{#1}\par}
  \end{tcolorbox}
}
% Section numérotée : pastille orange avec le numéro + titre Archivo
\newcounter{popsec}
\newcommand{\coursec}[1]{%
  \stepcounter{popsec}\setcounter{popsub}{0}%
  \par\vspace{16pt}\noindent
  \begin{minipage}{\linewidth}
  \tikz[baseline=-0.7ex]\node[draw=popink,line width=1.6pt,fill=poporange,rounded corners=4pt,minimum size=22pt,inner sep=0]{\popdisplay\fontsize{12}{12}\selectfont\color{popcream}\thepopsec};%
  \hspace{8pt}{\popdisplay\fontsize{15}{17}\selectfont\color{popink}\MakeUppercase{#1}}\par
  \vspace{4pt}\noindent{\color{poporange}\rule{36pt}{3.6pt}}
  \end{minipage}\par\nopagebreak\vspace{8pt}
}
\newcounter{popsub}
\newcommand{\courssub}[1]{%
  \stepcounter{popsub}%
  \par\vspace{9pt}\noindent{\popxbold\fontsize{11.5}{13}\selectfont\color{popink}{\color{poporange}\thepopsec.\thepopsub}\hspace{6pt}#1}\par\nopagebreak\vspace{4pt}
}

% ============================================================
% Boîtes pédagogiques — même recette : fond blanc, bordure épaisse colorée,
% ombre dure encre, badge plein. Titre optionnel [#1].
% ============================================================
\tcbset{popbase/.style={breakable,enhanced,colback=white,boxrule=2.3pt,arc=9pt,
  shadow={3pt}{-3pt}{0pt}{popink},left=14pt,right=14pt,top=11pt,bottom=11pt,before skip=11pt,after skip=15pt,
  fontupper=\popsemi\fontsize{10.6}{14.5}\selectfont\color{popink},
  fontlower=\popsemi\fontsize{10.6}{14.5}\selectfont\color{popink}}}
% Coche dessinée (le glyphe ✓ n'existe pas dans les polices du document)
\newcommand{\popcheck}[1]{\tikz[baseline=-0.5ex]\draw[#1,line width=1.6pt,line cap=round,line join=round](-0.13,0.0)--(-0.04,-0.09)--(0.13,0.1);}
\providecommand{\coche}{\popcheck{popgreen}}
\providecommand{\resultat}[1]{\par\smallskip\noindent\fcolorbox{popgreen}{popgreenL}{\parbox{\dimexpr\linewidth-2\fboxsep-2\fboxrule\relax}{\textbf{Résultat.} #1}}\par\smallskip}
\newcommand{\popboxhead}[3]{\poppill{#1}{#2}{white}\ifx\relax#3\relax\else\hspace{6pt}{\popxbold\fontsize{10.6}{12}\selectfont\color{popink}#3}\fi\par\vspace{7pt}}

\newtcolorbox{aretenir}[1][]{popbase,colframe=popgreen,before upper={\popboxhead{À retenir}{popgreen}{#1}}}
% Texte en espagnol (document de lecture, dialogue, poème, extrait) : cadre
% d'étude. Un titre optionnel (source, auteur) s'affiche dans l'en-tête.
\newtcolorbox{textebox}[1][]{popbase,colframe=popblue,before upper={\popboxhead{Texto}{popblue}{#1}\begin{otherlanguage}{spanish}},after upper={\end{otherlanguage}}}
% Marques ✓ / ✗ (forme correcte / fautive) souvent utilisées par les modèles
\providecommand{\cross}{\textcolor{poppink}{\ensuremath{\times}}}
\providecommand{\xmark}{\cross}\providecommand{\crossmark}{\cross}\providecommand{\cmark}{\ensuremath{\checkmark}}
% Ligne de réponse / justification à compléter (inventée par les modèles)
\providecommand{\justification}[1]{\par\noindent\textit{Justification~:} \dotfill\par}
% \textipa (package tipa non chargé) : la police ne couvre pas l'API -> simple passage
\providecommand{\textipa}[1]{#1}
% Soulignement (ulem non chargé) : \uline, \ul
\providecommand{\uline}[1]{\underline{#1}}\providecommand{\ul}[1]{\underline{#1}}
% Ordinaux français écrits en macro par les modèles : 1\ière, 3\ième
\newcommand{\ière}{\textsuperscript{re}}\newcommand{\ième}{\textsuperscript{e}}
% Mot ou phrase en espagnol à l'intérieur d'un paragraphe français
\newcommand{\esp}[1]{\ifmmode\text{\textit{#1}}\else\begingroup\selectlanguage{spanish}\itshape #1\endgroup\fi}
\newtcolorbox{exemplebox}[1][]{popbase,colframe=poporange,before upper={\popboxhead{Exemple}{poporange}{#1}}}
\newtcolorbox{definition}[1][]{popbase,colframe=popblue,before upper={\popboxhead{Définition}{popblue}{#1}}}
\newtcolorbox{propriete}[1][]{popbase,colframe=poppurple,before upper={\popboxhead{Propriété}{poppurple}{#1}}}
\newtcolorbox{methode}[1][]{popbase,colframe=popgold,before upper={\popboxhead{Méthode}{popgold}{#1}}}
\newtcolorbox{attention}[1][]{popbase,colframe=poppink,before upper={\popboxhead{Attention}{poppink}{#1}}}
\newtcolorbox{experience}[1][]{popbase,colframe=popblue,colback=popblueL,before upper={\popboxhead{Expérience}{popblue}{#1}}}
% « Le savais-tu ? » : carte violette pleine (contexte, culture, Cameroun)
\newtcolorbox{savaistu}[1][]{popbase,colback=poppurple,colframe=popink,
  fontupper=\popsemi\fontsize{10.4}{14.2}\selectfont\color{white},
  before upper={\poppill{Le savais-tu\,?}{white}{poppurple}\ifx\relax#1\relax\else\hspace{6pt}{\popxbold\color{white}#1}\fi\par\vspace{7pt}}}
% Exercice résolu : énoncé en haut, \tcblower puis la solution sur fond crème
\newcounter{popexo}\newcounter{popexr}
\newtcolorbox{exoresolu}[1][]{popbase,colframe=poporange,colbacklower=popcream,
  segmentation style={popink,line width=1.2pt,dashed},
  before upper={\stepcounter{popexr}\popboxhead{Exercice résolu \thepopexr}{poporange}{#1}},
  before lower={\poppill{Solution}{popink}{popcream}\par\vspace{6pt}}}
% Exercice (non résolu en place) — la correction est dans la partie Corrigés
\newtcolorbox{exercice}[1][]{popbase,colframe=popink,
  before upper={\stepcounter{popexo}\popboxhead{Exercice \thepopexo}{popink}{#1}}}
% Corrigé
\newtcolorbox{corrige}[1][]{popbase,colframe=popgreen,colback=popgreenL,
  before upper={\popboxhead{Corrigé}{popgreen}{#1}}}
% Objectifs / prérequis / bilan
\newtcolorbox{objectifs}{popbase,colframe=popink,colback=poporangeL,
  before upper={\popboxhead{Objectifs de la leçon}{popink}{}}}
\newtcolorbox{prerequis}{popbase,colframe=popink,colback=popblueL,
  before upper={\popboxhead{Prérequis}{popblue}{}}}
\newtcolorbox{bilan}{popbase,colframe=popink,colback=white,boxrule=2.8pt,
  before upper={\popboxhead{Fiche bilan}{poporange}{L'essentiel à connaître pour l'examen}}}
% Figure encadrée avec légende
\newtcolorbox{popfigure}[1][]{enhanced,breakable=false,colback=white,colframe=popline,boxrule=1.4pt,arc=8pt,
  left=10pt,right=10pt,top=10pt,bottom=8pt,before skip=10pt,after skip=12pt,halign=center,
  lower separated=false,halign lower=center,
  fontlower=\popsemi\fontsize{9}{11.5}\selectfont\color{popmuted},
  #1}
\newcommand{\legende}[1]{\par\vspace{6pt}{\popsemi\fontsize{9}{11.5}\selectfont\color{popmuted}#1}}

% ============================================================
% Signature OnBuch+
% ============================================================
\newcommand{\poplogoinline}[1]{{\poplogo\fontsize{#1}{#1}\selectfont\color{popink}OnBuch\color{poporange}+}}
\newcommand{\signatureonbuch}{%
  \begin{tcolorbox}[enhanced,colback=popink,colframe=popink,boxrule=2.2pt,arc=10pt,
      drop shadow=poporange,left=14pt,right=14pt,top=10pt,bottom=10pt,before skip=0pt,after skip=0pt]
    \tikz[baseline=-0.7ex]\node[circle,fill=popgreen,draw=popcream,line width=1.3pt,minimum size=22pt,inner sep=0]{\popcheck{white}};%
    \hspace{9pt}\begin{minipage}[c]{0.62\linewidth}
      {\popxbold\fontsize{12}{13}\selectfont\color{popcream}Signé\ \textbullet\ {\poplogo OnBuch\color{poporange}+}}\par\vspace{2pt}
      {\raggedright\popsemi\fontsize{8}{10.5}\selectfont\color{popline}\MakeUppercase{Équipe pédagogique OnBuch+}\\\MakeUppercase{Conforme au programme officiel MINESEC}\par}
    \end{minipage}\hfill
    {\poplogo\fontsize{22}{22}\selectfont\color{popcream}OnBuch\color{poporange}+}
  \end{tcolorbox}%
}

% ============================================================
% Page de garde
% #1 titre · #2 matière · #3 niveau · #4 module · #5 n° leçon · #6 durée ·
% #7 description / objectifs (texte ou itemize)
% ============================================================
\newcommand{\pagegarde}[7]{%
  \thispagestyle{empty}
  \newgeometry{a4paper,margin=1.7cm,top=1.6cm,bottom=1.4cm}%
  \begin{tikzpicture}[remember picture,overlay]
    \fill[poporange!18] ([xshift=-3.2cm,yshift=-3.6cm]current page.north east) circle (4.6cm);
    \fill[poppurple!14] ([xshift=2.4cm,yshift=7.6cm]current page.south west) circle (3.8cm);
  \end{tikzpicture}%
  {\poplogo\fontsize{30}{30}\selectfont\color{popink}OnBuch\color{poporange}+}\hfill
  \raisebox{0.6ex}{\poppill{Cours signé \textbullet\ Premium}{popink}{popcream}}\par
  \vspace{1pt}\popsquiggle{poppurple}\par
  \vspace{8pt}
  \begin{tcolorbox}[enhanced,colback=poporange,colframe=popink,boxrule=2.8pt,arc=13pt,
      drop shadow=popink,left=18pt,right=18pt,top=12pt,bottom=13pt,after skip=10pt]
    \poppill{#2 \textbullet\ #3}{popink}{popcream}\hspace{5pt}\poppill{#4}{white}{popink}\par\vspace{8pt}
    {\popdisplay\fontsize{14}{14}\selectfont\color{popink}LEÇON #5}\par\vspace{4pt}
    {\raggedright\hyphenpenalty=10000\exhyphenpenalty=10000\popdisplay\fontsize{25}{27}\selectfont\color{popcream}\MakeUppercase{#1}\par}
    \vspace{8pt}
    {\popxbold\fontsize{9.5}{9.5}\selectfont\color{popink}\MakeUppercase{Durée conseillée : #6}}
  \end{tcolorbox}
  \begin{tcolorbox}[enhanced,colback=white,colframe=popink,boxrule=2.2pt,arc=10pt,
      drop shadow=popink,left=16pt,right=16pt,top=10pt,bottom=10pt,after skip=8pt]
    \poppill{Dans cette leçon}{poppurple}{white}\par\vspace{5pt}
    \popsemi\fontsize{9.3}{12.6}\selectfont\color{popink}#7
  \end{tcolorbox}
  \noindent\begin{minipage}[t]{0.487\linewidth}
    \begin{tcolorbox}[enhanced,colback=popgreen,colframe=popink,boxrule=2.2pt,arc=10pt,
        drop shadow=popink,left=11pt,right=11pt,top=10pt,bottom=10pt,height=3.1cm,valign=center]
      \tcbox[colback=white,colframe=popink,boxrule=1.3pt,arc=5pt,boxsep=0pt,nobeforeafter,
        left=3pt,right=3pt,top=3pt,bottom=3pt,valign=center]{\qrcode[height=1.9cm]{https://play.google.com/store/apps/details?id=com.luvvix.onbuch&hl=fr}}%
      \hspace{8pt}\begin{minipage}[c]{0.5\linewidth}
        {\popdisplay\fontsize{11}{12.5}\selectfont\color{white}\MakeUppercase{Télécharge OnBuch}}\par\vspace{4pt}
        {\raggedright\popsemi\fontsize{8.4}{11}\selectfont\color{white}Gratuit \textbullet\ Google Play\\Tuteur IA, annales, concours, résultats\par}
      \end{minipage}
    \end{tcolorbox}
  \end{minipage}\hfill
  \begin{minipage}[t]{0.487\linewidth}
    \begin{tcolorbox}[enhanced,colback=poppurple,colframe=popink,boxrule=2.2pt,arc=10pt,
        drop shadow=popink,left=11pt,right=11pt,top=10pt,bottom=10pt,height=3.1cm,valign=center]
      \tikz[baseline=-0.7ex]\node[circle,fill=white,draw=popink,line width=1.3pt,minimum size=1.6cm,inner sep=0]{\popdisplay\fontsize{24}{24}\selectfont\color{poppurple}+};%
      \hspace{8pt}\begin{minipage}[c]{0.6\linewidth}
        {\popdisplay\fontsize{11}{12.5}\selectfont\color{white}\MakeUppercase{Passe à OnBuch+}}\par\vspace{4pt}
        {\raggedright\popsemi\fontsize{8.4}{11}\selectfont\color{white}Tous les cours signés, Tuteur IA illimité, fiches PDF — onglet OnBuch+ de l'app\par}
      \end{minipage}
    \end{tcolorbox}
  \end{minipage}
  \vspace{8pt}
  \popcontenu
  \vfill
  \signatureonbuch
  \restoregeometry
  \newpage
}

% Rangée « Ce cours contient » (page de garde)
\newcommand{\poptile}[3]{% #1 couleur #2 gros texte #3 légende
  \begin{tcolorbox}[enhanced,colback=#1,colframe=popink,boxrule=2pt,arc=9pt,shadow={2.5pt}{-2.5pt}{0pt}{popink},
      left=6pt,right=6pt,top=9pt,bottom=9pt,halign=center,valign=center,height=1.75cm,width=\linewidth,nobeforeafter]
    {\popdisplay\fontsize{12.5}{13}\selectfont\color{white}\MakeUppercase{#2}}\par\vspace{3pt}
    {\centering\popsemi\fontsize{7.8}{9.5}\selectfont\color{white}#3\par}
  \end{tcolorbox}}
\newcommand{\popcontenu}{%
  \par\noindent{\popxboldls\fontsize{7.5}{7.5}\selectfont\color{popmuted}CE COURS SIGNÉ CONTIENT}\par\vspace{7pt}
  \noindent\begin{minipage}[t]{0.235\linewidth}\poptile{popblue}{Cours}{complet, illustré, pas à pas}\end{minipage}\hfill
  \begin{minipage}[t]{0.235\linewidth}\poptile{poporange}{Méthodes}{et exercices résolus}\end{minipage}\hfill
  \begin{minipage}[t]{0.235\linewidth}\poptile{poppink}{Exercices}{type examen + corrigés}\end{minipage}\hfill
  \begin{minipage}[t]{0.235\linewidth}\poptile{popgreen}{Fiche bilan}{l'essentiel à retenir}\end{minipage}\par}

% Sommaire
\newtcolorbox{sommaire}{enhanced,colback=white,colframe=popink,boxrule=2.2pt,arc=10pt,
  drop shadow=popink,left=18pt,right=18pt,top=13pt,bottom=13pt,before skip=6pt,after skip=20pt,
  before upper={\poppill{Sommaire}{poppurple}{white}\par\vspace{9pt}\popsemi\fontsize{10.6}{14.5}\selectfont\color{popink}}}

% Signature de fin de cours — poussée tout en bas de la dernière page
\newcommand{\findecours}{%
  \par\vfill\nopagebreak
  \begin{center}\popsquiggle{poporange}\end{center}\vspace{4pt}
  \signatureonbuch
}
\AtBeginDocument{\def\llbracket{\mathopen{[\mkern-3.5mu[}}\def\rrbracket{\mathclose{]\mkern-3.5mu]}}} % intervalles d'entiers (glyphe ⟦ absent de la police)
% Glyphes absents de Fira Math : redéfinis à partir de glyphes disponibles
\AtBeginDocument{%
  \def\setminus{\mathbin{\backslash}}\let\smallsetminus\setminus
  \def\vdots{\mathord{\vbox{\baselineskip3.2pt\lineskiplimit0pt\kern4pt\hbox{.}\hbox{.}\hbox{.}}}}%
  \def\ddots{\mathinner{\mkern1mu\raise6.5pt\hbox{.}\mkern2mu\raise3.6pt\hbox{.}\mkern2mu\raise0.7pt\hbox{.}\mkern1mu}}%
  \def\triangle{\mathord{\scalebox{0.9}{$\symup{\Delta}$}}}%
  \def\nabla{\mathord{\raisebox{\depth}{\scalebox{1}[-1]{$\symup{\Delta}$}}}}%
  \def\checkmark{\popcheck{popdark}}%
  \def\star{\mathbin{\ast}}\def\bigstar{\mathord{\ast}}%
  \def\diamond{\mathbin{\scalebox{0.6}{$\Diamond$}}}%
  \def\bigcup{\mathop{\mathchoice{\raisebox{-0.3em}{\scalebox{1.7}{$\cup$}}}{\scalebox{1.25}{$\cup$}}{\cup}{\cup}}\displaylimits}%
  \def\bigcap{\mathop{\mathchoice{\raisebox{-0.3em}{\scalebox{1.7}{$\cap$}}}{\scalebox{1.25}{$\cap$}}{\cap}{\cap}}\displaylimits}%
  \def\lesseqqgtr{\mathrel{\lessgtr}}\def\bigoplus{\mathop{\oplus}\displaylimits}%
}
\providecommand{\ssi}{si et seulement si} % abréviation fréquente
\AtBeginDocument{\providecommand{\wideparen}{\overparen}} % arc de cercle

%% FIGFIX-BEGIN v3 — correctifs globaux des figures (docs/onbuch-generation/tools/figfix)
\usepackage{adjustbox}\usepackage{etoolbox}
% toute figure TikZ/pgfplots plus large que la ligne est réduite à la largeur disponible
\BeforeBeginEnvironment{tikzpicture}{\begin{adjustbox}{max width=\linewidth}}
\AfterEndEnvironment{tikzpicture}{\end{adjustbox}}
% légendes pgfplots lisibles par-dessus les courbes
\pgfplotsset{every axis/.append style={legend style={fill=white,fill opacity=0.92,text opacity=1,draw=black!15,inner sep=3pt}}}
% étiquettes d'arêtes (syntaxe edge["..."]) avec fond blanc
\tikzset{every edge quotes/.append style={fill=white,inner sep=1.2pt}}
% bulles de cartes mentales : texte à la ligne dans la bulle, bulle un peu plus grande
\tikzset{every concept/.append style={text width=2.0cm,align=center,minimum size=2.3cm,inner sep=2pt}}
%% FIGFIX-END
\begin{document}

\pagegarde{La familia : evolución y conflictos intergeneracionales}{Espagnol}{1ère A}{Módulo 1 : Vie familiale et intégration sociale}{EA01}{2 h}{Cette leçon mixte étudie l'évolution des modèles familiaux (famille nucléaire, monoparentale, élargie, recomposée) et les causes des conflits intergénérationnels. Elle fournit le lexique de la famille et les outils grammaticaux pour exprimer la cause, la conséquence et le conseil en espagnol.
\vspace{4pt}
\begin{multicols}{2}\small
\begin{itemize}
\item À la fin de cette leçon, je sais lire et comprendre un texte espagnol sur l'évolution de la famille.
\item À la fin de cette leçon, je sais définir et distinguer les différents types de familles : nuclear, monoparental, extensa, reconstituida.
\item À la fin de cette leçon, je sais utiliser le lexique des relations familiales et des générations (abuelos, padres, hijos, brecha generacional, autoridad, respeto, conflicto).
\item À la fin de cette leçon, je sais identifier dans un texte les causes des conflits intergénérationnels (NTIC, éducation, valeurs, urbanisation).
\item À la fin de cette leçon, je sais exprimer la cause avec porque, pues, ya que, como et gracias a / debido a / a causa de.
\item À la fin de cette leçon, je sais exprimer la conséquence avec por eso, así que, entonces, de modo que.
\end{itemize}\end{multicols}}

\begin{sommaire}
\begin{multicols}{2}
\begin{enumerate}
\item Texte support : « La familia cambia »
\item Lexique : la famille et les générations
\item Grammaire 1 : exprimer la cause
\item Grammaire 2 : exprimer la conséquence
\item Grammaire 3 : exprimer le conseil
\item Méthode du commentaire et commentaire modèle
\item Activité d'intégration
\item Méthodes et exercices résolus
\item Exercices
\item Corrigés des exercices
\item Fiche bilan
\end{enumerate}
\end{multicols}
\end{sommaire}

\begin{objectifs}
\begin{itemize}
\item À la fin de cette leçon, je sais lire et comprendre un texte espagnol sur l'évolution de la famille.
\item À la fin de cette leçon, je sais définir et distinguer les différents types de familles : nuclear, monoparental, extensa, reconstituida.
\item À la fin de cette leçon, je sais utiliser le lexique des relations familiales et des générations (abuelos, padres, hijos, brecha generacional, autoridad, respeto, conflicto).
\item À la fin de cette leçon, je sais identifier dans un texte les causes des conflits intergénérationnels (NTIC, éducation, valeurs, urbanisation).
\item À la fin de cette leçon, je sais exprimer la cause avec porque, pues, ya que, como et gracias a / debido a / a causa de.
\item À la fin de cette leçon, je sais exprimer la conséquence avec por eso, así que, entonces, de modo que.
\item À la fin de cette leçon, je sais donner un conseil avec deberías, hay que, es importante que + subjonctif (notion d'initiation).
\item À la fin de cette leçon, je sais rédiger quelques phrases pour proposer des solutions au dialogue familial.
\item À la fin de cette leçon, je sais traduire un court texte du français vers l'espagnol sur le thème de la famille.
\end{itemize}
\end{objectifs}

\begin{prerequis}
\begin{itemize}
\item Le présent de l'indicatif des verbes réguliers et irréguliers courants (Seconde).
\item Le lexique de base de la famille vu en Seconde (padre, madre, hermano, tío, primo...).
\item Les adjectifs possessifs et les articles définis (mi familia, los abuelos).
\item La construction de phrases simples avec connecteurs de base (y, pero, también).
\item La lecture méthodique d'un texte : repérage du thème, du locuteur, des idées principales (Seconde).
\end{itemize}
\end{prerequis}

\newpage

\coursec{Texte support : « La familia cambia »}

\courssub{Situation d'accroche et problématique}

Imaginez la scène. À Yaoundé, quartier Mvog-Ada, la famille Etoundi s'installe autour de la table du soir. Grand-père Samuel commence un récit sur son village d'Awae, mais Sandrine, seize ans, les yeux rivés sur son smartphone, répond à peine. « De mon temps, on écoutait les anciens ! » s'emporte le grand-père. « De ton temps, il n'y avait pas Internet ! » réplique l'adolescente. Le dîner se transforme en champ de bataille.

À Madrid, même scène ou presque : la famille García dîne ensemble, mais le fils Pablo consulte ses messages sous la table et la grand-mère Carmen soupire en se souvenant des longues conversations d'autrefois.

Au Cameroun, la famille élargie (\esp{familia extensa}) cède peu à peu la place à la famille nucléaire dans les grandes villes comme Douala ou Yaoundé, et les jeunes reprochent parfois aux aînés leur autorité « à l'ancienne ». En Espagne et en Amérique latine, on parle de \esp{brecha generacional} pour décrire ce fossé entre les générations.

\textbf{Problématique :} Comment la famille a-t-elle évolué en Espagne, en Amérique latine et en Afrique urbaine ? Pourquoi les conflits entre générations éclatent-ils, et comment rétablir le dialogue familial ? C'est tout l'objet de cette leçon. Pour y répondre, nous allons d'abord lire un texte, comme on vous le demandera au baccalauréat.

\courssub{Méthode : lire un texte au baccalauréat}

\begin{methode}[Cinq étapes pour comprendre un texte]
\begin{enumerate}
\item \textbf{Lire le titre et la source.} Le titre annonce le thème : ici, \esp{La familia cambia} (« La famille change ») annonce une évolution, pas une description fixe.
\item \textbf{Faire une première lecture globale, sans dictionnaire.} Cherchez le sens général : de quoi parle le texte ? Qui ? Où ?
\item \textbf{À la deuxième lecture, souligner les mots-clés.} Repérez les mots transparents (\esp{familia, conflicto, diálogo}) et les mots nouveaux que le contexte permet de deviner.
\item \textbf{Répondre aux questions en citant le texte.} Une réponse au Bac doit toujours être justifiée par une citation entre guillemets.
\item \textbf{Rédiger en phrases complètes}, en espagnol, avec sujet + verbe + complément : \esp{El texto dice que...} (« Le texte dit que... »).
\end{enumerate}
\end{methode}

\courssub{Lecture du texte}

\begin{textebox}[La familia cambia]
Antes, en muchos pueblos de España, de América Latina y de África, la familia extensa reunía bajo el mismo techo a abuelos, padres, hijos y tíos. Todos comían juntos, los ancianos contaban historias y los niños aprendían escuchando. La convivencia era la escuela de la vida.

Hoy, en las grandes ciudades como Madrid, Ciudad de México o Yaundé, domina la familia nuclear: solamente el padre, la madre y los hijos viven en la misma casa. Además, aparecen nuevos modelos: la familia monoparental, donde un solo progenitor cría a los niños, y la familia reconstituida, formada después de un divorcio, cuando dos adultos con hijos deciden vivir juntos.

Esta evolución tiene ventajas, pero también provoca conflictos. Los jóvenes viven conectados a sus teléfonos y a las redes sociales, mientras que los abuelos añoran las conversaciones de antes. Es la brecha digital. Los padres quieren mantener la autoridad y los valores tradicionales, pero los hijos, influenciados por la urbanización y por otras culturas, piensan de manera diferente. La educación moderna, más libre, choca a veces con la disciplina de antes.

Sin embargo, el conflicto no es inevitable. La solución existe: el diálogo. Si los padres escuchan a sus hijos y los hijos respetan a sus padres, la familia puede encontrar un equilibrio. La escucha activa, el respeto mutuo y el tiempo compartido son las llaves de la convivencia feliz. La familia cambia, sí, pero el amor que la une no tiene por qué desaparecer.
\end{textebox}

\textbf{Traduction du texte.} « Autrefois, dans beaucoup de villages d'Espagne, d'Amérique latine et d'Afrique, la famille élargie réunissait sous le même toit grands-parents, parents, enfants et oncles. Tous mangeaient ensemble, les anciens racontaient des histoires et les enfants apprenaient en écoutant. La cohabitation était l'école de la vie.

Aujourd'hui, dans les grandes villes comme Madrid, Mexico ou Yaoundé, domine la famille nucléaire : seuls le père, la mère et les enfants vivent dans la même maison. De plus, de nouveaux modèles apparaissent : la famille monoparentale, où un seul parent élève les enfants, et la famille recomposée, formée après un divorce, quand deux adultes avec enfants décident de vivre ensemble.

Cette évolution a des avantages, mais elle provoque aussi des conflits. Les jeunes vivent connectés à leurs téléphones et aux réseaux sociaux, tandis que les grands-parents regrettent les conversations d'autrefois. C'est la fracture numérique. Les parents veulent maintenir l'autorité et les valeurs traditionnelles, mais les enfants, influencés par l'urbanisation et par d'autres cultures, pensent différemment. L'éducation moderne, plus libre, entre parfois en conflit avec la discipline d'autrefois.

Cependant, le conflit n'est pas inévitable. La solution existe : le dialogue. Si les parents écoutent leurs enfants et si les enfants respectent leurs parents, la famille peut trouver un équilibre. L'écoute active, le respect mutuel et le temps partagé sont les clés d'une cohabitation heureuse. La famille change, oui, mais l'amour qui l'unit ne doit pas disparaître. »

\courssub{Glossaire du texte}

\begin{definition}[Vocabulaire essentiel du texte]
\begin{itemize}
\item \esp{la brecha} : le fossé, l'écart. \esp{la brecha generacional} : le fossé entre les générations ; \esp{la brecha digital} : la fracture numérique (écart entre ceux qui maîtrisent les technologies et ceux qui ne les maîtrisent pas).
\item \esp{la convivencia} : la cohabitation, le fait de vivre ensemble.
\item \esp{criar} : élever (des enfants). Attention : \esp{criar} signifie « élever », pas « créer ».
\item \esp{añorar} : regretter avec nostalgie, avoir la nostalgie de. \esp{Los abuelos añoran las conversaciones de antes.} « Les grands-parents regrettent les conversations d'autrefois. »
\item \esp{bajo el mismo techo} : sous le même toit.
\item \esp{las redes sociales} : les réseaux sociaux.
\item \esp{la escucha activa} : l'écoute active.
\item \esp{el respeto mutuo} : le respect mutuel.
\item \esp{chocar con} : entrer en conflit avec, heurter.
\item \esp{no tiene por qué desaparecer} : ne doit pas forcément disparaître.
\end{itemize}
\end{definition}

\begin{definition}[La brecha generacional]
La \esp{brecha generacional} est le fossé culturel, technologique et de valeurs qui sépare les générations : \esp{padres} et \esp{hijos}, \esp{abuelos} et \esp{nietos}. Elle se manifeste dans les goûts (musique, habits), les outils (smartphone contre radio), et les idées (autorité, liberté, mariage). Ce n'est pas un phénomène nouveau, mais les nouvelles technologies l'ont considérablement élargi.
\end{definition}

\courssub{Les quatre modèles familiaux du texte}

Le texte présente quatre modèles de famille. Observons-les dans un tableau avant de les mémoriser.

\begin{center}
\begin{tabularx}{\linewidth}{|X|X|X|}
\hline
\rowcolor{popblueL} Modelo & Definición & Ejemplo del texto \\
\hline
\esp{familia extensa} & famille élargie : grands-parents, parents, enfants, oncles sous le même toit & \esp{reunía bajo el mismo techo a abuelos, padres, hijos y tíos} \\
\hline
\esp{familia nuclear} & famille nucléaire : père, mère et enfants seulement & \esp{solamente el padre, la madre y los hijos} \\
\hline
\esp{familia monoparental} & famille monoparentale : un seul parent élève les enfants & \esp{un solo progenitor cría a los niños} \\
\hline
\esp{familia reconstituida} & famille recomposée : deux adultes avec enfants d'unions précédentes & \esp{formada después de un divorcio} \\
\hline
\end{tabularx}
\end{center}

\begin{attention}[Faux amis et pièges de traduction]
\begin{itemize}
\item \esp{la familia extensa} se dit « famille élargie » en français, et non « famille extensive ». On trouve aussi \esp{familia extendida} en Amérique latine.
\item \esp{el progenitor} est un mot soutenu pour « parent » (père ou mère). Dans la conversation courante, on dit plutôt \esp{el padre} ou \esp{la madre}.
\item \esp{criar} = élever (des enfants) ; « créer » se dit \esp{crear}. Ne confondez pas : \esp{Ella cría a sus hijos sola.} « Elle élève seule ses enfants. »
\item \esp{asistir} ne signifie pas « assister » au sens d'aider : c'est « être présent, assister à ». « Aider » se dit \esp{ayudar}.
\end{itemize}
\end{attention}

\courssub{Carte mentale : la famille, ses modèles, ses conflits et ses solutions}

\begin{popfigure}
\begin{tikzpicture}[
  mindmap,
  grow cyclic,
  every node/.style={concept, align=center, font=\small},
  root concept/.append style={concept color=popblue, font=\bfseries},
  level 1/.append style={level distance=3.4cm, sibling angle=90},
  level 2/.append style={level distance=2.6cm, sibling angle=60},
]
\node[root concept]{La familia}
  child[concept color=popgreen] { node {nuclear} }
  child[concept color=poporange] { node {monoparental} }
  child[concept color=poppurple] { node {extensa}
    child[concept color=poppink] { node[align=center]{causas de\\conflicto:\\NTIC, valores,\\educación,\\urbanización} }
    child[concept color=popgold] { node[align=center]{soluciones:\\diálogo, respeto\\mutuo, escucha\\activa} }
  }
  child[concept color=popblue!60] { node {reconstituida} };
\end{tikzpicture}
\legende{Carte mentale du texte : les quatre modèles familiaux, les causes des conflits et les solutions.}
\end{popfigure}

\courssub{Repérage dans le texte : causes des conflits et solutions}

\textbf{Les causes de conflits mentionnées dans le texte :}
\begin{itemize}
\item \textbf{Les NTIC} (nouvelles technologies) : \esp{« Los jóvenes viven conectados a sus teléfonos y a las redes sociales »} — c'est la \esp{brecha digital}.
\item \textbf{Les valeurs} : \esp{« Los padres quieren mantener la autoridad y los valores tradicionales »}.
\item \textbf{L'éducation} : \esp{« La educación moderna, más libre, choca a veces con la disciplina de antes »}.
\item \textbf{L'urbanisation} : \esp{« influenciados por la urbanización y por otras culturas »}.
\end{itemize}

\textbf{Les solutions proposées par le texte :}
\begin{itemize}
\item \esp{el diálogo} : le dialogue, « la solution existe : le dialogue » ;
\item \esp{el respeto mutuo} : le respect mutuel entre parents et enfants ;
\item \esp{la escucha activa} : l'écoute active, c'est-à-dire écouter vraiment l'autre, pas seulement l'entendre ;
\item \esp{el tiempo compartido} : le temps partagé en famille.
\end{itemize}

\begin{aretenir}[Structure argumentative du texte]
Le texte suit un plan en trois temps, typique des textes de réflexion : \textbf{1)} un constat (la famille a changé : quatre modèles) ; \textbf{2)} un problème (la \esp{brecha generacional} et ses causes) ; \textbf{3)} une solution (le dialogue, le respect, l'écoute). Retenez ce plan : il vous servira pour vos propres rédactions.
\end{aretenir}

\begin{savaistu}[La famille en Espagne aujourd'hui]
En Espagne, plus de deux millions de familles sont monoparentales, et dans la grande majorité des cas, c'est la mère qui élève seule les enfants. La \esp{familia reconstituida} (famille recomposée) est en forte augmentation depuis les années 2000, à la faveur des divorces et des nouvelles unions. En Guinée équatoriale, pays hispanophone voisin du Cameroun, la famille élargie reste très présente, comme en Afrique centrale : les oncles, tantes et grands-parents participent activement à l'éducation des enfants.
\end{savaistu}

\courssub{Questions de compréhension}

\begin{exercice}[Compréhension globale et détaillée]
\textbf{A. Vrai ou faux ? Justifiez chaque réponse par une citation du texte.}
\begin{enumerate}
\item \esp{Antes, la familia extensa vivía separada en varias casas.}
\item \esp{Hoy, en las grandes ciudades, domina la familia nuclear.}
\item \esp{Los abuelos están contentos con los teléfonos de los jóvenes.}
\item \esp{El texto propone el diálogo como solución.}
\end{enumerate}
\textbf{B. Questions ouvertes.}
\begin{enumerate}
\item \esp{¿Qué cuatro modelos de familia presenta el texto?}
\item \esp{¿Qué es la brecha digital, según el texto?}
\item \esp{¿Cuáles son las tres soluciones propuestas para la convivencia feliz?}
\item \esp{¿Y en tu familia, en Yaundé o en tu pueblo, hay conflictos entre generaciones? ¿Por qué?}
\end{enumerate}
\end{exercice}

\begin{exoresolu}[Modèle de réponse justifiée (question A.1)]
Énoncé : \esp{Antes, la familia extensa vivía separada en varias casas.} Vrai ou faux ? Justifiez par une citation.
\tcblower
\textbf{Solution.} Étape 1 : je relis le premier paragraphe, qui parle du passé (\esp{antes}). Étape 2 : je repère la phrase clé : \esp{« la familia extensa reunía bajo el mismo techo a abuelos, padres, hijos y tíos »}. Étape 3 : « sous le même toit » est le contraire de « séparée en plusieurs maisons » : l'affirmation est donc \textbf{fausse}.

Rédaction en phrase complète : \esp{Falso. El texto dice que « la familia extensa reunía bajo el mismo techo a abuelos, padres, hijos y tíos ». Por lo tanto, la familia vivía junta, no separada.}

Traduction : « Faux. Le texte dit que "la famille élargie réunissait sous le même toit grands-parents, parents, enfants et oncles". Par conséquent, la famille vivait ensemble, pas séparée. » Remarquez la formule \esp{por lo tanto} (« par conséquent »), très utile pour conclure une justification.
\end{exoresolu}

\begin{corrige}[Questions A.2 à A.4 et B]
\textbf{A.2.} Vrai : \esp{« Hoy, en las grandes ciudades... domina la familia nuclear »}. \textbf{A.3.} Faux : \esp{« los abuelos añoran las conversaciones de antes »} — les grands-parents regrettent le passé, ils ne sont donc pas contents. \textbf{A.4.} Vrai : \esp{« La solución existe: el diálogo »}.

\textbf{B.1.} \esp{El texto presenta la familia extensa, la familia nuclear, la familia monoparental y la familia reconstituida.} \textbf{B.2.} \esp{La brecha digital es el fossé tecnológico entre los jóvenes, conectados a sus teléfonos, y los abuelos, que añoran las conversaciones de antes.} \textbf{B.3.} \esp{Las soluciones son la escucha activa, el respeto mutuo y el tiempo compartido.} \textbf{B.4.} Réponse personnelle ; modèle : \esp{En mi familia hay conflictos porque mis padres no comprenden mi pasión por el fútbol y por Internet.} « Dans ma famille, il y a des conflits parce que mes parents ne comprennent pas ma passion pour le football et pour Internet. »
\end{corrige}

\begin{attention}[Erreurs fréquentes des francophones dans les réponses]
\begin{itemize}
\item Ne traduisez pas mot à mot « selon le texte » par \esp{según de el texto} : on dit simplement \esp{según el texto}.
\item « Citer le texte » se dit \esp{citar el texto} ; la citation s'introduit avec \esp{El texto dice que...} suivi de la phrase entre guillemets.
\item Attention au genre : \esp{la familia} est féminin, donc \esp{la familia extensa}, \esp{una familia monoparental} (l'adjectif \esp{monoparental} est invariable en genre).
\item En espagnol, les guillemets s'écrivent « ... » ou « ... », mais jamais avec une espace avant les deux-points : on écrit \esp{existe: el diálogo} et non « existe : el diálogo ».
\end{itemize}
\end{attention}

\begin{experience}[À vous de parler : mini-débat de classe]
Par groupes de quatre, préparez en cinq minutes un mini-débat. Deux élèves jouent les \esp{abuelos} (ils défendent les traditions et les conversations en famille), deux élèves jouent les \esp{jóvenes} (ils défendent le smartphone et la liberté). Chacun doit utiliser au moins deux mots du glossaire : \esp{brecha generacional, convivencia, respeto mutuo, escucha activa}. Puis, toute la classe vote pour la meilleure solution de dialogue familial.
\end{experience}

\begin{aretenir}[L'essentiel de la section]
\begin{itemize}
\item La famille a évolué : \esp{familia extensa} (traditionnelle, rurale) → \esp{familia nuclear} (urbaine), avec deux nouveaux modèles : \esp{monoparental} et \esp{reconstituida}.
\item Les conflits entre générations (\esp{brecha generacional}) ont quatre causes principales : les NTIC (\esp{brecha digital}), l'éducation, les valeurs, l'urbanisation.
\item Les solutions : \esp{el diálogo, el respeto mutuo, la escucha activa, el tiempo compartido}.
\item Au Bac, toute réponse de compréhension doit être rédigée en phrase complète et justifiée par une citation du texte.
\end{itemize}
\end{aretenir}

\coursec{Lexique : la famille et les générations}

Dans la section précédente, nous avons lu le texte \esp{« La familia cambia »}, qui décrivait les grandes transformations de la famille aujourd'hui : familles nucléaires, monoparentales, élargies ou recomposées, et tensions entre parents et adolescents. Pour analyser ce texte et, plus tard, rédiger vos propres productions, il faut d'abord maîtriser un vocabulaire précis. Cette section vous donne toutes les \cle{palabras clave} (mots-clés) du thème de la famille et des relations entre générations, organisées en quatre champs lexicaux, avec leurs pièges de traduction.

\courssub{Champ 1 : les membres de la famille}

\begin{definition}[Les membres de la famille]
\esp{Los miembros de la familia} désignent toutes les personnes unies par des liens de parenté. En espagnol, chaque mot possède une forme masculine et une forme féminine, et le pluriel masculin peut désigner un groupe mixte : \esp{los hijos} = les enfants (garçons et filles).
\end{definition}

\begin{center}
\begin{tabularx}{\linewidth}{|X|X|X|}
\hline
\rowcolor{popblueL} Español & Français & Remarque \\
\hline
el abuelo / la abuela & le grand-père / la grand-mère & \esp{los abuelos} = les grands-parents \\
\hline
el padre / la madre & le père / la mère & \esp{los padres} = les parents \\
\hline
el hijo / la hija & le fils / la fille & \esp{los hijos} = les enfants \\
\hline
el nieto / la nieta & le petit-fils / la petite-fille & \esp{los nietos} = les petits-enfants \\
\hline
el tío / la tía & l'oncle / la tante & accent sur le í : \esp{tío} \\
\hline
el primo / la prima & le cousin / la cousine & \esp{primo hermano} = cousin germain \\
\hline
el hermano / la hermana & le frère / la sœur & \esp{los hermanos} = les frères et sœurs \\
\hline
el padrastro / la madrastra & le beau-père / la belle-mère (remariage) & famille recomposée \\
\hline
el hermanastro / la hermanastra & le demi-frère / la demi-sœur & sans lien de sang \\
\hline
el suegro / la suegra & le beau-père / la belle-mère (par alliance) & parents du conjoint \\
\hline
el pariente / la pariente & le parent, la proche & famille élargie (épicène) \\
\hline
\end{tabularx}
\end{center}

\begin{attention}[Faux ami : « los padres » et « los parientes »]
\esp{Los padres} signifie « les parents » (le père \textbf{et} la mère), et non « les pères ». Pour parler de plusieurs pères, on dit \esp{los padres de familia} ou le contexte doit le préciser. De même, \esp{los parientes} est un faux ami : il ne désigne pas le père et la mère, mais les \textbf{parents au sens large} (oncles, cousins, beaux-frères...). Ne traduisez donc jamais « mes parents viennent me voir » par \esp{mis parientes vienen a verme} si vous parlez de votre père et de votre mère : dites \esp{mis padres vienen a verme}.
\end{attention}

\begin{exemplebox}[Exemples traduits]
\esp{Mis abuelos viven con nosotros.} « Mes grands-parents vivent avec nous. »

\esp{Mi tío y mis primos llegan de Douala el sábado.} « Mon oncle et mes cousins arrivent de Douala samedi. »

\esp{Su madrastra tiene dos hijos de un primer matrimonio.} « Sa belle-mère a deux enfants d'un premier mariage. »

\esp{Mis suegros celebran sus bodas de oro.} « Mes beaux-parents fêtent leurs noces d'or. »
\end{exemplebox}

\begin{popfigure}
\begin{center}
\begin{tikzpicture}[
  every node/.style={draw, rounded corners=2pt, align=center, font=\small, inner sep=4pt},
  abuelo/.style={fill=popgoldL},
  padre/.style={fill=popblueL},
  hijo/.style={fill=popgreenL},
  lado/.style={fill=poppurpleL},
  lin/.style={draw=popmuted, thick}]
  \node[abuelo, align=center] (abu) at (0,4){los abuelos\\(el abuelo / la abuela)};
  \node[padre, align=center] (pad) at (0,2.2){los padres\\(el padre / la madre)};
  \node[lado, align=center] (tio) at (4.6,2.2){el tío / la tía\\(hermanos del padre)};
  \node[hijo, align=center] (hij) at (0,0.4){los hijos\\(el hijo / la hija)};
  \node[lado, align=center] (pri) at (4.6,0.4){los primos\\(hijos de los tíos)};
  \node[hijo, align=center] (nie) at (0,-1.4){los nietos\\(el nieto / la nieta)};
  \draw[lin] (abu) -- (pad);
  \draw[lin] (pad) -- (hij);
  \draw[lin] (hij) -- (nie);
  \draw[lin, dashed] (pad) -- (tio);
  \draw[lin, dashed] (tio) -- (pri);
  \draw[lin, dashed] (hij) -- (pri);
\end{tikzpicture}
\end{center}
\legende{Arbre généalogique simplifié : les générations et les liens de parenté.}
\end{popfigure}

\courssub{Champ 2 : les types de familles}

Le texte de la section 1 présentait plusieurs modèles familiaux. Retenez les adjectifs correspondants : ils s'accordent avec \esp{familia} (féminin singulier).

\begin{center}
\begin{tabularx}{\linewidth}{|X|X|X|}
\hline
\rowcolor{popblueL} Español & Français & Définition rapide \\
\hline
la familia nuclear & la famille nucléaire & père, mère et enfants \\
\hline
la familia monoparental & la famille monoparentale & un seul parent avec ses enfants \\
\hline
la familia extensa & la famille élargie & plusieurs générations sous un même toit \\
\hline
la familia reconstituida & la famille recomposée & après un divorce ou un remariage \\
\hline
la familia tradicional & la famille traditionnelle & modèle classique, souvent religieux \\
\hline
la familia numerosa & la famille nombreuse & avec beaucoup d'enfants \\
\hline
\end{tabularx}
\end{center}

\begin{definition}[La familia reconstituida]
Une \esp{familia reconstituida} est une famille recomposée, formée après un divorce ou un remariage. Elle comprend souvent des \esp{hijastros} (beaux-enfants), un \esp{padrastro} ou une \esp{madrastra}, et parfois des \esp{hermanastros} (demi-frères et demi-sœurs sans lien de sang).
\end{definition}

\begin{savaistu}[La famille élargie, une réalité partagée]
Au Cameroun comme dans les villages hispanophones d'Espagne, d'Amérique latine ou de Guinée équatoriale, la famille élargie inclut souvent oncles, tantes et cousins sous le même toit : c'est exactement la \esp{familia extensa}. Ce modèle favorise la solidarité, mais peut aussi multiplier les avis différents sur l'éducation des enfants.
\end{savaistu}

\courssub{Champ 3 : les relations et les valeurs}

Pour parler de l'éducation et de la vie en famille, voici les noms essentiels. Remarquez que beaucoup se terminent en \esp{-ción} ou \esp{-dad}, comme en français.

\begin{center}
\begin{tabularx}{\linewidth}{|X|X|X|}
\hline
\rowcolor{popblueL} Español & Français & Exemple \\
\hline
la autoridad & l'autorité & \esp{El padre ejerce su autoridad.} \\
\hline
el respeto & le respect & \esp{El respeto es fundamental.} \\
\hline
la obediencia & l'obéissance & \esp{La obediencia no lo es todo.} \\
\hline
la libertad & la liberté & \esp{Los jóvenes piden más libertad.} \\
\hline
la independencia & l'indépendance & \esp{Busca su independencia.} \\
\hline
la confianza & la confiance & \esp{La confianza se construye.} \\
\hline
la convivencia & la vie en commun & \esp{La convivencia exige paciencia.} \\
\hline
la educación & l'éducation & \esp{La educación empieza en casa.} \\
\hline
los valores & les valeurs & \esp{Transmitir valores a los hijos.} \\
\hline
\end{tabularx}
\end{center}

\begin{attention}[Attention aux terminaisons]
Les noms en \esp{-dad} (\esp{la libertad}, \esp{la autoridad}) sont toujours féminins, comme en français. Mais méfiez-vous : \esp{el respeto} est masculin alors que « le respect » finit aussi par une consonne en français ; et \esp{la educación} prend un accent écrit sur la ó. Ne confondez pas non plus \esp{educación} (éducation, formation morale) et \esp{enseñanza} (enseignement scolaire).
\end{attention}

\courssub{Champ 4 : le conflit et le dialogue}

C'est le cœur du thème de la leçon : les tensions entre générations et les moyens de les résoudre.

\begin{definition}[La brecha generacional]
La \esp{brecha generacional} (littéralement « la brèche générationnelle ») désigne le fossé qui sépare les générations : différences de valeurs, de goûts, d'usage des nouvelles technologies. On dit aussi \esp{el choque generacional} (le choc des générations).
\end{definition}

\begin{center}
\begin{tabularx}{\linewidth}{|X|X|}
\hline
\rowcolor{popblueL} Le conflit & Le dialogue et la paix \\
\hline
el conflicto (le conflit) & el diálogo (le dialogue) \\
\hline
la discusión (la dispute, le débat) & la escucha (l'écoute) \\
\hline
la pelea (la querelle) & la comprensión (la compréhension) \\
\hline
el malentendido (le malentendu) & la reconciliación (la réconciliation) \\
\hline
la brecha generacional (le fossé générationnel) & la solución (la solution) \\
\hline
\end{tabularx}
\end{center}

\begin{exemplebox}[Exemples traduits]
\esp{Hay una brecha generacional entre mi hermano y mi padre.} « Il y a un fossé générationnel entre mon frère et mon père. »

\esp{El malentendido duró una semana.} « Le malentendu a duré une semaine. »

\esp{Después de la pelea, hubo una reconciliación.} « Après la querelle, il y a eu une réconciliation. »

\esp{El diálogo es la mejor solución.} « Le dialogue est la meilleure solution. »
\end{exemplebox}

\courssub{Les verbes utiles}

\begin{center}
\begin{tabularx}{\linewidth}{|X|X|X|}
\hline
\rowcolor{popblueL} Verbo & Français & Particularité \\
\hline
criar & élever (des enfants) & \esp{Mis padres me criaron en Yaoundé.} \\
\hline
educar & éduquer & régulier \\
\hline
obedecer & obéir & \esp{obedezco} à la 1re personne \\
\hline
respetar & respecter & régulier \\
\hline
discutir & discuter, se disputer & pas seulement « débattre » \\
\hline
pelearse & se disputer, se battre & verbe pronominal \\
\hline
reconciliarse & se réconcilier & verbe pronominal \\
\hline
escuchar & écouter & \esp{escuchar a alguien} \\
\hline
comprender & comprendre & régulier \\
\hline
aconsejar & conseiller & \esp{aconsejar algo a alguien} \\
\hline
\end{tabularx}
\end{center}

\begin{propriete}[Trois particularités à retenir]
\begin{itemize}
\item \esp{Obedecer} est irrégulier à la première personne du présent de l'indicatif : \esp{yo obedezco}, mais \esp{tú obedeces}, \esp{él obedece} (et à tout le présent du subjonctif : \esp{que yo obedezca}).
\item \esp{Pelearse} et \esp{reconciliarse} sont pronominaux : \esp{me peleo con mi hermano} (je me dispute avec mon frère), \esp{se reconciliaron} (ils se sont réconciliés).
\item \esp{Discutir} est un demi-ami : il signifie souvent « se disputer », et pas seulement « discuter calmement ». \esp{Discutir con los padres} = se disputer avec ses parents.
\end{itemize}
\end{propriete}

\begin{textebox}[Un malentendido en casa]
Ayer por la noche, en casa de los Mbarga, hubo una discusión. Andrés, de dieciséis años, llegó tarde sin avisar. Su padre, muy enfadado, le dijo: « ¡No respetas la autoridad de esta casa! » Andrés respondió que sus amigos tienen más libertad y que él ya no es un niño. Su madre, mujer de gran paciencia, propuso un diálogo familiar después de la cena. La abuela Rosa, que vive con ellos, recordó que los valores cambian con el tiempo, pero que el respeto y la confianza siempre son necesarios. Al final, Andrés comprendió que avisar no es perder la independencia, sino respetar la convivencia. Padre e hijo se reconciliaron y decidieron escuchar más al otro. Fue un malentendido, no una guerra.

\textbf{Glossaire} : \esp{avisar} = prévenir ; \esp{enfadado} = en colère ; \esp{ya no es un niño} = il n'est plus un enfant ; \esp{al final} = finalement ; \esp{sino} = mais (après négation).
\end{textebox}

Questions de compréhension :
\begin{enumerate}
\item ¿Por qué discuten Andrés y su padre?
\item ¿Quién propone el diálogo familiar?
\item ¿Qué recuerda la abuela Rosa?
\item ¿Cuál es la solución final del conflicto?
\end{enumerate}

\begin{exoresolu}[Compléter avec le mot juste]
Énoncé. Complétez chaque phrase avec le mot du vocabulaire qui convient : \esp{brecha generacional}, \esp{madrastra}, \esp{parientes}, \esp{respeto}, \esp{reconciliarse}.
\begin{enumerate}
\item Después del divorcio de mi padre, su nueva esposa es mi \dots
\item Todos mis \dots{} vienen a la fiesta familiar : tíos, primos, sobrinos.
\item Entre los jóvenes y los ancianos existe a veces una \dots
\item Después de la pelea, los dos hermanos decidieron \dots
\item El \dots{} hacia los mayores es un valor importante.
\end{enumerate}
\tcblower
Solution détaillée.
\begin{enumerate}
\item \esp{madrastra} : la nouvelle épouse du père est la belle-mère (famille recomposée).
\item \esp{parientes} : il s'agit de la famille élargie (oncles, cousins, neveux), et non du père et de la mère.
\item \esp{brecha generacional} : le fossé entre jeunes et anciens.
\item \esp{reconciliarse} : verbe pronominal, infinitif après \esp{decidieron}.
\item \esp{respeto} : nom masculin, \esp{el respeto hacia los mayores}.
\end{enumerate}
\end{exoresolu}

\begin{exercice}[À vous de jouer]
Traduisez en espagnol :
\begin{enumerate}
\item Mes grands-parents vivent avec nous dans une famille élargie.
\item Il y a un conflit entre mon cousin et son beau-père.
\item Le dialogue est la meilleure solution pour se réconcilier.
\item Mes parents (père et mère) respectent ma liberté.
\end{enumerate}
\end{exercice}

\begin{corrige}[Exercice « À vous de jouer »]
\begin{enumerate}
\item \esp{Mis abuelos viven con nosotros en una familia extensa.}
\item \esp{Hay un conflicto entre mi primo y su padrastro.}
\item \esp{El diálogo es la mejor solución para reconciliarse.}
\item \esp{Mis padres respetan mi libertad.} (et non \esp{mis parientes} : faux ami !)
\end{enumerate}
\end{corrige}

\begin{aretenir}[L'essentiel du lexique]
\begin{itemize}
\item \esp{Los padres} = les parents (père et mère) ; \esp{los parientes} = les parents au sens large (faux ami).
\item Quatre types de familles : \esp{nuclear}, \esp{monoparental}, \esp{extensa}, \esp{reconstituida}.
\item Le conflit : \esp{conflicto}, \esp{discusión}, \esp{pelea}, \esp{malentendido}, \esp{brecha generacional}. La paix : \esp{diálogo}, \esp{escucha}, \esp{reconciliación}, \esp{solución}.
\item Verbes clés : \esp{criar}, \esp{obedecer} (\esp{obedezco}), \esp{pelearse}, \esp{reconciliarse}, \esp{aconsejar}.
\end{itemize}
\end{aretenir}

Ce vocabulaire maîtrisé, nous pouvons maintenant passer à la grammaire : comment \textbf{exprimer la cause} pour expliquer ces conflits familiaux.

\coursec{Grammaire 1 : exprimer la cause}

Dans la section précédente, nous avons enrichi notre vocabulaire : \esp{la brecha generacional}, \esp{la autoridad}, \esp{el respeto}, \esp{el conflicto}... Nous savons maintenant \emph{nommer} les tensions familiales. Mais pour analyser un texte ou rédiger un paragraphe, il ne suffit pas de nommer un conflit : il faut pouvoir en expliquer les \cle{causes}. Pourquoi les jeunes discutent-ils avec leurs parents ? Pourquoi le dialogue familial est-il parfois rompu ? C'est toute la grammaire de cette section : les \cle{connecteurs de cause} espagnols.

\courssub{Observation : la cause dans le texte}

Reprenons quelques phrases tirées de notre texte support \emph{La familia cambia} et de son univers.

\begin{exemplebox}[Phrases d'observation]
\esp{Los jóvenes discuten con sus padres porque pasan mucho tiempo en el teléfono.}\\
« Les jeunes se disputent avec leurs parents parce qu'ils passent beaucoup de temps au téléphone. »

\esp{Como no hay diálogo, hay conflictos.}\\
« Comme il n'y a pas de dialogue, il y a des conflits. »

\esp{Ya que no me escuchas, no te hablo.}\\
« Puisque tu ne m'écoutes pas, je ne te parle pas. »

\esp{Muchos abuelos se sienten solos a causa de la urbanización.}\\
« Beaucoup de grands-parents se sentent seuls à cause de l'urbanisation. »

\esp{Gracias al diálogo, la familia se reconcilió.}\\
« Grâce au dialogue, la famille s'est réconciliée. »
\end{exemplebox}

Observons attentivement. Deux grandes familles de constructions apparaissent :

\begin{itemize}
\item \esp{porque}, \esp{como}, \esp{ya que} sont suivis d'un \cle{verbe conjugué} : \esp{pasan}, \esp{hay}, \esp{escuchas}. Ce sont des \cle{conjonctions} : elles introduisent une proposition subordonnée de cause, exactement comme « parce que », « comme », « puisque » en français.
\item \esp{a causa de}, \esp{gracias a} sont suivis d'un \cle{nom} : \esp{la urbanización}, \esp{el diálogo}. Ce sont des \cle{locutions prépositives} : elles introduisent une cause nominale, comme « à cause de », « grâce à » en français.
\end{itemize}

C'est la grande distinction à retenir : \textbf{cause verbale} (avec conjonction) ou \textbf{cause nominale} (avec locution prépositive + nom).

\courssub{La cause verbale : porque, pues, ya que, como}

\begin{propriete}[Porque + indicatif]
\esp{Porque} + verbe conjugué à l'\cle{indicatif} répond à la question \esp{¿Por qué...?} (« Pourquoi... ? »). Il se place normalement \textbf{après} la proposition principale.

\esp{¿Por qué discuten los jóvenes? — Porque pasan mucho tiempo en el teléfono.}\\
« Pourquoi les jeunes se disputent-ils ? — Parce qu'ils passent beaucoup de temps au téléphone. »
\end{propriete}

\begin{attention}[Porque, por qué, porqué : trois mots à ne pas confondre]
Les francophones confondent souvent ces trois formes, pourtant bien distinctes :
\begin{itemize}
\item \esp{porque} (un seul mot, sans accent) = « parce que », conjonction de cause : \esp{No salgo porque llueve.} « Je ne sors pas parce qu'il pleut. »
\item \esp{por qué} (deux mots, accent sur \esp{qué}) = « pourquoi », interrogatif direct ou indirect : \esp{¿Por qué llegas tarde?} « Pourquoi arrives-tu en retard ? » ; \esp{No sé por qué está enfadado.} « Je ne sais pas pourquoi il est fâché. »
\item \esp{el porqué} (un seul mot, accentué, précédé de l'article) = « la cause, le pourquoi », nom masculin : \esp{No entiendo el porqué de este conflicto.} « Je ne comprends pas la cause de ce conflit. »
\end{itemize}
Et surtout : ne dites jamais \esp{*porque de} ! \esp{Porque} se construit directement avec le verbe, sans \esp{de}.
\end{attention}

\begin{propriete}[Pues et ya que : la cause justificative]
\esp{Pues} et \esp{ya que} introduisent une cause présentée comme une \cle{justification}, une évidence que le locuteur rappelle (« puisque », « car », « vu que »).
\begin{itemize}
\item \esp{Pues} se place toujours \textbf{à l'intérieur} de la phrase, après la proposition principale ; il ne commence jamais une phrase en tête de texte soutenu : \esp{No puedo ayudarte, pues tengo muchos deberes.} « Je ne peux pas t'aider, car j'ai beaucoup de devoirs. »
\item \esp{Ya que}, en revanche, peut se placer \textbf{en tête de phrase} : \esp{Ya que no me escuchas, no te hablo.} « Puisque tu ne m'écoutes pas, je ne te parle pas. »
\end{itemize}
\end{propriete}

\begin{propriete}[Como en tête de phrase]
\esp{Como} + verbe conjugué, placé \textbf{en tête de phrase} et suivi d'une virgule, équivaut à « comme / puisque ». C'est un emploi très fréquent à l'écrit.

\esp{Como no hay diálogo, hay conflictos.} « Comme il n'y a pas de dialogue, il y a des conflits. »

\esp{Como mis padres no entienden las redes sociales, discutimos a menudo.} « Comme mes parents ne comprennent pas les réseaux sociaux, nous discutons souvent. »

\textbf{Attention à la place} : en tête de phrase, \esp{como} = « puisque » ; à l'intérieur de la phrase, il garde son sens de comparaison (« comme ») ou de manière. Comparez : \esp{Trabaja como su padre.} « Il travaille comme son père. » (comparaison, pas de cause !)
\end{propriete}

\courssub{La cause nominale : gracias a, debido a, a causa de, por culpa de}

Quand la cause est un \cle{nom} (et non une phrase avec verbe conjugué), l'espagnol choisit la locution selon la \textbf{valeur} de la cause : positive, neutre ou négative. C'est une nuance que le français marque aussi, mais pas exactement de la même manière.

\begin{propriete}[Le choix dépend de la valeur de la cause]
\begin{itemize}
\item \esp{gracias a} + nom : cause \cle{positive} (ou ironique) : \esp{Gracias al diálogo, la familia se reconcilió.} « Grâce au dialogue, la famille s'est réconciliée. »
\item \esp{debido a} + nom : cause \cle{neutre}, registre plutôt soutenu : \esp{La reunión se aplazó debido a la lluvia.} « La réunion a été reportée en raison de la pluie. »
\item \esp{a causa de} + nom : cause \cle{négative} : \esp{Llegó tarde a causa del tráfico.} « Il est arrivé en retard à cause de la circulation. »
\item \esp{por culpa de} + nom : cause négative avec idée de \cle{reproche}, de responsabilité : \esp{Por culpa del teléfono, los hijos no hablan con sus padres.} « À cause du téléphone, les enfants ne parlent pas avec leurs parents. »
\end{itemize}
Ces locutions sont suivies d'un \textbf{nom} (ou d'un infinitif), jamais directement d'un verbe conjugué. Pour exprimer une cause verbale négative, on repasse à \esp{porque} : \esp{Llegó tarde porque había mucho tráfico.}
\end{propriete}

\begin{popfigure}
\begin{center}
\begin{tikzpicture}[
  box/.style={draw=popblue, thick, rounded corners, fill=popblueL, align=center, inner sep=6pt, font=\small},
  lab/.style={align=center, font=\small\bfseries, text=popdark}
]
\node[box, fill=popgreenL, draw=popgreen, align=center] (pos) at (0,2){\esp{gracias a} + nom\\ \esp{Gracias al diálogo...}\\ « Grâce au dialogue... »};
\node[box, align=center] (neu) at (0,0){\esp{debido a} + nom\\ \esp{...debido a la lluvia.}\\ « ...en raison de la pluie. »};
\node[box, fill=poppinkL, draw=poppink, align=center] (neg) at (0,-2.4){\esp{a causa de / por culpa de} + nom\\ \esp{...a causa del tráfico.}\\ « ...à cause de la circulation. »};
\node[lab, text=popgreen, align=center] at (-4.2,2){CAUSA\\ POSITIVA};
\node[lab, text=popblue, align=center] at (-4.2,0){CAUSA\\ NEUTRA};
\node[lab, text=poppink, align=center] at (-4.2,-2.4){CAUSA\\ NEGATIVA};
\draw[-{Stealth}, very thick, popgreen] (-3.1,2) -- (-1.9,2);
\draw[-{Stealth}, very thick, popblue] (-3.1,0) -- (-1.9,0);
\draw[-{Stealth}, very thick, poppink] (-3.1,-2.4) -- (-1.9,-2.4);
\end{tikzpicture}
\end{center}
\legende{Choisir la locution de cause nominale selon la valeur de la cause.}
\end{popfigure}

\begin{attention}[Les pièges des francophones]
\begin{itemize}
\item « À cause de » n'est pas toujours \esp{a causa de} ! Le français « à cause de » peut être neutre ; l'espagnol distingue strictement : \esp{gracias a} est \textbf{toujours positif} (ou ironique). On ne dit pas \esp{*Gracias al tráfico, llegó tarde} sauf ironie mordante.
\item Ne dites jamais \esp{*debido de} : c'est \esp{debido a}, avec la préposition \esp{a}.
\item Contraction obligatoire : \esp{a} + \esp{el} = \esp{al} ; \esp{de} + \esp{el} = \esp{del}. D'où \esp{gracias al diálogo}, \esp{a causa del tráfico}, \esp{debido al ruido}.
\item Ne traduisez pas « parce que » par \esp{*por causa de} devant un verbe : \esp{por causa de} n'existe pas ; devant un verbe conjugué, utilisez \esp{porque}.
\end{itemize}
\end{attention}

\courssub{Comparaison français / espagnol}

\begin{center}
\begin{tabularx}{\linewidth}{|X|X|X|}
\hline
\rowcolor{popblueL} Français & Español & Remarque \\
\hline
parce que + verbe & \esp{porque} + indicatif & répond à « pourquoi » \\
\hline
puisque / comme + verbe (en tête) & \esp{ya que} / \esp{como} + verbe & \esp{como} uniquement en tête de phrase \\
\hline
car (à l'intérieur) & \esp{pues} & jamais en tête de phrase soutenue \\
\hline
grâce à + nom & \esp{gracias a} + nom & cause positive \\
\hline
en raison de + nom & \esp{debido a} + nom & cause neutre, soutenu \\
\hline
à cause de + nom & \esp{a causa de / por culpa de} + nom & cause négative ; \esp{por culpa de} = reproche \\
\hline
pourquoi ? & \esp{¿por qué?} & deux mots, accent \\
\hline
la cause, le pourquoi & \esp{el porqué} & nom masculin \\
\hline
\end{tabularx}
\end{center}

\begin{aretenir}[Tableau récapitulatif des connecteurs de cause]
\begin{center}
\begin{tabularx}{\linewidth}{|l|X|l|X|}
\hline
\rowcolor{popblueL} Connecteur & Place & Registre & Ejemplo \\
\hline
\esp{porque} + verbe & après la principale & courant & \esp{Discuten porque pasan mucho tiempo en el teléfono.} « Ils se disputent parce qu'ils passent beaucoup de temps au téléphone. » \\
\hline
\esp{pues} + verbe & à l'intérieur & courant & \esp{No salgo, pues estoy cansado.} « Je ne sors pas, car je suis fatigué. » \\
\hline
\esp{ya que} + verbe & tête ou intérieur & soutenu & \esp{Ya que no me escuchas, no te hablo.} « Puisque tu ne m'écoutes pas, je ne te parle pas. » \\
\hline
\esp{como} + verbe & tête de phrase & écrit & \esp{Como no hay diálogo, hay conflictos.} « Comme il n'y a pas de dialogue, il y a des conflits. » \\
\hline
\esp{gracias a} + nom & tête ou intérieur & courant & \esp{Gracias al diálogo, la familia se reconcilió.} « Grâce au dialogue, la famille s'est réconciliée. » \\
\hline
\esp{debido a} + nom & tête ou intérieur & soutenu & \esp{Se aplazó debido a la lluvia.} « Reporté en raison de la pluie. » \\
\hline
\esp{a causa de} + nom & tête ou intérieur & courant & \esp{Llegó tarde a causa del tráfico.} « Il est arrivé en retard à cause de la circulation. » \\
\hline
\esp{por culpa de} + nom & tête ou intérieur & familier & \esp{Por culpa del teléfono, no hablamos.} « À cause du téléphone, nous ne parlons pas. » \\
\hline
\end{tabularx}
\end{center}
\end{aretenir}

\courssub{Méthode : choisir le bon connecteur de cause}

\begin{methode}[Quatre questions pour ne jamais se tromper]
\begin{enumerate}
\item Ma cause contient-elle un \textbf{verbe conjugué} ? $\rightarrow$ Oui : utilisez une conjonction (\esp{porque}, \esp{ya que}, \esp{como} en tête, \esp{pues} à l'intérieur). Non (c'est un simple nom) : passez à l'étape 2.
\item La cause est-elle \textbf{positive} ? $\rightarrow$ \esp{gracias a}.
\item La cause est-elle \textbf{négative} ? $\rightarrow$ \esp{a causa de} ; avec reproche : \esp{por culpa de}.
\item La cause est-elle \textbf{neutre} (simple constat) ? $\rightarrow$ \esp{debido a}.
\end{enumerate}
Vérifiez enfin les contractions : \esp{a + el = al}, \esp{de + el = del}.
\end{methode}

\courssub{Exercice résolu}

\begin{exoresolu}[Compléter avec le bon connecteur de cause]
Complétez avec : \esp{porque}, \esp{como}, \esp{ya que}, \esp{gracias a}, \esp{debido a}, \esp{a causa de}.

1. \esp{............ no hay respeto, hay conflictos en la familia.} (en tête de phrase, verbe conjugué)

2. \esp{Los abuelos están contentos ............ la visita de sus nietos.} (cause positive, nom)

3. \esp{María no fue a la escuela ............ estaba enferma.} (répond à « pourquoi », verbe)

4. \esp{El partido de fútbol se canceló ............ la lluvia.} (cause neutre, nom)

5. \esp{............ eres mi hijo, debes respetarme.} (en tête, justification)

6. \esp{Muchos jóvenes de Yaoundé llegaron tarde ............ los atascos.} (cause négative, nom)
\tcblower
\textbf{Solution détaillée.}

1. \esp{\textbf{Como} no hay respeto, hay conflictos en la familia.} — Verbe conjugué \esp{hay} en tête de phrase : seul \esp{como} (ou \esp{ya que}) convient ; ici \esp{como} est attendu. « Comme il n'y a pas de respect, il y a des conflits dans la famille. »

2. \esp{...contentos \textbf{gracias a} la visita de sus nietos.} — Cause nominale (\esp{la visita}) et positive (les grands-parents sont contents) : \esp{gracias a}. « Grâce à la visite de leurs petits-enfants. »

3. \esp{...a la escuela \textbf{porque} estaba enferma.} — Cause verbale qui répond à « pourquoi n'y est-elle pas allée ? » : \esp{porque}. « Parce qu'elle était malade. »

4. \esp{...se canceló \textbf{debido a} la lluvia.} — Cause nominale neutre, registre soutenu (annonce officielle) : \esp{debido a}. « En raison de la pluie. »

5. \esp{\textbf{Ya que} eres mi hijo, debes respetarme.} — Justification en tête de phrase : \esp{ya que}. « Puisque tu es mon fils, tu dois me respecter. »

6. \esp{...llegaron tarde \textbf{a causa de} los atascos.} — Cause nominale négative (\esp{los atascos}, « les embouteillages ») : \esp{a causa de}. « À cause des embouteillages. »
\end{exoresolu}

\begin{exercice}[À vous de jouer]
Traduisez en espagnol :
\begin{enumerate}
\item Les jeunes se disputent avec leurs parents parce qu'ils passent trop de temps sur les réseaux sociaux.
\item Comme il n'y a pas de dialogue, les conflits augmentent.
\item Grâce au respect mutuel, la famille vit en paix.
\item La réunion familiale a été annulée en raison du voyage du grand-père.
\item À cause du bruit, je ne peux pas faire mes devoirs.
\end{enumerate}
\end{exercice}

\begin{corrige}[Exercice « À vous de jouer »]
1. \esp{Los jóvenes discuten con sus padres porque pasan demasiado tiempo en las redes sociales.} (cause verbale $\rightarrow$ \esp{porque})

2. \esp{Como no hay diálogo, los conflictos aumentan.} (tête de phrase $\rightarrow$ \esp{como})

3. \esp{Gracias al respeto mutuo, la familia vive en paz.} (cause nominale positive ; contraction \esp{al})

4. \esp{La reunión familiar se canceló debido al viaje del abuelo.} (cause neutre ; \esp{debido al}, \esp{del abuelo})

5. \esp{A causa del ruido, no puedo hacer mis deberes.} (cause nominale négative ; contraction \esp{del})
\end{corrige}

\begin{savaistu}[Et dans la vraie famille hispanophone ?]
En Espagne comme en Amérique latine, la formule \esp{¡Porque lo digo yo!} (« Parce que je le dis ! ») est la réponse classique — et redoutée — des parents quand l'autorité remplace le dialogue. Les jeunes hispanophones lui répondent souvent par un \esp{¿Por qué?} insistant. Toute la grammaire de cette section tient dans ce petit duel : la question \esp{¿por qué?} appelle la réponse \esp{porque}... et le dialogue peut commencer !
\end{savaistu}

Maintenant que nous savons dire \emph{pourquoi} les conflits apparaissent, il reste à apprendre à en exprimer les \emph{conséquences} : ce sera l'objet de la section suivante, avec \esp{por eso}, \esp{entonces}, \esp{así que} et \esp{de modo que}.

\coursec{Grammaire 2 : exprimer la conséquence}

Dans la section précédente, nous avons appris à exprimer la \cle{cause} avec \esp{porque}, \esp{como}, \esp{pues} ou \esp{a causa de} : \esp{Hay conflictos porque los padres no hablan con sus hijos.} « Il y a des conflits parce que les parents ne parlent pas avec leurs enfants. » Mais on peut renverser la perspective : au lieu de remonter du résultat vers la cause, on part de la cause et on annonce le \cle{résultat}, la \cle{conséquence}. C'est exactement ce que fait notre phrase d'observation, tirée du texte support :

\begin{exemplebox}[Observation]
\esp{Los padres no hablan con sus hijos, por eso hay conflictos.}\\
« Les parents ne parlent pas avec leurs enfants, c'est pourquoi il y a des conflits. »
\end{exemplebox}

Ici, la cause (\esp{los padres no hablan con sus hijos}) est d'abord énoncée, puis \esp{por eso} introduit la conséquence (\esp{hay conflictos}). Toute cette section est consacrée aux \cle{connecteurs de conséquence} : \esp{por eso}, \esp{así que}, \esp{entonces}, \esp{de modo que}, \esp{tan... que}, \esp{tanto... que}.

\courssub{Observation : du texte à la règle}

Lisons d'abord ce court texte original, qui reprend le thème de la leçon et multiplie les expressions de la conséquence.

\begin{textebox}[La brecha generacional en casa de los Mbarga]
En casa de la familia Mbarga, en Yaoundé, el abuelo Pierre no entiende las nuevas tecnologías. Sus nietos pasan horas delante del teléfono, por eso el abuelo piensa que ya no respetan a los mayores. La madre trabaja todo el día en el mercado de Mokolo, así que apenas tiene tiempo para hablar con sus hijos. Los jóvenes hablan tan rápido y con tantas palabras nuevas que los padres no les entienden. Además, la familia se ha mudado del pueblo a la ciudad, de modo que los primos y los tíos viven lejos y la familia extensa se ha roto. El padre está muy cansado por la noche, entonces no cena con los niños. Hay tanto ruido en la casa que nadie escucha a nadie. Sin embargo, la abuela Rosa ha tenido una idea : cada domingo, toda la familia come junta sin teléfonos. Poco a poco, el diálogo vuelve : los abuelos cuentan historias del pueblo y los jóvenes explican Internet a los mayores. Así, la familia descubre que la brecha generacional no es un muro, sino un puente que hay que cruzar con paciencia y respeto.
\end{textebox}

\begin{description}
\item[\esp{la brecha generational}] le fossé entre les générations
\item[\esp{apenas}] à peine
\item[\esp{mudarse}] déménager
\item[\esp{los mayores}] les aînés, les personnes âgées
\item[\esp{un muro}] un mur
\item[\esp{un puente}] un pont
\item[\esp{cruzar}] traverser
\end{description}

\textbf{Questions de compréhension :}
\begin{enumerate}
\item ¿Por qué piensa el abuelo que los nietos no respetan a los mayores ?
\item ¿Qué consecuencia tiene el trabajo de la madre ?
\item ¿Cuál es la solución de la abuela Rosa ?
\end{enumerate}

Soulignons les connecteurs du texte : \esp{por eso}, \esp{así que}, \esp{de modo que}, \esp{entonces}, \esp{tan... que}, \esp{tanto... que}. Tous annoncent une \textbf{conséquence}, c'est-à-dire le résultat logique d'un fait exprimé juste avant. Observons leur place : ils se trouvent \textbf{au début de la deuxième proposition}, après une virgule ou un point-virgule.

\courssub{La règle générale}

\begin{propriete}[Règle de la conséquence]
Les connecteurs \esp{por eso} (c'est pourquoi), \esp{así que} (donc, alors), \esp{entonces} (donc, alors) et \esp{de modo que} / \esp{de manera que} (si bien que, de sorte que) introduisent le \textbf{résultat} d'une cause exprimée dans la proposition précédente. Ils sont suivis de l'\textbf{indicatif}, car le résultat est présenté comme un fait réel.\\
Structure : \textbf{CAUSE} + , + \textbf{connecteur} + \textbf{conséquence (indicatif)}.\\
\esp{La madre trabaja mucho, por eso está cansada.} « La mère travaille beaucoup, c'est pourquoi elle est fatiguée. »
\end{propriete}

\begin{popfigure}
\begin{tikzpicture}[node distance=0.6cm and 1.6cm,
causa/.style={rectangle, rounded corners, draw=popblue, fill=popblueL, align=center, text width=4.6cm, inner sep=6pt},
consec/.style={rectangle, rounded corners, draw=poporange, fill=poporangeL, align=center, text width=4.6cm, inner sep=6pt},
flecha/.style={-{Stealth[length=3mm]}, very thick, popdark}]
\node[causa, align=center] (c1){\textbf{CAUSA}\\ \esp{Los padres no hablan}\\\esp{con sus hijos.}\\ (porque...)};
\node[consec, right=of c1, align=center] (r1){\textbf{CONSECUENCIA}\\ \esp{por eso hay conflictos.}\\ (indicatif)};
\draw[flecha] (c1) -- node[above, font=\small\itshape]{por eso} (r1);
\node[causa, below=1.1cm of c1, align=center] (c2){\textbf{CAUSA}\\ \esp{No había diálogo}\\\esp{en la familia.}};
\node[consec, right=of c2, align=center] (r2){\textbf{CONSECUENCIA}\\ \esp{así que la familia}\\\esp{se separó.}};
\draw[flecha] (c2) -- node[above, font=\small\itshape]{así que} (r2);
\end{tikzpicture}
\legende{De la cause à la conséquence : le connecteur annonce le résultat.}
\end{popfigure}

\courssub{Étude détaillée des connecteurs}

\textbf{1. \esp{Por eso} : la conséquence logique}

\esp{Por eso} signifie littéralement « pour cela », donc « c'est pourquoi, c'est pour cette raison ». C'est le connecteur de conséquence le plus neutre et le plus utile à l'écrit. Il se place en début de deuxième proposition, le plus souvent après une virgule ou un point-virgule, parfois après un point.

\begin{itemize}
\item \esp{Mi hermano pasa todo el día en Internet, por eso sus notas bajan.} « Mon frère passe toute la journée sur Internet, c'est pourquoi ses notes baissent. »
\item \esp{Los jóvenes viven en la ciudad ; por eso olvidan las costumbres del pueblo.} « Les jeunes vivent en ville ; c'est pourquoi ils oublient les coutumes du village. »
\item \esp{El abuelo no entiende el francés. Por eso prefiere hablar en ewondo.} « Le grand-père ne comprend pas le français. C'est pourquoi il préfère parler en ewondo. »
\end{itemize}

\textbf{2. \esp{Así que} et \esp{entonces} : la conséquence orale}

\esp{Así que} (« donc, alors ») est extrêmement fréquent à l'oral, dans les conversations familiales. \esp{Entonces} (« alors, donc ») est un peu plus soutenu et s'emploie aussi bien à l'oral qu'à l'écrit. Les deux sont suivis de l'indicatif.

\begin{itemize}
\item \esp{No había diálogo, así que la familia se separó.} « Il n'y avait pas de dialogue, donc la famille s'est séparée. »
\item \esp{Hoy es domingo, así que comemos todos juntos.} « Aujourd'hui c'est dimanche, alors nous mangeons tous ensemble. »
\item \esp{El padre llega tarde, entonces no ve a sus hijos.} « Le père rentre tard, alors il ne voit pas ses enfants. »
\end{itemize}

\begin{attention}[Orthographe de \esp{así que}]
\esp{Así} prend un accent sur le \esp{í} (mot aigu terminé par une voyelle) : on écrit \esp{así que}, jamais \esp{asi que}. Et ne confondez pas \esp{así que} (donc) avec \esp{aunque} (bien que) : deux connecteurs au sens opposé !
\end{attention}

\textbf{3. \esp{De modo que} / \esp{de manera que} : la conséquence écrite}

Ces deux connecteurs, équivalents à « si bien que, de sorte que », appartiennent plutôt à la langue écrite et soignée. Suivis de l'\textbf{indicatif}, ils expriment un résultat réel :

\begin{itemize}
\item \esp{La familia se ha mudado a Douala, de modo que los abuelos se quedaron solos en el pueblo.} « La famille a déménagé à Douala, si bien que les grands-parents sont restés seuls au village. »
\item \esp{Hablaban dos lenguas distintas, de manera que no se entendían.} « Ils parlaient deux langues différentes, de sorte qu'ils ne se comprenaient pas. »
\end{itemize}

\begin{attention}[Indicatif ou subjonctif après \esp{de modo que} ?]
Après \esp{de modo que} / \esp{de manera que}, l'espagnol choisit le mode selon le sens : \textbf{indicatif} si l'on exprime un \textbf{résultat réel} (la conséquence, notre leçon), \textbf{subjonctif} si l'on exprime un \textbf{but} (la finalité, hors programme détaillé en Première, mais il faut le savoir repérer). Comparez :\\
\esp{No estudió, de modo que suspendió el examen.} « Il n'a pas étudié, si bien qu'il a raté l'examen » (résultat, indicatif).\\
\esp{Te lo explico de modo que lo entiendas.} « Je te l'explique de sorte que tu comprennes » (but, subjonctif).\\
Au niveau Première, retenez : conséquence = indicatif.
\end{attention}

\textbf{4. \esp{Tan... que} et \esp{tanto... que} : la conséquence d'intensité}

Quand la cause est une \textbf{intensité} (« tellement... que »), l'espagnol utilise une construction en deux morceaux :

\begin{itemize}
\item \esp{tan} + adjectif ou adverbe + \esp{que} : \esp{Habla tan rápido que no le entiendo.} « Il parle si vite que je ne le comprends pas. »
\item \esp{tanto / tanta / tantos / tantas} + nom + \esp{que} : \esp{Hay tanto ruido que nadie escucha.} « Il y a tellement de bruit que personne n'écoute. » ; \esp{Tiene tantos hijos que no puede pagar todas las escuelas.} « Il a tant d'enfants qu'il ne peut pas payer toutes les écoles. »
\item \esp{tanto} seul après un verbe : \esp{Los adolescentes discuten tanto que los padres se cansan.} « Les adolescents discutent tellement que les parents se fatiguent. »
\end{itemize}

Attention à l'accord de \esp{tanto} avec le nom : \esp{tanta paciencia} (féminin singulier), \esp{tantos problemas} (masculin pluriel).

\courssub{Comparaison français / espagnol}

\begin{center}
\begin{tabularx}{\linewidth}{|X|X|X|}
\hline
\rowcolor{popblueL} Français & Español & Remarque \\
\hline
c'est pourquoi & por eso & neutre, écrit et oral \\
\hline
donc, alors (oral) & así que & très fréquent à l'oral \\
\hline
donc, alors & entonces & oral et écrit \\
\hline
si bien que, de sorte que & de modo que / de manera que & écrit ; indicatif = résultat \\
\hline
si / tellement... que (+ adj.) & tan... que & \esp{tan rápido que} \\
\hline
tellement de... que (+ nom) & tanto/a/os/as... que & accord avec le nom \\
\hline
par conséquent, c'est pourquoi & por lo tanto / por consiguiente & soutenu, dissertation \\
\hline
\end{tabularx}
\end{center}

Le français « donc » a donc deux traductions principales : \esp{por eso} (plus écrit) et \esp{entonces} ou \esp{así que} (plus oral). Dans une rédaction de Première, variez les connecteurs pour enrichir votre expression.

\begin{attention}[Le piège du Bac : \esp{por eso} contre \esp{porque}]
Ne confondez jamais \esp{por eso} (c'est pourquoi → conséquence) et \esp{porque} (parce que → cause). C'est l'erreur classique des francophones, sanctionnée au Bac :\\
\esp{Hay conflictos \textbf{porque} no hay diálogo.} « Il y a des conflits \textbf{parce qu'}il n'y a pas de dialogue » (cause).\\
\esp{No hay diálogo, \textbf{por eso} hay conflictos.} « Il n'y a pas de dialogue, \textbf{c'est pourquoi} il y a des conflits » (conséquence).\\
Astuce : \esp{porque} répond à la question « ¿por qué ? » (pourquoi ?) ; \esp{por eso} répond à la question « ¿qué pasa entonces ? » (que se passe-t-il alors ?).
\end{attention}

\begin{exoresolu}[Reliez par le connecteur de conséquence adapté]
Énoncé. Reliez les deux propositions avec \esp{por eso}, \esp{así que}, \esp{de modo que} ou \esp{tan... que}, puis traduisez.\\
1. Los padres trabajan todo el día. / No ven a sus hijos.\\
2. La abuela habla muy despacio. / Todos la entienden.\\
3. La familia se mudó a la ciudad. / Perdió el contacto con el pueblo.\\
4. Mi hermano pequeño grita mucho. / No puedo estudiar.
\tcblower
Solution détaillée.\\
1. \esp{Los padres trabajan todo el día, por eso no ven a sus hijos.} « Les parents travaillent toute la journée, c'est pourquoi ils ne voient pas leurs enfants. » (conséquence neutre, écrit).\\
2. \esp{La abuela habla tan despacio que todos la entienden.} « La grand-mère parle si lentement que tout le monde la comprend. » (intensité : \esp{tan} + adverbe + \esp{que}).\\
3. \esp{La familia se mudó a la ciudad, de modo que perdió el contacto con el pueblo.} « La famille a déménagé en ville, si bien qu'elle a perdu le contact avec le village. » (style écrit, indicatif car résultat réel).\\
4. \esp{Mi hermano pequeño grita mucho, así que no puedo estudiar.} « Mon petit frère crie beaucoup, alors je ne peux pas étudier. » (registre oral, conversation).
\end{exoresolu}

\begin{exercice}[À vous]
Transformez les phrases de cause en phrases de conséquence, puis traduisez-les.\\
1. \esp{Hay conflictos porque los padres no escuchan a los hijos.}\\
2. \esp{Los jóvenes no respetan a los mayores porque no hablan con ellos.}\\
3. \esp{La familia está unida porque cena junta cada domingo.}
\end{exercice}

\begin{corrige}[Exercice « À vous »]
1. \esp{Los padres no escuchan a los hijos, por eso hay conflictos.} « Les parents n'écoutent pas les enfants, c'est pourquoi il y a des conflits. »\\
2. \esp{Los jóvenes no hablan con los mayores, por eso no los respetan.} « Les jeunes ne parlent pas avec les aînés, c'est pourquoi ils ne les respectent pas. » (attention : \esp{los} remplace \esp{a los mayores}).\\
3. \esp{La familia cena junta cada domingo, por eso está unida.} « La famille dîne ensemble chaque dimanche, c'est pourquoi elle est unie. »
\end{corrige}

\begin{aretenir}[L'essentiel]
\begin{itemize}
\item La conséquence s'exprime avec \esp{por eso} (c'est pourquoi), \esp{así que} et \esp{entonces} (donc, oral), \esp{de modo que} / \esp{de manera que} (si bien que, écrit).
\item Ces connecteurs se placent en début de deuxième proposition, après une virgule ou un point-virgule, et sont suivis de l'\textbf{indicatif}.
\item L'intensité s'exprime avec \esp{tan} + adjectif/adverbe + \esp{que} et \esp{tanto/a/os/as} + nom + \esp{que}.
\item \esp{Porque} = cause ; \esp{por eso} = conséquence. Ne les inversez jamais.
\end{itemize}
\end{aretenir}

\begin{savaistu}[« Justifique tu respuesta » : le mode d'emploi du Bac]
Au Bac, les questions de compréhension demandent souvent de \emph{justifier} sa réponse : \esp{Justifica tu respuesta con frases del texto.} C'est l'occasion rêvée d'employer \esp{porque} et \esp{por eso} dans une réponse rédigée : \esp{El abuelo está triste porque sus nietos no le escuchan ; por eso piensa que no respetan a los mayores.} Une phrase qui combine cause et conséquence montre à l'examinateur une vraie maîtrise de l'argumentation... et rapporte des points précieux !
\end{savaistu}

\coursec{Grammaire 3 : exprimer le conseil}

Dans la section précédente, nous avons appris à exprimer la \cle{conséquence} (\esp{por eso}, \esp{de modo que}, \esp{entonces}) : \esp{Los jóvenes pasan mucho tiempo en el móvil, por eso los padres se preocupan.} Mais une fois le problème identifié et ses conséquences décrites, reste l'essentiel : que faire ? Comment proposer des \cle{solutions} au conflit intergénérationnel ? C'est le rôle de l'expression du \cle{conseil}, dernière étape de notre boîte à outils grammaticale et compétence clé pour l'expression écrite au Bac.

\courssub{Observation : partir des exemples}

Lisons d'abord ce court dialogue entre une mère et sa fille à Yaoundé.

\begin{textebox}[Diálogo : un consejo de madre]
--- Hija, deberías escuchar a tus padres. Tenemos experiencia de la vida.\\
--- Pero mamá, ¡vosotros no entendéis nada de las redes sociales!\\
--- Quizás, pero hay que dialogar en la familia. Tienes que respetar a tus abuelos también.\\
--- Está bien, mamá. ¿Por qué no hablamos todos juntos esta noche?\\
--- Buena idea. Escucha a tu abuela: ella siempre da buenos consejos. Es importante que dialoguéis con ella.
\end{textebox}

\textbf{Glossaire :} \esp{las redes sociales} = les réseaux sociaux ; \esp{quizás} = peut-être ; \esp{esta noche} = ce soir ; \esp{dar consejos} = donner des conseils ; \esp{vosotros} = vous (forme péninsulaire, équivalent à \esp{ustedes} en Amérique latine et en Guinée équatoriale).

\textbf{Questions d'observation :}
\begin{enumerate}
\item Relevez toutes les phrases qui expriment un conseil ou une obligation.
\item Quels verbes ou expressions reviennent avant l'infinitif ?
\item Quelle forme verbale suit toujours \esp{deberías}, \esp{tienes que}, \esp{hay que} ?
\end{enumerate}

On observe quatre stratégies : \esp{deberías + infinitif} (conseil personnel atténué), \esp{tienes que + infinitif} (obligation), \esp{hay que + infinitif} (règle générale impersonnelle), et l'\cle{impératif} (\esp{escucha}, conseil direct). Étudions-les une à une.

\courssub{Deber + infinitif : le conseil personnel}

\begin{propriete}[Deber + infinitif : forme et emploi]
\textbf{Forme :} verbe \esp{deber} conjugué + \textbf{infinitif}.\\
\textbf{Emploi :} exprime le \cle{conseil} ou le devoir moral adressé à une personne précise.
\begin{itemize}
\item Au \textbf{présent} (\esp{debo, debes, debe...}) : devoir assez fort, proche de « tu dois ».
\item Au \textbf{conditionnel} (\esp{debería, deberías, debería...}) : conseil \textbf{atténué}, poli, proche de « tu devrais ». C'est la forme la plus utilisée pour conseiller sans imposer.
\end{itemize}
\esp{Deberías escuchar a tus padres.} « Tu devrais écouter tes parents. »\\
\esp{Debes ayudar en casa.} « Tu dois aider à la maison. »
\end{propriete}

\begin{popfigure}
\begin{center}
\begin{tikzpicture}
\node[draw=popblue, fill=popblueL, rounded corners, align=center, font=\bfseries] (titre) at (0,0) {DEBER + infinitif};
\node[draw=popgreen, fill=popgreenL, rounded corners, align=center] (pres) at (-4.2,-2.6) {\textbf{Présent} (= tu dois)\\[2pt]
yo \textbf{debo}\\ tú \textbf{debes}\\ él/ella \textbf{debe}\\ nosotros \textbf{debemos}\\ vosotros \textbf{debéis}\\ ellos \textbf{deben}};
\node[draw=poporange, fill=poporangeL, rounded corners, align=center] (cond) at (4.2,-2.6) {\textbf{Conditionnel} (= tu devrais)\\[2pt]
yo \textbf{debería}\\ tú \textbf{deberías}\\ él/ella \textbf{debería}\\ nosotros \textbf{deberíamos}\\ vosotros \textbf{deberíais}\\ ellos \textbf{deberían}};
\draw[-{Stealth}, thick, popblue] (titre) -- (pres);
\draw[-{Stealth}, thick, popblue] (titre) -- (cond);
\node[align=center, font=\itshape] at (0,-5.6) {Devoir fort : \esp{Debes respetar.} \quad Conseil doux : \esp{Deberías respetar.}};
\end{tikzpicture}
\legende{Conjugaison de \esp{deber} au présent (devoir) et au conditionnel (conseil atténué).}
\end{center}
\end{popfigure}

\begin{attention}[« Tu devrais » ≠ « tu dois »]
En français, « tu devrais » (conditionnel) est un conseil poli ; « tu dois » (présent) est plus direct. L'espagnol fait exactement la même distinction : \esp{deberías} = tu devrais (conseil) ; \esp{debes} = tu dois (devoir). Au Bac, pour \emph{proposer des solutions} sans paraître autoritaire, préférez le \textbf{conditionnel} : \esp{Los padres deberían hablar más con sus hijos.} « Les parents devraient parler davantage avec leurs enfants. »
\end{attention}

\courssub{Tener que + infinitif : l'obligation}

\begin{propriete}[Tener que + infinitif]
\textbf{Forme :} \esp{tener} conjugué + \esp{que} + \textbf{infinitif}.\\
\textbf{Emploi :} exprime une \cle{obligation} personnelle, plus forte que le simple conseil : on n'a pas le choix.\\
\esp{Tienes que respetar a tus abuelos.} « Tu dois respecter tes grands-parents. »\\
\esp{Tenemos que ayudar a nuestros padres.} « Nous devons aider nos parents. »
\end{propriete}

Rappel de la conjugaison de \esp{tener} (irrégulier) au présent : \esp{tengo, tienes, tiene, tenemos, tenéis, tienen}. Attention à l'alternance \textbf{e → ie} (\esp{tengo} excepté).

\begin{center}
\begin{tabularx}{\linewidth}{|X|X|X|}
\hline
\rowcolor{popblueL} Conseil (atténué) & Obligation (forte) & Remarque \\
\hline
\esp{Deberías estudiar más.} & \esp{Tienes que estudiar más.} & Tu devrais / tu dois étudier plus. \\
\hline
\esp{Deberíamos dialogar.} & \esp{Tenemos que dialogar.} & Nous devrions / nous devons dialoguer. \\
\hline
\end{tabularx}
\end{center}

\courssub{Hay que + infinitif : le conseil impersonnel}

\begin{propriete}[Hay que + infinitif]
\textbf{Forme :} \esp{hay que} (invariable) + \textbf{infinitif}.\\
\textbf{Emploi :} exprime une règle générale, une nécessité qui vaut pour tout le monde, sans désigner personne : « il faut ».\\
\esp{Hay que dialogar.} « Il faut dialoguer. »\\
\esp{Hay que escuchar antes de juzgar.} « Il faut écouter avant de juger. »
\end{propriete}

\begin{attention}[Pièges des francophones avec « hay que »]
\begin{itemize}
\item \esp{Hay que} est \textbf{impersonnel et invariable} : jamais \esp{*han que}, jamais \esp{*hay que a}. On ne conjugue pas \esp{haber} ici.
\item Après \esp{deber}, \esp{tener que} et \esp{hay que}, le verbe est \textbf{toujours à l'infinitif} : \esp{Deberías escuchar}, jamais \esp{*deberías escuchas}.
\item Ne confondez pas \esp{hay que} (il faut) et \esp{hay} (il y a) : \esp{Hay un problema. Hay que resolverlo.} « Il y a un problème. Il faut le résoudre. »
\end{itemize}
\end{attention}

\courssub{Formes directes du conseil : impératif et « ¿Por qué no...? »}

Pour conseiller de façon plus directe et vivante (à l'oral, dans un dialogue, dans une lettre à un ami), l'espagnol utilise l'impératif et la suggestion interrogative.

\begin{propriete}[L'impératif affirmatif (tú) régulier et « ¿Por qué no...? »]
\textbf{1. L'impératif affirmatif} à la 2\textsuperscript{e} personne du singulier (\esp{tú}) : pour les verbes réguliers, il reprend la 3\textsuperscript{e} personne du singulier du présent de l'indicatif :\\
\esp{escuchar → escucha} ; \esp{hablar → habla} ; \esp{respetar → respeta} ; \esp{comer → come} ; \esp{escribir → escribe}.\\
\esp{Escucha a tu abuela, tiene experiencia.} « Écoute ta grand-mère, elle a de l'expérience. »\\
\textbf{2. La tournure interrogative} \esp{¿Por qué no + présent ?} : conseil sous forme de suggestion aimable.\\
\esp{¿Por qué no hablas con tu madre?} « Pourquoi ne parles-tu pas avec ta mère ? »
\end{propriete}

\begin{propriete}[Impératifs affirmatifs irréguliers (tú) — à connaître par cœur]
\begin{center}
\begin{tabular}{|l|l|l|}
\hline
\rowcolor{popblueL} Infinitif & Impératif (tú) & Traduction \\
\hline
\esp{tener} & \esp{ten} & Aie / Tiens \\
\hline
\esp{poner} & \esp{pon} & Pose / Mets \\
\hline
\esp{salir} & \esp{sal} & Sors \\
\hline
\esp{venir} & \esp{ven} & Viens \\
\hline
\esp{decir} & \esp{di} & Dis \\
\hline
\esp{hacer} & \esp{haz} & Fais \\
\hline
\esp{ir} & \esp{ve} & Va \\
\hline
\esp{ser} & \esp{sé} & Sois \\
\hline
\end{tabular}
\end{center}
\esp{Ten paciencia con tu hermano.} « Sois patient avec ton frère. » \quad \esp{¡Ven a comer!} « Viens manger ! »
\end{propriete}

\begin{exemplebox}[Exemples traduits]
\begin{itemize}
\item \esp{Deberías pasar menos tiempo en el móvil.} « Tu devrais passer moins de temps sur le portable. »
\item \esp{Hay que respetar la autoridad de los padres.} « Il faut respecter l'autorité des parents. »
\item \esp{Escucha a tu abuela, tiene experiencia.} « Écoute ta grand-mère, elle a de l'expérience. »
\item \esp{¿Por qué no pasas el domingo con tu familia?} « Pourquoi ne passes-tu pas le dimanche avec ta famille ? »
\item \esp{Tenéis que comprender a vuestros padres también.} « Vous devez comprendre vos parents aussi. »
\item \esp{¡Ten cuidado!} « Fais attention ! » (irrégulier)
\end{itemize}
\end{exemplebox}

\courssub{Pour aller plus loin (complément Bac) : « Es importante que + subjonctif »}

\begin{savaistu}[Le subjonctif : simple reconnaissance]
Après une expression impersonnelle d'opinion ou de nécessité (\esp{es importante que}, \esp{es necesario que}, \esp{es bueno que}, \esp{es fundamental que}), l'espagnol emploie le \textbf{subjonctif} :\\
\esp{Es importante que dialoguéis.} « Il est important que vous dialoguiez. »\\
\esp{Es necesario que los jóvenes respeten a los mayores.} « Il est nécessaire que les jeunes respectent les aînés. »\\
Au niveau de la 1\textsuperscript{re}, on vous demande seulement de \textbf{reconnaître} cette structure dans un texte, pas encore de la manipuler librement. Notez que le français utilise aussi le subjonctif : « il faut que tu \emph{viennes} ».
\end{savaistu}

\courssub{Tableau récapitulatif : les outils du conseil}

\begin{center}
\begin{tabularx}{\linewidth}{|X|X|X|}
\hline
\rowcolor{popblueL} Structure & Emploi & Exemple traduit \\
\hline
\esp{debería(s) + inf.} & conseil atténué & \esp{Deberías hablar con él.} = Tu devrais lui parler. \\
\hline
\esp{debe(s) + inf.} & devoir & \esp{Debes estudiar.} = Tu dois étudier. \\
\hline
\esp{tener que + inf.} & obligation forte & \esp{Tienes que volver temprano.} = Tu dois rentrer tôt. \\
\hline
\esp{hay que + inf.} & règle impersonnelle & \esp{Hay que dialogar.} = Il faut dialoguer. \\
\hline
impératif (tú) & conseil direct & \esp{¡Habla con tu padre!} = Parle à ton père ! \\
\hline
\esp{¿Por qué no + prés. ?} & suggestion & \esp{¿Por qué no lo intentas?} = Pourquoi n'essaies-tu pas ? \\
\hline
\esp{es importante que + subj.} & nécessité (reconnaissance Bac) & \esp{Es importante que escuches.} = Il est important que tu écoutes. \\
\hline
\end{tabularx}
\end{center}

\courssub{Exercice résolu}

\begin{exoresolu}[Conseils à un ami]
Énoncé : Votre ami Carlos est en conflit avec son père parce qu'il passe trop de temps sur son téléphone. Complétez les conseils avec la structure indiquée, puis traduisez.\\
1. (deber, conditionnel) \esp{Tú \ldots\ pasar menos tiempo en el móvil.}\\
2. (hay que) \esp{\ldots\ respetar a los padres.}\\
3. (impératif de \esp{hablar}) \esp{¡\ldots\ con tu padre esta noche!}\\
4. (tener que, tú) \esp{\ldots\ ayudar en casa.}
\tcblower
Solution détaillée :\\
1. \esp{deberías} : conditionnel de \esp{deber}, 2\textsuperscript{e} pers. du sg. → \esp{Tú deberías pasar menos tiempo en el móvil.} « Tu devrais passer moins de temps sur le portable. »\\
2. \esp{Hay que} : invariable, impersonnel → \esp{Hay que respetar a los padres.} « Il faut respecter les parents. »\\
3. \esp{Habla} : impératif affirmatif de \esp{hablar} (3\textsuperscript{e} pers. du sg. du présent) → \esp{¡Habla con tu padre esta noche!} « Parle avec ton père ce soir ! »\\
4. \esp{Tienes que} : présent de \esp{tener} (e → ie) + \esp{que} → \esp{Tienes que ayudar en casa.} « Tu dois aider à la maison. »\\
Vérification : après chaque structure, le verbe est bien à l'\textbf{infinitif} (\esp{pasar, respetar, ayudar}) ; seul l'impératif est une forme conjuguée autonome.
\end{exoresolu}

\courssub{Méthode : rédiger un paragraphe de solutions au Bac}

\begin{methode}[Rédiger un paragraphe de solutions aux conflits familiaux]
\begin{enumerate}
\item \textbf{Annoncer le problème} avec une cause (\esp{porque}, \esp{como}) : \esp{Hay conflictos en la familia porque los jóvenes y los padres no se comprenden.}
\item \textbf{Proposer deux ou trois conseils} en variant les structures : un conseil personnel (\esp{Los jóvenes deberían escuchar a sus padres}), une règle générale (\esp{Hay que dialogar con respeto}), éventuellement un impératif ou une suggestion (\esp{¿Por qué no coméis juntos los domingos?}).
\item \textbf{Conclure par la conséquence positive} avec les connecteurs de la section précédente (\esp{por eso}, \esp{de modo que}, \esp{así}) : \esp{De modo que, con diálogo y respeto, la familia será más unida.}
\end{enumerate}
Modèle complet : \esp{En muchas familias hay conflictos porque los padres y los hijos tienen valores diferentes. Los jóvenes deberían escuchar a sus mayores y los padres deberían comprender la nueva realidad de sus hijos. Además, hay que dialogar con calma y con respeto. Por eso, el diálogo es la mejor solución para una familia unida.}
\end{methode}

\begin{aretenir}[L'essentiel de la section]
Pour exprimer le conseil en espagnol : \esp{debería(s) + infinitif} = conseil atténué (« tu devrais ») ; \esp{debe(s) + infinitif} = devoir ; \esp{tener que + infinitif} = obligation (« tu dois ») ; \esp{hay que + infinitif} = règle impersonnelle invariable (« il faut ») ; l'impératif (\esp{escucha, habla, respeta, ten, ven, di, haz, ve, sé}) et \esp{¿Por qué no + présent ?} pour le conseil direct. Après toutes ces structures (sauf l'impératif), le verbe est toujours à l'infinitif. Au Bac, combinez cause (\esp{porque}) + conseils (\esp{deberías}, \esp{hay que}) + conséquence (\esp{por eso}) pour un paragraphe de solutions réussi.
\end{aretenir}

\coursec{Méthode du commentaire et commentaire modèle}

Dans les sections précédentes, tu as appris à exprimer la cause (\esp{porque}, \esp{como}, \esp{a causa de}), la conséquence (\esp{por eso}, \esp{entonces}, \esp{así que}) et le conseil (\esp{deberías}, \esp{hay que}, \esp{es importante}). Ces outils ne sont pas de simples exercices de grammaire : ce sont les \cle{briques} qui permettent de construire un texte organisé. Cette dernière section te montre comment les assembler dans l'exercice roi de la classe de Première : le \cle{commentaire de texte} (\esp{comentario de texto}).

\courssub{Qu'est-ce qu'un commentaire de texte ?}

\begin{definition}[Le comentario de texto]
Le \cle{commentaire de texte} est un exercice d'expression écrite dans lequel l'élève présente un texte, dégage son idée principale, l'analyse en plusieurs \cle{axes} (\esp{ejes}) et conclut. Contrairement au résumé, il ne s'agit pas seulement de répéter le texte, mais de l'\emph{expliquer} et de donner son avis argumenté.
\end{definition}

Un bon commentaire suit toujours la même architecture, en forme d'\emph{entonnoir} : on part du général (la présentation du texte) pour aller vers le précis (l'analyse), puis on revient au général (la conclusion).

\begin{popfigure}
\begin{tikzpicture}[font=\small]
\fill[popblue!25] (0,4.4) -- (10,4.4) -- (9.2,3.4) -- (0.8,3.4) -- cycle;
\node[align=center] at (5,3.9) {\textbf{Introducción} : presentación del texto + problemática};
\fill[popgreen!25] (1.2,3.2) -- (8.8,3.2) -- (8.2,2.2) -- (1.8,2.2) -- cycle;
\node[align=center] at (5,2.7) {\textbf{Eje 1} : los cambios de la familia};
\fill[poporange!30] (2,2.0) -- (8,2.0) -- (7.4,1.0) -- (2.6,1.0) -- cycle;
\node[align=center] at (5,1.5) {\textbf{Eje 2} : los conflictos entre generaciones};
\fill[poppurple!20] (2.8,0.8) -- (7.2,0.8) -- (6.6,-0.2) -- (3.4,-0.2) -- cycle;
\node[align=center] at (5,0.3) {\textbf{Eje 3} : las soluciones};
\fill[popgold!30] (3.6,-0.4) -- (6.4,-0.4) -- (5.9,-1.4) -- (4.1,-1.4) -- cycle;
\node[align=center] at (5,-0.9) {\textbf{Conclusión}};
\end{tikzpicture}
\legende{La structure en entonnoir du commentaire de texte : du général vers le précis, puis retour au général.}
\end{popfigure}

\courssub{La méthode en cinq étapes}

\begin{methode}[Comment réussir un commentaire de texte]
\begin{enumerate}
\item \textbf{Lire le texte deux fois.} La première fois pour comprendre le sens général ; la deuxième fois en soulignant les mots-clés (\esp{familia}, \esp{conflicto}, \esp{diálogo}...) et les connecteurs (\esp{porque}, \esp{por eso}, \esp{sin embargo}).
\item \textbf{Rédiger l'introduction.} Présente le texte (type, thème, source éventuelle) et formule une \cle{problématique}, c'est-à-dire la question à laquelle ton commentaire va répondre. Exemple : \esp{¿Cómo ha cambiado la familia y por qué hay conflictos entre padres e hijos?} « Comment la famille a-t-elle changé et pourquoi y a-t-il des conflits entre parents et enfants ? »
\item \textbf{Développer deux ou trois axes.} Pour chaque axe : \textbf{cite} le texte en espagnol entre guillemets, \textbf{traduis} brièvement la citation (dans ton brouillon ou mentalement), \textbf{analyse}-la en espagnol (que montre-t-elle ? pourquoi est-ce important ?), puis \textbf{enchaîne} avec un connecteur (\esp{además}, \esp{sin embargo}, \esp{por eso}).
\item \textbf{Rédiger la conclusion.} Résume l'idée principale et donne ton avis personnel ou une ouverture (\esp{En mi opinión...}, \esp{Para terminar...}).
\item \textbf{Relire.} Vérifie les accords, les accents et la présence des connecteurs de cause, de conséquence et de conseil.
\end{enumerate}
\end{methode}

\begin{aretenir}[La règle d'or : citation + traduction + analyse]
Chaque citation doit suivre le schéma en trois temps :
\begin{itemize}
\item \textbf{Citer} (en espagnol) : \esp{El texto dice: «la familia extensa ha desaparecido poco a poco».}
\item \textbf{Traduire} (en français, pour toi) : « le texte dit : la famille élargie a peu à peu disparu ».
\item \textbf{Analyser} (en espagnol) : \esp{Esta frase muestra que la urbanización ha transformado las costumbres.} « Cette phrase montre que l'urbanisation a transformé les coutumes. »
\end{itemize}
\end{aretenir}

\courssub{Les trois axes appliqués à notre leçon}

Reprenons le texte support de cette leçon (\emph{Un domingo en familia}), qui illustre le thème « La familia cambia ». Voici comment construire chaque axe.

\textbf{Axe 1 : les mutations de la famille.}\par
Le texte oppose deux modèles : la \esp{familia extensa} (famille élargie : grands-parents, oncles, tantes, cousins vivant ensemble ou près les uns des autres) et la \esp{familia nuclear} (famille nucléaire : seulement les parents et les enfants). On observe aussi l'apparition de la \esp{familia monoparental} (famille monoparentale) et de la \esp{familia recompuesta} (famille recomposée). Le verbe clé est \esp{cambiar} : \esp{la familia ha cambiado mucho}.

\textbf{Axe 2 : les causes des conflits intergénérationnels.}\par
C'est ici que tu réutilises la grammaire de la cause : \esp{Los conflictos nacen \textbf{porque} los jóvenes y los padres no comparten los mismos valores} « Les conflits naissent parce que les jeunes et les parents ne partagent pas les mêmes valeurs ». Les trois causes principales à mobiliser : les TIC (\esp{las TIC} / \esp{las nuevas tecnologías} : \esp{los jóvenes pasan horas en el teléfono}), les valeurs (\esp{la autoridad y el respeto}) et l'urbanisation (\esp{los hijos dejan el pueblo para vivir en la ciudad}).

\textbf{Axe 3 : les solutions.}\par
C'est le domaine du conseil : \esp{Los padres \textbf{deberían} escuchar a sus hijos}, \esp{\textbf{hay que} practicar el diálogo}, \esp{\textbf{es importante} respetar a los mayores}. Les mots-clés : \esp{diálogo}, \esp{respeto}, \esp{escucha} (écoute), \esp{paciencia}.

\begin{center}
\begin{tabularx}{\linewidth}{|X|X|X|}
\hline
\rowcolor{popblueL} Partie du commentaire & Español (expressions utiles) & Français \\
\hline
Introduction & \esp{Este texto trata de... / El tema es...} & Ce texte traite de... / Le thème est... \\
\hline
Problématique & \esp{¿Por qué hay conflictos entre generaciones?} & Pourquoi y a-t-il des conflits entre générations ? \\
\hline
Citer & \esp{El autor dice: «...» / Según el texto...} & L'auteur dit : « ... » / Selon le texte... \\
\hline
Ajouter & \esp{además, también} & de plus, aussi \\
\hline
Opposer & \esp{sin embargo, pero, antes... hoy...} & cependant, mais, autrefois... aujourd'hui... \\
\hline
Cause & \esp{porque, como, a causa de} & parce que, comme, à cause de \\
\hline
Conséquence & \esp{por eso, entonces, así que} & c'est pourquoi, donc, si bien que \\
\hline
Conseil & \esp{deberías, hay que, es importante} & tu devrais, il faut, il est important \\
\hline
Conclusion & \esp{En conclusión... / En mi opinión...} & En conclusion... / À mon avis... \\
\hline
\end{tabularx}
\end{center}

\courssub{Texte support : Un domingo en familia}

Voici le texte original sur lequel porte le commentaire modèle. Il reprend tout le lexique de la leçon.

\begin{textebox}[Un domingo en familia]
En la casa de los Mbarga, en Yaoundé, los domingos son días de reunión. El abuelo Samuel, de setenta y cinco años, prepara el ndolé mientras sus nietos miran sus teléfonos. «Antes, los jóvenes escuchaban las historias de los mayores», suspira el abuelo. «Hoy, solo miran las pantallas.»

Su nieta Amina, de dieciséis años, no está de acuerdo: «Abuelo, con mi teléfono aprendo español y hablo con mis primos de Malabo. ¡El mundo ha cambiado!»

La madre de Amina, señora Mbarga, intenta calmar a todos. Ella piensa que los conflictos entre generaciones nacen porque nadie escucha a nadie. Los padres quieren imponer su autoridad, y los jóvenes quieren más libertad. Por eso, cada domingo, la familia organiza ahora un momento de diálogo: cada uno puede hablar de sus problemas, y los demás deben escuchar sin interrumpir.

Poco a poco, la tensión disminuye. El abuelo Samuel ha aceptado aprender a usar WhatsApp, y Amina ha prometido dejar su teléfono durante las comidas. «El respeto y la escucha son la base de la familia», concluye la señora Mbarga, sonriendo.
\end{textebox}

\textbf{Glossaire :} \esp{reunión} = réunion ; \esp{mientras} = pendant que ; \esp{suspirar} = soupirer ; \esp{pantalla} = écran ; \esp{no estar de acuerdo} = ne pas être d'accord ; \esp{calmar} = calmer ; \esp{imponer} = imposer ; \esp{sin interrumpir} = sans interrompre ; \esp{disminuir} = diminuer ; \esp{prometer} = promettre ; \esp{sonreír} = sourire.

\textbf{Questions de compréhension :}
\begin{enumerate}
\item ¿Qué hacen los nietos mientras el abuelo cocina? (Que font les petits-enfants pendant que le grand-père cuisine ?)
\item ¿Por qué hay conflictos en la familia, según la señora Mbarga?
\item ¿Qué solución organiza la familia cada domingo?
\end{enumerate}

\begin{corrige}[Questions de compréhension]
\begin{enumerate}
\item \esp{Los nietos miran sus teléfonos mientras el abuelo cocina.} « Les petits-enfants regardent leurs téléphones pendant que le grand-père cuisine. »
\item \esp{Según la señora Mbarga, los conflictos nacen porque nadie escucha a nadie.} « Selon Mme Mbarga, les conflits naissent parce que personne n'écoute personne. »
\item \esp{Cada domingo, la familia organiza un momento de diálogo donde cada uno puede hablar y los demás deben escuchar.} « Chaque dimanche, la famille organise un moment de dialogue où chacun peut parler et où les autres doivent écouter. »
\end{enumerate}
\end{corrige}

\courssub{Commentaire modèle}

Voici un commentaire court (10 à 15 lignes), rédigé en espagnol, qui réutilise le lexique et les connecteurs de la leçon. Observe comment chaque citation est suivie d'une analyse.

\begin{exemplebox}[Comentario modelo]
\esp{Este texto trata de la vida de una familia de Yaoundé y muestra que la familia ha cambiado mucho: antes era extensa, hoy es nuclear, y las costumbres son diferentes. El tema principal es el conflicto entre generaciones. El abuelo dice: «Antes, los jóvenes escuchaban las historias de los mayores»; esta frase muestra la nostalgia de los mayores. Sin embargo, Amina responde que el teléfono también sirve para aprender y comunicarse. Los conflictos nacen porque los jóvenes y los padres no comparten los mismos valores: los mayores defienden la autoridad y el respeto, y los jóvenes quieren más libertad. Por eso, el diálogo es necesario. La madre propone una solución: un momento de diálogo cada domingo. En mi opinión, esta idea es excelente, porque hay que escuchar para comprender. En conclusión, la familia cambia, pero el respeto y la escucha deben quedar.}
\end{exemplebox}

\textbf{Traduction :} « Ce texte traite de la vie d'une famille de Yaoundé et montre que la famille a beaucoup changé : autrefois elle était élargie, aujourd'hui elle est nucléaire, et les coutumes sont différentes. Le thème principal est le conflit entre générations. Le grand-père dit : "Autrefois, les jeunes écoutaient les histoires des aînés" ; cette phrase montre la nostalgie des aînés. Cependant, Amina répond que le téléphone sert aussi à apprendre et à communiquer. Les conflits naissent parce que les jeunes et les parents ne partagent pas les mêmes valeurs : les aînés défendent l'autorité et le respect, et les jeunes veulent plus de liberté. C'est pourquoi le dialogue est nécessaire. La mère propose une solution : un moment de dialogue chaque dimanche. À mon avis, cette idée est excellente, car il faut écouter pour comprendre. En conclusion, la famille change, mais le respect et l'écoute doivent rester. »

Remarque la mécanique : \textbf{citation} (« Antes, los jóvenes escuchaban... ») → \textbf{analyse} (\esp{esta frase muestra...}) → \textbf{connecteur} (\esp{sin embargo}, \esp{por eso}). C'est exactement la règle d'or de la méthode (la traduction se fait dans ton brouillon).

\begin{attention}[Les trois fautes qui coûtent cher au commentaire]
\begin{itemize}
\item \textbf{Recopier de longs passages sans analyse.} Une citation ne vaut rien seule : elle doit toujours être suivie d'une explication (\esp{esta frase muestra que...}, \esp{esto significa que...}).
\item \textbf{Traduire mot à mot dans la copie finale.} Le commentaire s'écrit \emph{en espagnol}. \esp{Antes era extensa, hoy es nuclear} ne s'analyse pas en français dans la copie. Attention au faux ami : \esp{extenso/a} signifie « vaste, élargi », pas « extensif » au sens agricole.
\item \textbf{Oublier les connecteurs.} Un commentaire sans \esp{porque}, \esp{sin embargo}, \esp{por eso} est une liste de phrases, pas un texte. Chaque axe doit être relié au suivant.
\end{itemize}
\end{attention}

\begin{exoresolu}[Transformer des phrases en mini-commentaire]
\textbf{Énoncé.} À partir de la phrase du texte \esp{«El respeto y la escucha son la base de la familia»}, rédige trois lignes de commentaire en suivant le schéma : citation + traduction + analyse avec un connecteur de conséquence.
\tcblower
\textbf{Solution détaillée.}
\begin{enumerate}
\item \textbf{Citation} : \esp{La madre concluye: «El respeto y la escucha son la base de la familia».}
\item \textbf{Traduction} (brouillon) : « la mère conclut : le respect et l'écoute sont la base de la famille ».
\item \textbf{Analyse + conséquence} : \esp{Esta frase resume la idea principal del texto: por eso, el diálogo es la mejor solución para los conflictos entre generaciones.} « Cette phrase résume l'idée principale du texte : c'est pourquoi le dialogue est la meilleure solution aux conflits entre générations. »
\end{enumerate}
On remarque l'accord : \esp{la base} (féminin singulier), \esp{la mejor solución} (féminin singulier), et le connecteur \esp{por eso} placé après la virgule.
\end{exoresolu}

\courssub{Production guidée : « Soluciones para el diálogo familiar »}

À ton tour ! Rédige un paragraphe de 6 à 8 lignes en espagnol en suivant ce plan détaillé, ligne par ligne.

\begin{methode}[Plan détaillé de ton paragraphe]
\begin{enumerate}
\item \textbf{Phrase d'accroche} (le problème) : \esp{En muchas familias, hay conflictos entre padres e hijos porque...} (complète avec une cause : los teléfonos, los valores, la autoridad).
\item \textbf{Première solution} (conseil aux parents) : \esp{Los padres deberían escuchar más a sus hijos.}
\item \textbf{Deuxième solution} (conseil aux jeunes) : \esp{Los jóvenes deberían respetar a los mayores y pasar menos tiempo en el teléfono.}
\item \textbf{Solution collective} (obligation générale) : \esp{Hay que organizar momentos de diálogo en familia.}
\item \textbf{Conséquence} : \esp{Por eso / Así, la familia será más unida y feliz.}
\item \textbf{Avis personnel} : \esp{En mi opinión, el diálogo y el respeto son la base de la familia.}
\end{enumerate}
\end{methode}

\begin{exercice}[À toi de jouer]
Rédige ton paragraphe \esp{Soluciones para el diálogo familiar} en suivant le plan ci-dessus. Contraintes : utilise au moins une fois \esp{porque}, une fois \esp{por eso}, une fois \esp{deberían} et une fois \esp{hay que}. Vérifie les accents (\esp{diálogo}, \esp{opinión}, \esp{además}) et les accords (\esp{la familia unida}, \esp{los padres}, \esp{las soluciones}).
\end{exercice}

\begin{experience}[Débat en clase : ¿El teléfono destruye la familia?]
Par groupes de quatre : deux élèves défendent les \esp{abuelos} (\esp{El teléfono destruye la comunicación familiar}), deux élèves défendent les \esp{jóvenes} (\esp{El teléfono sirve para aprender y comunicarse}). Chaque élève doit prononcer au moins deux phrases avec \esp{porque} ou \esp{por eso}, puis le groupe propose ensemble une solution avec \esp{deberíamos} ou \esp{hay que}. Le professeur note les connecteurs utilisés au tableau.
\end{experience}

\begin{savaistu}[Le commentaire au Bac]
Au baccalauréat camerounais, l'expression écrite porte très souvent sur des sujets de société comme la famille, et elle vaut plusieurs points de la note finale. Bonne nouvelle : tu n'as pas besoin d'un vocabulaire compliqué pour réussir. Maîtriser parfaitement trois outils — \esp{porque} (la cause), \esp{por eso} (la conséquence) et \esp{deberías / hay que} (le conseil) — suffit à construire un paragraphe cohérent, logique et bien noté. Les correcteurs récompensent la clarté et la structure (introduction, axes, conclusion) bien plus que les phrases longues et risquées. En Guinée équatoriale, voisin hispanophone du Cameroun, la \esp{familia extensa} reste très présente : c'est un excellent exemple de comparaison culturelle à citer dans ta conclusion pour montrer ton ouverture sur le monde hispanique.
\end{savaistu}

\begin{aretenir}[L'essentiel de la section]
\begin{itemize}
\item Le commentaire suit la structure en entonnoir : \textbf{introduction} (présentation + problématique) → \textbf{2 ou 3 axes} → \textbf{conclusion} avec avis personnel.
\item Chaque axe suit la règle d'or : \textbf{citer} en espagnol entre guillemets, \textbf{traduire} brièvement (brouillon), \textbf{analyser} en espagnol, puis \textbf{enchaîner} avec un connecteur.
\item Les trois axes de notre leçon : les \textbf{changements} de la famille, les \textbf{causes} des conflits, les \textbf{solutions} (diálogo, respeto, escucha).
\item Les connecteurs \esp{porque}, \esp{por eso}, \esp{deberías}, \esp{hay que} sont la colonne vertébrale de ton texte.
\end{itemize}
\end{aretenir}

\newpage

\coursec{Activité d'intégration}

\coursec{EA01 : La familia : evolución y conflictos intergeneracionales}

\courssub{Objectifs de l'activité}

À l'issue de cette activité d'intégration, tu seras capable de :

\begin{itemize}
\item mobiliser dans une situation de communication authentique le \cle{lexique de la famille} et des relations entre générations (\esp{abuelos, padres, hijos, brecha generacional, autoridad, respeto, conflicto}) ;
\item exprimer la \cle{cause} d'un conflit familial avec \esp{porque}, \esp{ya que}, \esp{como} ;
\item exprimer la \cle{conséquence} avec \esp{por eso}, \esp{de modo que}, \esp{entonces} ;
\item formuler des \cle{conseils} avec \esp{deberías + infinitif}, \esp{hay que + infinitif} et l'impératif ;
\item négocier à l'oral en espagnol et rédiger en groupe un texte court, cohérent et correct : une \esp{carta de diálogo familiar} (charte du dialogue familial) ;
\item développer des compétences de vie : écoute, respect de la parole de l'autre, recherche de compromis.
\end{itemize}

\courssub{Situation}

Tu es élève en classe de Première A au lycée de Nkol-Eton, à Yaoundé. Dans le quartier, les disputes entre parents et adolescents sont fréquentes : téléphone portable, réseaux sociaux, sorties, études, respect des aînés... L'association des parents d'élèves organise une \esp{mesa redonda familiar} (table ronde familiale) pour réconcilier les générations. Ton professeur d'espagnol a été invité à animer l'atelier des jeunes, et il a choisi ta classe pour préparer et présenter des solutions.

Voici le document qui a servi d'invitation à cette rencontre :

\begin{textebox}[Invitación de la Asociación de Padres de Alumnos]
\textbf{MESA REDONDA FAMILIAR : « HABLEMOS, ESCUCHÉMOS, ENTENDÁMONOS »}\par
\medskip
Queridas familias:\par
\medskip
Nuestra asociación organiza una mesa redonda sobre los conflictos entre padres e hijos. Hoy, muchos adolescentes pasan horas con el teléfono \textbf{móvil} y no hablan con sus padres. Muchos padres, por eso, se enfadan y castigan sin dialogar. La brecha generacional crece porque las costumbres cambian muy rápido.\par
\medskip
Invitamos a los alumnos de español a preparar una \emph{carta de diálogo familiar} con tres soluciones concretas. Las mejores cartas se leerán durante la reunión y se colgarán en el tablón de anuncios del barrio.\par
\medskip
¡El diálogo es el camino! ¡Contamos con vosotros!\par
\medskip
\textit{La presidenta de la asociación}
\end{textebox}

\textbf{Traduction du document :} « Chères familles : notre association organise une table ronde sur les conflits entre parents et enfants. Aujourd'hui, beaucoup d'adolescents passent des heures avec le téléphone portable et ne parlent pas avec leurs parents. Beaucoup de parents, c'est pourquoi, se fâchent et punissent sans dialoguer. Le fossé générationnel grandit parce que les coutumes changent très vite. Nous invitons les élèves d'espagnol à préparer une charte du dialogue familial avec trois solutions concrètes. Les meilleures chartes seront lues pendant la réunion et affichées au panneau d'affichage du quartier. Le dialogue est le chemin ! Nous comptons sur vous ! »

\courssub{Consignes}

Travail en groupes de quatre élèves. Chaque groupe incarne une famille en conflit : \esp{un abuelo} (le grand-père, défenseur des traditions), \esp{un padre} (le père, sévère mais fatigué), \esp{una madre} (la mère, médiatrice), \esp{un hijo adolescente} (un adolescent accro à son téléphone).

\begin{enumerate}
\item \textbf{Compréhension du document.} Lis l'invitation de l'association. Souligne dans le texte : deux mots du lexique de la famille, une expression de la cause et une expression de la conséquence. Recopie-les sur ton cahier.

\item \textbf{Préparation des rôles (10 minutes).} Chaque membre du groupe prépare trois phrases à la première personne pour défendre son point de vue. Chaque phrase doit contenir une cause introduite par \esp{porque} ou \esp{ya que}. Exemple pour l'adolescent : \esp{Uso mucho el móvil porque hablo con mis amigos y busco información para los deberes.}

\item \textbf{Jeu de rôle : la table ronde (10 minutes).} Chacun expose son point de vue. Les autres écoutent sans interrompre, puis réagissent avec un conseil : \esp{Deberías...}, \esp{Hay que...} ou un impératif (\esp{Escucha a tu padre}, \esp{Habla con calma}).

\item \textbf{Négociation.} Le groupe se met d'accord sur \textbf{trois solutions concrètes} acceptables par tous (par exemple : pas de téléphone pendant les repas, une heure de discussion familiale chaque soir, aide aux devoirs avant les réseaux sociaux).

\item \textbf{Production écrite.} Rédigez ensemble la \esp{carta de diálogo familiar} de \textbf{8 à 10 lignes}. Elle doit contenir obligatoirement :
\begin{itemize}
\item un titre et une formule d'ouverture ;
\item au moins deux expressions de la cause (\esp{porque}, \esp{ya que}, \esp{como}) ;
\item trois conseils avec \esp{deberías}, \esp{hay que} et un impératif ;
\item au moins une expression de la conséquence (\esp{por eso}, \esp{de modo que}) ;
\item une formule de clôture.
\end{itemize}

\item \textbf{Lecture et affichage.} Chaque groupe lit sa charte devant la classe. La classe vote pour les trois meilleures chartes, qui seront affichées au tableau du C.D.I.
\end{enumerate}

\begin{attention}[Critères de réussite]
Ta charte sera évaluée sur : la correction des structures de cause, conséquence et conseil ; l'usage du lexique de la famille ; la cohérence du texte ; la participation de chaque membre du groupe à l'oral. Attention aux accents (\esp{está}, \esp{adiós}, \esp{teléfono}, \esp{móvil}) et à la ponctuation espagnole (¡... !, ¿... ?).
\end{attention}

\courssub{Ressources à mobiliser}

\textbf{Lexique utile}

\begin{center}
\begin{tabularx}{\linewidth}{|X|X|}\hline
\rowcolor{popblueL} Español & Français \\ \hline
la brecha generacional & le fossé générationnel \\ \hline
la autoridad / el respeto & l'autorité / le respect \\ \hline
el conflicto / la discusión & le conflit / la dispute \\ \hline
castigar / prohibir & punir / interdire \\ \hline
escuchar / comprender / dialogar & écouter / comprendre / dialoguer \\ \hline
estar de acuerdo / llegar a un acuerdo & être d'accord / aboutir à un accord \\ \hline
pasar tiempo con la familia & passer du temps avec la famille \\ \hline
\end{tabularx}
\end{center}

\textbf{Structures grammaticales}

\begin{center}
\begin{tabularx}{\linewidth}{|X|X|X|}\hline
\rowcolor{popblueL} Fonction & Estructura & Ejemplo \\ \hline
Cause & porque, ya que, como + verbe conjugué & \esp{Me enfado porque no me escuchas.} \\ \hline
Conséquence & por eso, de modo que, entonces & \esp{No me escuchas, por eso me enfado.} \\ \hline
Conseil & deberías / deberíamos + infinitif & \esp{Deberías hablar con tu hijo.} \\ \hline
Obligation générale & hay que + infinitif & \esp{Hay que respetar a los abuelos.} \\ \hline
Conseil direct & impératif affirmatif & \esp{Escucha, habla, respeta.} \\ \hline
\end{tabularx}
\end{center}

\begin{aretenir}[Rappels essentiels]
\begin{itemize}
\item \esp{porque} répond à la question \esp{¿por qué?} ; \esp{ya que} et \esp{como} (en tête de phrase) sont des synonymes élégants de \esp{porque}.
\item \esp{por eso} reprend une idée déjà exprimée et annonce son résultat : \esp{El chico usa el móvil toda la noche; por eso está cansado.} « Le garçon utilise le portable toute la nuit ; c'est pourquoi il est fatigué. »
\item \esp{deberías} est le conditionnel de \esp{deber} : il exprime le conseil, jamais l'ordre. Ne confonds pas avec \esp{debes} (indicatif = obligation forte).
\item Impératif affirmatif de \esp{tú} : comme la 3\textsuperscript{e} personne du singulier du présent (\esp{escuchar → escucha}, \esp{hablar → habla}), avec des irréguliers : \esp{hacer → haz}, \esp{ir → ve}, \esp{ser → sé}, \esp{tener → ten}, \esp{poner → pon}, \esp{decir → di}, \esp{salir → sal}.
\end{itemize}
\end{aretenir}

\begin{attention}[Erreurs fréquentes des francophones]
\begin{itemize}
\item Ne traduis pas « parce que » par \esp{*para que} : \esp{para que} exprime le but (et demande le subjonctif), pas la cause. La cause = \esp{porque}.
\item Ne dis pas \esp{*por eso de que} : \esp{por eso} seul suffit.
\item Attention au faux ami : \esp{los padres} signifie « les parents » (père et mère), pas seulement « les pères ».
\item \textbf{Faux amis anglais :} \esp{actualmente} signifie « actuellement / de nos jours », pas « en fait » (=\esp{en realidad}) ; \esp{asistir a} signifie « assister à / être présent à », pas « aider » (=\esp{ayudar}).
\end{itemize}
\end{attention}

\courssub{Proposition de production et corrigé commenté}

\begin{exoresolu}[Carta de diálogo familiar : production modèle]
\textbf{Énoncé.} Rédige avec ton groupe une \esp{carta de diálogo familiar} de 8 à 10 lignes proposant trois solutions au conflit entre un adolescent accro au téléphone et sa famille, avec deux expressions de cause, trois conseils et une expression de conséquence.
\tcblower
\textbf{Production modèle :}\par
\medskip
\esp{\textbf{CARTA DE DIÁLOGO FAMILIAR}}\par
\medskip
\esp{Querida familia:}\par
\medskip
\esp{En nuestra casa hay conflictos porque el hijo usa el móvil todo el día y ya que los padres castigan sin dialogar, nadie se entiende. Por eso proponemos tres soluciones. Primero, hay que apagar los teléfonos durante las comidas para hablar juntos. Segundo, hijo, deberías terminar los deberes antes de usar las redes sociales. Tercero, padres, escuchad a vuestro hijo con calma y comprended su mundo. De modo que, si todos respetamos estas reglas, habrá más diálogo, más respeto y más amor en la familia.}\par
\medskip
\esp{¡El diálogo es el camino!}\par
\medskip
\esp{La familia Ngo Bell}\par
\medskip
\textbf{Traduction :} « Chère famille : dans notre maison il y a des conflits parce que le fils utilise le portable toute la journée et puisque les parents punissent sans dialoguer, personne ne se comprend. C'est pourquoi nous proposons trois solutions. Premièrement, il faut éteindre les téléphones pendant les repas pour parler ensemble. Deuxièmement, mon fils, tu devrais finir tes devoirs avant d'utiliser les réseaux sociaux. Troisièmement, parents, écoutez votre fils avec calme et comprenez son monde. De sorte que, si tous respectons ces règles, il y aura plus de dialogue, plus de respect et plus d'amour dans la famille. Le dialogue est le chemin ! La famille Ngo Bell. »\par
\medskip
\textbf{Commentaire des choix de langue :}
\begin{itemize}
\item \textbf{Causes} : \esp{porque} introduit la cause principale ; \esp{ya que} varie le style et évite la répétition — compétence visée : ne pas aligner trois \esp{porque} dans le même texte.
\item \textbf{Conséquence} : \esp{por eso} relie le constat au projet de solutions ; \esp{de modo que} annonce le résultat attendu à la fin. Les deux connecteurs structurent le texte.
\item \textbf{Conseils} : trois registres différents sont mobilisés — \esp{hay que + infinitif} (règle générale valable pour tous), \esp{deberías + infinitif} (conseil personnel adressé au fils), et deux impératifs pluriels \esp{escuchad, comprended} (vosotros, car on s'adresse aux deux parents).
\item \textbf{Lexique} : \esp{conflictos, castigar, dialogar, respeto, diálogo} reprennent le vocabulaire de la leçon ; \esp{apagar los teléfonos} et \esp{las redes sociales} ancrent le texte dans la vie quotidienne.
\item \textbf{Orthographe} : accents respectés (\esp{móvil, deberías, diálogo, habrá}) et points d'exclamation encadrants \esp{¡... !}.
\end{itemize}
\end{exoresolu}

\courssub{Bilan de l'activité}

\begin{aretenir}[Ce que tu as appris à faire]
\begin{itemize}
\item Tu as lu et compris un document authentique (une invitation associative) en repérant le lexique de la famille et les connecteurs logiques.
\item Tu as pris la parole à la première personne pour défendre un point de vue, en justifiant tes arguments par \esp{porque} et \esp{ya que}.
\item Tu as écouté les autres, réagi par des conseils (\esp{deberías}, \esp{hay que}, impératif) et négocié un compromis : compétence de vie essentielle.
\item Tu as rédigé en groupe un texte cohérent de 8 à 10 lignes, structuré par les connecteurs de cause et de conséquence, avec une ouverture, trois solutions numérotées (\esp{primero, segundo, tercero}) et une clôture.
\end{itemize}
\end{aretenir}

\textbf{Auto-évaluation.} Coche ce que tu sais faire désormais :

\begin{itemize}
\item[$\square$] Je peux nommer les membres de la famille et parler de la \esp{brecha generacional}.
\item[$\square$] Je peux expliquer la cause d'un conflit avec \esp{porque} et \esp{ya que} sans les confondre avec \esp{para que}.
\item[$\square$] Je peux exprimer une conséquence avec \esp{por eso} et \esp{de modo que}.
\item[$\square$] Je peux donner un conseil avec \esp{deberías}, \esp{hay que} et l'impératif.
\item[$\square$] Je peux participer à une discussion familiale en espagnol et rédiger une courte charte de dialogue.
\end{itemize}

\begin{savaistu}[Pour aller plus loin]
En Guinée équatoriale, seul pays hispanophone d'Afrique, la famille élargie joue un rôle central : grands-parents, oncles et tantes participent à l'éducation des enfants, comme dans beaucoup de familles camerounaises. Le dialogue entre générations y est une valeur fondamentale : on dit souvent que \esp{la palabra del anciano es una lámpara} — « la parole de l'ancien est une lampe ». Retiens ce proverbe : il pourra enrichir ta conclusion lors du prochain commentaire de texte sur la famille.
\end{savaistu}

\newpage

\coursec{Méthodes et exercices résolus}

\coursec{Leçon EA01 : La familia : evolución y conflictos intergeneracionales}

\courssub{Texte support : « La familia cambia »}

\begin{textebox}[La familia cambia]
En Yaundé y Duala, como en muchas ciudades de África y del mundo hispano, la estructura familiar ha cambiado mucho en veinte años. Antes, la familia extensa —abuelos, tíos, primos y padres viviendo bajo el mismo techo— era la norma. Hoy, la familia nuclear (padre, madre e hijos) predomina en los barrios urbanos. También aumentan las familias monoparentales, donde un solo progenitor cría a los hijos, y las familias reconstituidas, formadas tras una separación. La migración del campo a la ciudad y el coste de la vida obligan a los padres a trabajar largas jornadas. Los abuelos, antes pilares de la educación, se sienten a veces marginados. Esta transformación genera a veces una brecha generacional: los abuelos reclaman respeto y autoridad, mientras los jóvenes, conectados a Internet y a las redes sociales, exigen más autonomía. Los teléfonos móviles cambian la comunicación: los padres trabajan lejos de casa y los hijos estudian en línea. Sin embargo, los expertos coinciden: el diálogo sincero y la escucha mutua son el único puente para superar los conflictos y mantener la unión familiar.
\end{textebox}

\textbf{Traduction et glossaire.}\par
« À Yaoundé et Douala, comme dans de nombreuses villes d'Afrique et du monde hispanophone, la structure familiale a beaucoup changé en vingt ans. Autrefois, la famille élargie — grands-parents, oncles, cousins et parents vivant sous le même toit — était la norme. Aujourd'hui, la famille nucléaire (père, mère et enfants) prédomine dans les quartiers urbains. Augmentent aussi les familles monoparentales, où un seul parent élève les enfants, et les familles recomposées, formées après une séparation. La migration du village vers la ville et le coût de la vie obligent les parents à travailler de longues journées. Les grands-parents, autrefois piliers de l'éducation, se sentent parfois mis à l'écart. Cette transformation crée parfois un fossé générationnel : les grands-parents réclament respect et autorité, tandis que les jeunes, connectés à Internet et aux réseaux sociaux, exigent plus d'autonomie. Les téléphones portables changent la communication : les parents travaillent loin de la maison et les enfants étudient en ligne. Cependant, les experts s'accordent : le dialogue sincère et l'écoute mutuelle sont le seul pont pour surmonter les conflits et maintenir l'union familiale. »

\begin{center}
\begin{tabularx}{\linewidth}{|X|X|X|}
\hline
\rowcolor{popblueL} \textbf{Mot / Locución} & \textbf{Traduction} & \textbf{Note} \\
\hline
\esp{la estructura familiar} & la structure familiale & \emph{estructura} = structure, organisation \\
\esp{bajo el mismo techo} & sous le même toit & locution figée \\
\esp{la norma} & la norme, la règle & \\
\esp{predomina} & prédomine (verbe \esp{predominar}) & présent indicatif, 3e pers. sg. \\
\esp{un solo progenitor} & un seul parent (père ou mère) & \esp{progenitor} = parent biologique \\
\esp{criar} & élever (des enfants) & \esp{criar a los hijos} \\
\esp{tras una separación} & après une séparation & \esp{tras} = après (préposition) \\
\esp{la migración} & la migration & \esp{migrar} = migrer \\
\esp{el coste de la vida} & le coût de la vie & \\
\esp{jornadas} & journées (de travail) & \esp{jornada laboral} \\
\esp{pilares} & piliers & figé : \esp{ser pilar de...} \\
\esp{marginados} & mis à l'écart, marginalisés & participe passé adjectivé \\
\esp{la brecha generacional} & le fossé générationnel & \emph{brecha} = brèche, fossé \\
\esp{reclaman} & réclament (verbe \esp{reclamar}) & \\
\esp{la autonomía} & l'autonomie & \\
\esp{la escucha mutua} & l'écoute mutuelle & \esp{escuchar} = écouter \\
\esp{el puente} & le pont & image : \esp{puente para superar} \\
\hline
\end{tabularx}
\end{center}

\courssub{Lexique : la famille et les générations}

\begin{definition}[Membres de la famille et types de familles]
\begin{center}
\begin{tabularx}{\linewidth}{|X|X|X|}
\hline
\rowcolor{popblueL} \textbf{Français} & \textbf{Español} & \textbf{Remarque / Piège} \\
\hline
Le grand-père / La grand-mère & \esp{el abuelo / la abuela} & Collectif : \esp{los abuelos} (masc. plur.) \\
Le père / La mère & \esp{el padre / la madre} & \cle{Collectif : \esp{los padres} = les parents (père ET mère), pas « les pères »} \\
Le fils / La fille & \esp{el hijo / la hija} & Collectif : \esp{los hijos} = les enfants \\
Le frère / La sœur & \esp{el hermano / la hermana} & Collectif : \esp{los hermanos} \\
L'oncle / La tante & \esp{el tío / la tía} & Collectif : \esp{los tíos} \\
Le cousin / La cousine & \esp{el primo / la prima} & Collectif : \esp{los primos} \\
Le neveu / La nièce & \esp{el sobrino / la sobrina} & \\
\hline
\multicolumn{3}{|c|}{\textbf{Types de familles}} \\
\hline
La famille élargie & \esp{la familia extensa} & Syn. : \esp{la familia grande} \\
La famille nucléaire & \esp{la familia nuclear} & Père, mère, enfants \\
La famille monoparentale & \esp{la familia monoparental} & Un seul \esp{progenitor} \\
La famille recomposée & \esp{la familia reconstituida} & Après divorce / séparation \\
\hline
\multicolumn{3}{|c|}{\textbf{Relations, conflits, valeurs}} \\
\hline
Le fossé générationnel & \esp{la brecha generacional} & \esp{brecha} = brèche, écart \\
L'autorité & \esp{la autoridad} & \esp{ejercer la autoridad} \\
Le respect & \esp{el respeto} & \esp{mostrar respeto} \\
Le conflit / La dispute & \esp{el conflicto / la discusión} & \esp{resolver un conflicto} \\
Le dialogue & \esp{el diálogo} & \textbf{Accent sur le \emph{o}} \\
La communication & \esp{la comunicación} & \esp{falta de comunicación} \\
L'autonomie & \esp{la autonomía} & \esp{buscar autonomía} \\
Les valeurs & \esp{los valores} & \esp{valores tradicionales / modernos} \\
Les traditions & \esp{las tradiciones} & \\
\hline
\end{tabularx}
\end{center}
\end{definition}

\begin{aretenir}[Verbes utiles pour la leçon]
\begin{center}
\begin{tabular}{|l|l|l|}
\hline
\rowcolor{popblueL} \textbf{Infinitif} & \textbf{Sens} & \textbf{Présent (yo / tú / él) / Part. passé} \\
\hline
\esp{vivir} & habiter, vivre & \esp{vivo, vives, vive} / \esp{vivido} \\
\esp{convivir} & cohabiter, vivre ensemble & \esp{convivimos, convivís, conviven} \\
\esp{cambiar} & changer & \esp{cambio, cambias, cambia} / \esp{cambiado} \\
\esp{crear} & créer & \esp{creo, creas, crea} / \esp{creado} \\
\esp{generar} & générer, créer & \esp{genero, generas, genera} / \esp{generado} \\
\esp{obligar} & obliger & \esp{obligo, obligas, obliga} / \esp{obligado} \\
\esp{sentir(se)} & se sentir & \esp{siento, sientes, siente} / \esp{sentido} \\
\esp{reclamar} & réclamer & \esp{reclamo, reclamas, reclama} \\
\esp{exigir} & exiger & \esp{exijo, exiges, exige} (g $\rightarrow$ j) \\
\esp{superar} & surmonter & \esp{supero, superas, supera} \\
\esp{mantener} & maintenir & \esp{mantengo, mantienes, mantiene} (g $\rightarrow$ ng) \\
\esp{dialogar} & dialoguer & \esp{dialogo, dialogas, dialoga} \\
\esp{escuchar} & écouter & \esp{escucho, escuchas, escucha} \\
\esp{entender} & comprendre & \esp{entiendo, entiendes, entiende} (e $\rightarrow$ ie) \\
\hline
\end{tabular}
\end{center}
\end{aretenir}

\courssub{Grammaire 1 : Exprimer la cause}

\begin{propriete}[Les connecteurs de cause]
En espagnol, la cause s'exprime principalement de quatre façons. Le choix dépend du registre (courant / soutenu) et de la nature du mot qui suit (verbe ou nom).

\begin{enumerate}
\item \textbf{\esp{Porque} + indicatif} (cause réelle, explication nouvelle). \emph{Le plus courant.}
\esp{Los jóvenes se van \textbf{porque} no hay trabajo en el pueblo.}
« Les jeunes partent \textbf{parce qu'}il n'y a pas de travail au village. »

\item \textbf{\esp{Como} + indicatif} (cause connue, évidente, souvent en tête de phrase).
\esp{\textbf{Como} llueve, no salimos.} « \textbf{Comme} il pleut, on ne sort pas. »

\item \textbf{\esp{Por} + nom / pronom / infinitif} (cause, motif, raison).
\esp{Lo hice \textbf{por} ti.} « Je l'ai fait \textbf{pour} toi / \textbf{à cause de} toi. »
\esp{El conflicto es \textbf{por} la herencia.} « Le conflit est \textbf{à cause de} l'héritage. »

\item \textbf{\esp{Debido a} / \esp{A causa de} + nom} (registre soutenu, écrit, journalistique).
\esp{\textbf{Debido a} la crisis, muchas familias emigran.} « \textbf{En raison de} la crise, beaucoup de familles émigrent. »
\esp{\textbf{A causa de} la lluvia, se canceló el partido.} « \textbf{À cause de} la pluie, le match a été annulé. »
\end{enumerate}

\cle{Attention :} \esp{Por qué} (deux mots, accent) = « pourquoi ? » (question). \esp{Porque} (un mot, sans accent) = « parce que ». \esp{Porqué} (un mot, accent, nom) = « le pourquoi, la raison » (\esp{el porqué}). \esp{Por que} (deux mots, sans accent) = \esp{por} + \esp{que} (relatif ou conjonction).
\end{propriete}

\begin{exemplebox}[Exemples contrastés cause / but]
\begin{itemize}
\item \esp{Estudio \textbf{porque} quiero aprender.} (Cause : raison présente)
\item \esp{Estudio \textbf{para} aprender.} (But : finalité future, \esp{para} + infinitif)
\item \esp{No vine \textbf{por} la lluvia.} (Cause : nom)
\item \esp{No vine \textbf{debido a} la lluvia.} (Cause soutenue)
\end{itemize}
\end{exemplebox}

\courssub{Grammaire 2 : Exprimer la conséquence}

\begin{propriete}[Les connecteurs de conséquence]
Ils relient une cause (souvent énoncée juste avant) à son résultat logique.

\begin{center}
\begin{tabularx}{\linewidth}{|X|X|X|}
\hline
\rowcolor{popblueL} \textbf{Connecteur} & \textbf{Niveau / Emploi} & \textbf{Exemple traduit} \\
\hline
\esp{Por eso} & Courant, neutre. « C'est pourquoi, donc. » & \esp{No hay trabajo, \textbf{por eso} los jóvenes se van.} \\
\esp{Entonces} & Courant, oral. « Alors, dans ce cas. » & \esp{¿\textbf{Entonces} qué hacemos?} \\
\esp{Así que} & Courant, familier. « Du coup, donc. » & \esp{Es tarde, \textbf{así que} me voy.} \\
\esp{Por lo tanto} & Soutenu, écrit. « Par conséquent, donc. » & \esp{Hay crisis; \textbf{por lo tanto}, suben los precios.} \\
\esp{Por consiguiente} & Très soutenu. « Par conséquent. » & \esp{Estudió mucho; \textbf{por consiguiente}, aprobó.} \\
\hline
\end{tabularx}
\end{center}

\cle{Règle de ponctuation :} Ces connecteurs sont souvent précédés d'un point-virgule (;) ou d'un point (.) et suivis d'une virgule (,) s'ils sont en tête de proposition. \esp{No hay dinero; por eso, no compramos.}
\end{propriete}

\courssub{Grammaire 3 : Exprimer le conseil et l'obligation}

\begin{propriete}[Conseil, obligation, nécessité : structures clés]
Au niveau A2+/B1, on distingue l'obligation générale (impersonnelle) et le conseil personnel.

\begin{enumerate}
\item \textbf{\esp{Hay que} + infinitif} (Impersonnel, présent) = « Il faut / Il est nécessaire de ». \emph{Invariable.}
\esp{\textbf{Hay que} dialogar más.} « \textbf{Il faut} dialoguer davantage. »

\item \textbf{\esp{Deber} + infinitif} (Personnel) = « Devoir (obligation morale, conseil) ».
\begin{itemize}
\item Présent : \esp{Debes escuchar a tus padres.} « Tu dois / Tu devrais écouter tes parents. »
\item \cle{Conditionnel (plus poli, conseil) :} \esp{Deberías escuchar a tus padres.} « \textbf{Tu devrais} écouter tes parents. » \emph{Forme recommandée pour le conseil.}
\end{itemize}

\item \textbf{\esp{Tener que} + infinitif} (Personnel) = « Avoir à / Devoir (obligation concrète, contrainte) ».
\esp{Tengo que irme ya.} « Je dois partir maintenant. »

\item \textbf{\esp{Es necesario / Es importante / Es fundamental} + infinitif} (Impersonnel, évaluation).
\esp{\textbf{Es importante} respetar a los mayores.} « \textbf{Il est important} de respecter les aînés. »

\item \textbf{\esp{Es necesario / Es importante / Te aconsejo / Te recomiendo} + \textbf{que} + \textbf{subjonctif}} (Personnel, sujet différent).
\esp{\textbf{Es necesario que} los padres \textbf{escuchen} a sus hijos.} (Subj. \esp{escuchen}).
\esp{\textbf{Te aconsejo que} \textbf{hables} con tu madre.} (Subj. \esp{hables}).
\end{enumerate}

\cle{Point de grammaire (B1) :} Après les expressions d'opinion, de conseil, de volonté, d'émotion (\esp{es importante que, quiero que, te aconsejo que, me gusta que}), le verbe suivant est au \textbf{subjonctif présent} si les deux sujets sont différents. Si le sujet est le même, on utilise l'infinitif (\esp{Es importante escuchar} / \esp{Quiero escuchar}).
\end{propriete}

\begin{exemplebox}[Résumé : Cause $\rightarrow$ Conséquence $\rightarrow$ Conseil]
\esp{Los abuelos no entienden Internet \textbf{porque} tienen otros valores. \textbf{Por eso}, hay conflictos. \textbf{Hay que} hablar. \textbf{Deberían} escucharse mutuamente.}
« Les grands-parents ne comprennent pas Internet \textbf{parce qu'}ils ont d'autres valeurs. \textbf{C'est pourquoi} il y a des conflits. \textbf{Il faut} parler. \textbf{Ils devraient} s'écouter mutuellement. »
\end{exemplebox}

\courssub{Méthode du commentaire et commentaire modèle}

\courssub{Méthode 1 : Lire et comprendre un texte sur la famille}

\begin{methode}[Comprendre un texte sur la famille : la démarche en 5 étapes]
\begin{enumerate}
\item \textbf{Lire le titre et la source.} Le titre (\esp{La familia cambia}) annonce le thème : repère tout de suite les mots-clés possibles (\esp{familia, padres, hijos, abuelos, conflicto, generación}).
\item \textbf{Première lecture globale.} Lis le texte entier sans dictionnaire. Souligne les mots transparents (proches du français : \esp{nuclear, autoridad, urbanización, comunicación}).
\item \textbf{Repérer la structure.} Un texte sur la famille suit souvent le plan : \emph{constat} (la famille a changé) $\rightarrow$ \emph{causes} (\esp{porque, debido a, a causa de}) $\rightarrow$ \emph{conséquences} (\esp{por eso, entonces, así que}) $\rightarrow$ \emph{solutions} (\esp{hay que, deberíamos}).
\item \textbf{Construire un mini-glossaire.} Pour chaque mot inconnu, devine le sens grâce au contexte avant de vérifier. Exemple : \esp{la brecha generacional} = le fossé entre les générations.
\item \textbf{Répondre aux questions en phrases complètes}, en reprenant les mots de la question, sans jamais recopier tout le texte. Commence par \esp{Según el texto...}, \esp{El autor dice que...}, \esp{Porque...}.
\end{enumerate}
\textbf{Réflexes à retenir :} une réponse = une phrase complète avec sujet + verbe conjugué ; justifier par une courte citation entre guillemets ; ne jamais répondre en français si la question est en espagnol.
\end{methode}

\begin{exoresolu}[Compréhension de texte : questions type Bac]
\textbf{Énoncé.} Lis le texte suivant, puis réponds en espagnol par des phrases complètes.

\esp{En muchas ciudades africanas, como Yaundé o Duala, la familia tradicional está cambiando. Antes, los abuelos, los padres y los hijos vivían juntos en la misma casa: era la familia extensa. Hoy, muchos jóvenes dejan el pueblo para estudiar o trabajar en la ciudad, y la familia nuclear —padre, madre e hijos— es cada vez más frecuente. Además, hay familias monoparentales, con un solo padre o una sola madre, y familias reconstituidas después de un divorcio. Esta evolución crea a veces conflictos entre las generaciones: los abuelos piensan que los jóvenes no respetan las tradiciones, y los jóvenes creen que sus padres no comprenden el mundo moderno de Internet y los teléfonos móviles. Sin embargo, el diálogo sigue siendo la mejor solución para unir a la familia.}

\begin{enumerate}
\item \esp{¿Qué tipos de familia menciona el texto?}
\item \esp{¿Por qué los jóvenes dejan el pueblo?}
\item \esp{¿Cuál es la causa de los conflictos entre las generaciones?}
\item \esp{¿Qué solución propone el autor?}
\end{enumerate}
\tcblower
\textbf{Raisonnement.} Il s'agit de questions de compréhension directe : les réponses se trouvent dans le texte, mais il faut les reformuler en phrases complètes, sujet + verbe conjugué.

\textbf{Question 1.} Le texte cite quatre types : \esp{familia extensa}, \esp{familia nuclear}, \esp{familia monoparental} et \esp{familia reconstituida}. Réponse rédigée :

\esp{El texto menciona cuatro tipos de familia: la familia extensa, la familia nuclear, la familia monoparental y la familia reconstituida.} « Le texte mentionne quatre types de famille : la famille élargie, la famille nucléaire, la famille monoparentale et la famille recomposée. »

\textbf{Question 2.} Le texte dit : \esp{dejan el pueblo para estudiar o trabajar}. Le \esp{para} + infinitif exprime le but. Réponse :

\esp{Los jóvenes dejan el pueblo para estudiar o trabajar en la ciudad.} « Les jeunes quittent le village pour étudier ou travailler en ville. »

\textbf{Question 3.} Attention : il faut dégager la cause générale (l'évolution de la famille et le fossé des valeurs), pas seulement un détail. Réponse :

\esp{La causa de los conflictos es que los abuelos y los jóvenes no tienen los mismos valores: los abuelos quieren respetar las tradiciones y los jóvenes viven en el mundo moderno de Internet.} « La cause des conflits est que les grands-parents et les jeunes n'ont pas les mêmes valeurs. »

\textbf{Question 4.} Dernière phrase du texte. Réponse :

\esp{Según el autor, el diálogo es la mejor solución para unir a la familia.} « Selon l'auteur, le dialogue est la meilleure solution pour unir la famille. »

\textbf{Conclusion méthodologique.} Chaque réponse commence par un verbe conjugué, reprend les mots de la question et reste courte. Aucune phrase n'est recopiée mot à mot sur plusieurs lignes : c'est ce que le correcteur attend.
\end{exoresolu}

\courssub{Méthode 2 : Identifier les causes des conflits et les solutions}

\begin{methode}[Repérer et formuler causes et solutions dans un texte]
\begin{enumerate}
\item \textbf{Repère les connecteurs de cause} : \esp{porque} (parce que), \esp{como} (comme, en début de phrase), \esp{debido a} + nom, \esp{a causa de} + nom, \esp{por} + nom.
\item \textbf{Repère les connecteurs de conséquence} : \esp{por eso} (c'est pourquoi), \esp{entonces}, \esp{así que}, \esp{por lo tanto} (donc, soutenu).
\item \textbf{Repère les formules de conseil/solution} : \esp{hay que} + infinitif (il faut), \esp{deberías} + infinitif (tu devrais), \esp{es importante} + infinitif, \esp{te aconsejo que} + subjonctif.
\item \textbf{Classe tes idées dans un tableau à deux colonnes} : \emph{causas del conflicto} / \emph{soluciones}. Cela t'aide à structurer ta réponse écrite.
\item \textbf{Pour rédiger}, enchaîne : constat $\rightarrow$ cause (\esp{porque}) $\rightarrow$ conséquence (\esp{por eso}) $\rightarrow$ conseil (\esp{hay que}).
\end{enumerate}
\textbf{Structures à retenir :} \esp{El conflicto existe porque...} ; \esp{Por eso, los padres deben...} ; \esp{Hay que escuchar a los jóvenes.}
\end{methode}

\begin{exoresolu}[Causes et solutions : analyse guidée]
\textbf{Énoncé.} Voici un extrait de dialogue entre une mère et sa fille à Yaundé :

\esp{— Hija, pasas demasiado tiempo con el teléfono móvil. Por eso no terminas tus deberes.\\
— Mamá, uso el móvil porque busco información para la escuela.\\
— Entonces, deberías explicarme lo que haces. Hay que hablar más en esta familia.}

\begin{enumerate}
\item Releve dans le dialogue : un connecteur de conséquence, un connecteur de cause, deux expressions du conseil.
\item Rédige en espagnol deux phrases : l'une qui résume le conflit, l'autre qui propose une solution.
\end{enumerate}
\tcblower
\textbf{Question 1.} Cherchons les connecteurs appris :
\begin{itemize}
\item Conséquence : \esp{Por eso} (« c'est pourquoi »), dans \esp{Por eso no terminas tus deberes}.
\item Cause : \esp{porque}, dans \esp{uso el móvil porque busco información}.
\item Conseil : \esp{deberías} + infinitif (\esp{deberías explicarme}, « tu devrais m'expliquer ») et \esp{Hay que} + infinitif (\esp{Hay que hablar más}, « il faut parler davantage »).
\end{itemize}

\textbf{Question 2.} Rédaction. Phrase 1 (résumé du conflit, avec cause) :

\esp{La madre está enfadada porque su hija pasa mucho tiempo con el móvil y no termina sus deberes.} « La mère est en colère parce que sa fille passe beaucoup de temps avec son portable et ne finit pas ses devoirs. »

Phrase 2 (solution, avec \esp{hay que} et \esp{deberían}) :

\esp{Para resolver el conflicto, hay que dialogar: los padres y los hijos deberían escucharse y explicarse con respeto.} « Pour résoudre le conflit, il faut dialoguer : les parents et les enfants devraient s'écouter et s'expliquer avec respect. »

\textbf{Conclusion.} La bonne réponse combine un connecteur correct (\esp{porque}, \esp{para} + infinitif) et une formule de conseil (\esp{hay que}, \esp{deberían} + infinitif). Remarque l'accord : \esp{está enfadada} (féminin, car c'est la mère).
\end{exoresolu}

\courssub{Méthode 3 : Traduire du français vers l'espagnol (thème) : cause, conséquence, conseil}

\begin{methode}[Réussir le thème sur le thème de la famille]
\begin{enumerate}
\item \textbf{Analyse la phrase française} : repère le verbe principal, son temps, et le connecteur logique (parce que, donc, il faut).
\item \textbf{Choisis le bon connecteur espagnol} : « parce que » = \esp{porque} ; « donc / c'est pourquoi » = \esp{por eso} ; « il faut » = \esp{hay que} + infinitif ; « tu devrais » = \esp{deberías} + infinitif.
\item \textbf{Conjugue correctement} : présent pour les vérités générales (\esp{los jóvenes usan}), imparfait pour le passé descriptif (\esp{antes vivían juntos}).
\item \textbf{Méfie-toi des calques} : « passer du temps » = \esp{pasar tiempo} (et non \esp{pasar el tiempo} partout) ; « respecter » = \esp{respetar} ; « une famille nombreuse » = \esp{una familia numerosa}.
\item \textbf{Relis} : accents (\esp{está, familia, móvil, diálogo}), accords (\esp{las familias monoparentales}), ponctuation espagnole (¿...? ¡...!).
\end{enumerate}
\end{methode}

\begin{exoresolu}[Thème : phrases à traduire]
\textbf{Énoncé.} Traduis en espagnol :
\begin{enumerate}
\item « Aujourd'hui, beaucoup de familles vivent en ville parce que les jeunes y trouvent du travail. »
\item « Les grands-parents ne comprennent pas Internet, c'est pourquoi il y a des conflits. »
\item « Il faut écouter les jeunes et tu devrais parler avec tes parents. »
\end{enumerate}
\tcblower
\textbf{Phrase 1.} Analyse : présent de vérité générale ; « beaucoup de familles » = \esp{muchas familias} ; « parce que » = \esp{porque} ; « y trouvent » : le pronom « y » ne se traduit pas ici, on dit simplement \esp{encuentran trabajo en la ciudad}.

\esp{Hoy, muchas familias viven en la ciudad porque los jóvenes encuentran trabajo allí.} (ou \esp{...porque los jóvenes encuentran trabajo en ella.})

\textbf{Phrase 2.} « C'est pourquoi » = \esp{por eso} (jamais \esp{por qué}, qui signifie « pourquoi ? »). « Il y a » = \esp{hay} (verbe \esp{haber}, invariable).

\esp{Los abuelos no comprenden Internet, por eso hay conflictos.}

\textbf{Phrase 3.} « Il faut » = \esp{hay que} + infinitif (\esp{escuchar}) ; « tu devrais » = conditionnel de \esp{deber} : \esp{deberías} + infinitif (\esp{hablar}) ; « avec tes parents » = \esp{con tus padres}.

\esp{Hay que escuchar a los jóvenes y deberías hablar con tus padres.}

\textbf{Justifications grammaticales.} (a) \esp{escuchar a los jóvenes} : la préposition \esp{a} est obligatoire devant un complément de personne (\emph{a} personnelle) — erreur classique des francophones. (b) \esp{deberías} porte un accent sur le \emph{i} (conditionnel). (c) \esp{jóvenes} prend un accent au pluriel pour conserver l'accentuation du singulier \esp{joven}.

\textbf{Conclusion.} Une bonne version de thème = connecteur juste + verbe au bon temps + accents corrects. Relis toujours chaque phrase en te demandant : « où est l'accent ? où est l'accord ? »
\end{exoresolu}

\courssub{Méthode 4 : Traduire de l'espagnol vers le français (version)}

\begin{methode}[Réussir la version : de l'espagnol au français]
\begin{enumerate}
\item \textbf{Repère le sujet et le verbe} de chaque proposition avant de traduire ; l'espagnol omet souvent le sujet (\esp{Viven en Duala} = « Ils vivent à Douala »).
\item \textbf{Traduis les connecteurs avec précision} : \esp{porque} = « parce que » ; \esp{por eso} = « c'est pourquoi » ; \esp{sin embargo} = « cependant » ; \esp{además} = « de plus ».
\item \textbf{Attention aux temps} : l'imparfait espagnol (\esp{vivían, pensaban}) = imparfait français ; le passé simple (\esp{cambió}) = passé composé ou passé simple français.
\item \textbf{Évite les faux amis} : \esp{actualmente} = « actuellement » (pas « en fait ») ; \esp{los parientes} = « les parents / la famille » au sens large (pas seulement père et mère : \esp{padres} = père et mère) ; \esp{asistir a} = « assister à ».
\item \textbf{Rédige en français naturel} : reformule si nécessaire, sans trahir le sens.
\end{enumerate}
\end{methode}

\begin{exoresolu}[Version : paragraphe à traduire]
\textbf{Énoncé.} Traduis en français :

\esp{Antes, las familias vivían juntas en el pueblo y los abuelos cuidaban a los niños. Hoy, muchos jóvenes se van a la ciudad porque allí hay más trabajo. Por eso, la familia extensa desaparece poco a poco. Sin embargo, los padres y los hijos deben hablar todos los días para mantener la unión familiar.}
\tcblower
\textbf{Analyse phrase par phrase.}

\emph{Phrase 1.} \esp{Antes} = « Autrefois » ; \esp{vivían} = imparfait (\esp{vivir}) ; \esp{cuidaban a los niños} = « gardaient les enfants » (\esp{cuidar a} = « s'occuper de », avec \emph{a} personnelle).

\emph{Phrase 2.} \esp{se van} = présent du pronominal \esp{irse} (« s'en aller, partir ») ; \esp{porque allí hay más trabajo} = « parce qu'il y a plus de travail là-bas ».

\emph{Phrase 3.} \esp{Por eso} = « C'est pourquoi » ; \esp{desaparece poco a poco} = « disparaît peu à peu » (verbe \esp{desaparecer}).

\emph{Phrase 4.} \esp{Sin embargo} = « Cependant » ; \esp{deben hablar} = « doivent parler » (\esp{deber} + infinitif) ; \esp{para mantener} = « pour maintenir » (\esp{para} + infinitif = but) ; \esp{la unión familiar} = « l'union familiale ».

\textbf{Traduction complète.}

« Autrefois, les familles vivaient ensemble au village et les grands-parents gardaient les enfants. Aujourd'hui, beaucoup de jeunes partent en ville parce qu'il y a plus de travail là-bas. C'est pourquoi la famille élargie disparaît peu à peu. Cependant, les parents et les enfants doivent parler tous les jours pour maintenir l'union familiale. »

\textbf{Conclusion.} Vérifie toujours : imparfait = imparfait, connecteurs rendus avec exactitude, et français fluide. Ici, \esp{padres} = « parents » (père et mère), pas « pères » : piège classique.
\end{exoresolu}

\courssub{Méthode 5 : Rédiger un court texte pour proposer le dialogue familial}

\begin{methode}[Rédiger 5 à 8 phrases : « Proposer le dialogue en famille »]
\begin{enumerate}
\item \textbf{Introduction (1 phrase)} : pose le problème. Modèle : \esp{En muchas familias hay conflictos entre padres e hijos.}
\item \textbf{Causes (2 phrases)} avec \esp{porque} ou \esp{a causa de} : \esp{Los conflictos existen porque los jóvenes y los padres no tienen las mismas ideas.}
\item \textbf{Conséquence (1 phrase)} avec \esp{por eso} : \esp{Por eso, es necesario buscar soluciones.}
\item \textbf{Solutions / conseils (2-3 phrases)} avec \esp{hay que}, \esp{deberíamos}, \esp{es importante} : \esp{Hay que dialogar con respeto. Los padres deberían escuchar a sus hijos.}
\item \textbf{Conclusion (1 phrase)} : \esp{Así, la familia será más unida y feliz.}
\end{enumerate}
\textbf{Réflexes :} phrases courtes ; un connecteur par phrase ; vocabulaire de la leçon (\esp{respeto, autoridad, brecha generacional, diálogo}) ; relecture des accords et des accents.
\end{methode}

\begin{exoresolu}[Production écrite guidée]
\textbf{Énoncé.} Rédige en espagnol un court texte (5 à 7 phrases) intitulé \esp{El diálogo en la familia}, dans lequel tu présentes le problème des conflits entre générations et tu proposes des solutions.
\tcblower
\textbf{Préparation.} Plan : constat $\rightarrow$ cause $\rightarrow$ conséquence $\rightarrow$ deux conseils $\rightarrow$ conclusion. Connecteurs prévus : \esp{porque}, \esp{por eso}, \esp{hay que}, \esp{deberían}, \esp{además}.

\textbf{Production corrigée.}

\esp{El diálogo en la familia}

\esp{Hoy, en muchas familias hay conflictos entre los padres y los hijos. Estos conflictos existen porque las generaciones no tienen los mismos valores: los jóvenes usan el móvil e Internet, y los abuelos prefieren las tradiciones. Por eso, la comunicación es a veces difícil. Para resolver este problema, hay que dialogar todos los días con respeto. Además, los padres deberían escuchar a sus hijos y los hijos deberían respetar la autoridad de sus padres. Así, la familia será más unida y todos serán más felices.}

« Aujourd'hui, dans beaucoup de familles, il y a des conflits entre les parents et les enfants. Ces conflits existent parce que les générations n'ont pas les mêmes valeurs : les jeunes utilisent le portable et Internet, et les grands-parents préfèrent les traditions. C'est pourquoi la communication est parfois difficile. Pour résoudre ce problème, il faut dialoguer tous les jours avec respect. De plus, les parents devraient écouter leurs enfants et les enfants devraient respecter l'autorité de leurs parents. Ainsi, la famille sera plus unie et tous seront plus heureux. »

\textbf{Justifications.} (a) \esp{hay} (il y a) est invariable ; (b) \esp{deberían escuchar} : conditionnel + infinitif pour le conseil ; (c) \esp{escuchar a sus hijos} : \emph{a} personnelle ; (d) \esp{será / serán} : futur simple pour la conclusion, correctement accentué ; (e) \esp{comunicación, diálogo, unida} : accents et accords vérifiés.

\textbf{Conclusion.} Ce texte obtient la note maximale car il respecte le plan, utilise trois connecteurs différents, et ne contient ni faute d'accord ni faute d'accent.
\end{exoresolu}

\begin{attention}[Les 5 erreurs qui coûtent des points au Bac]
\begin{enumerate}
\item \textbf{Oublier les accents.} \esp{está} (il est) $\neq$ \esp{esta} (cette) ; \esp{joven} $\rightarrow$ \esp{jóvenes} au pluriel ; \esp{deberías, diálogo, móvil, comunicación} prennent toujours leur accent. Une copie sans accents perd des points à chaque ligne.
\item \textbf{Confondre \esp{porque} et \esp{por qué}.} \esp{porque} (un seul mot, sans accent) = « parce que » ; \esp{¿Por qué?} (deux mots, accent) = « pourquoi ? ». Et \esp{por eso} = « c'est pourquoi ».
\item \textbf{Oublier l'\emph{a} personnelle.} On écrit \esp{escucho a mis padres}, \esp{los abuelos cuidan a los niños}. Devant un complément de personne, l'espagnol exige \esp{a} : le français ne le montre pas, d'où le calque fautif.
\item \textbf{Traduire « il faut » par \esp{es necesario que} + indicatif ou par un calque.} La forme simple et sûre est \esp{hay que} + infinitif : \esp{Hay que dialogar.} Après \esp{es importante que}, \esp{es necesario que} ou \esp{te aconsejo que}, il faut le \textbf{subjonctif} : \esp{Te aconsejo que hables con tu madre.}
\item \textbf{Faux amis et accords.} \esp{los padres} = « les parents » (père et mère), pas « les pères » ; \esp{los parientes} = « les membres de la famille » ; \esp{actualmente} = « actuellement ». Et accorde toujours : \esp{las familias monoparentales}, \esp{mi madre está enfadada}.
\end{enumerate}
\end{attention}

\newpage

\coursec{Exercices}

\courssub{Je vérifie mes connaissances}

\begin{exercice}[QCM : la famille et les générations]
Pour chaque question, entoure la bonne réponse.
\begin{enumerate}
\item Une famille formée uniquement par le père, la mère et leurs enfants est une famille :
\begin{enumerate}
\item extensa \item nuclear \item reconstituida \item monoparental
\end{enumerate}
\item \esp{La brecha generacional} désigne :
\begin{enumerate}
\item la distance entre la ville et le village \item le fossé entre les générations \item un conflit entre voisins \item une fête de famille
\end{enumerate}
\item Dans la phrase \esp{Los jóvenes discuten porque pasan mucho tiempo en el móvil}, \esp{porque} exprime :
\begin{enumerate}
\item la conséquence \item le conseil \item la cause \item le but
\end{enumerate}
\item Pour donner un conseil à un ami, on dit :
\begin{enumerate}
\item \esp{Deberías escuchar a tus padres.} \item \esp{Escuchas a tus padres.} \item \esp{Por eso escuchas a tus padres.} \item \esp{Escuchaste a tus padres.}
\end{enumerate}
\end{enumerate}
\end{exercice}

\begin{exercice}[Vrai ou faux ?]
Indique si chaque affirmation est vraie (V) ou fausse (F), puis justifie ta réponse par une phrase en espagnol.
\begin{enumerate}
\item \esp{Una familia monoparental está formada por un solo progenitor y sus hijos.}
\item \esp{La urbanización no influye en las relaciones familiares.}
\item \esp{Los abuelos transmiten tradiciones y valores a los nietos.}
\item \esp{El diálogo es una solución para los conflictos intergeneracionales.}
\end{enumerate}
\end{exercice}

\begin{exercice}[Vocabulaire : les membres de la famille]
Complète chaque phrase avec le mot qui convient, choisi dans la liste : \esp{abuelos} — \esp{padres} — \esp{hijos} — \esp{nietos} — \esp{hermanos}. Attention, un mot de la liste n'est pas utilisé.
\begin{enumerate}
\item Los \esp{\ldots\ldots} de mis \esp{\ldots\ldots} son mis bisabuelos.
\item Mis \esp{\ldots\ldots} son los hijos de mis abuelos.
\item Los hijos de mis tíos son mis primos, y los hijos de mis \esp{\ldots\ldots} son mis sobrinos.
\item Los \esp{\ldots\ldots} de mis padres son mis hermanos y yo.
\end{enumerate}
\end{exercice}

\begin{exercice}[Conjugaison à compléter]
Conjugue les verbes entre parenthèses au présent de l'indicatif.
\begin{enumerate}
\item Los jóvenes \esp{\ldots\ldots} (pasar) mucho tiempo en las redes sociales.
\item Mis padres \esp{\ldots\ldots} (querer) que yo estudie más.
\item Tú \esp{\ldots\ldots} (deber) respetar a tus abuelos.
\item Nosotros \esp{\ldots\ldots} (tener) que dialogar en familia.
\item Los conflictos \esp{\ldots\ldots} (ser) frecuentes entre generaciones.
\end{enumerate}
\end{exercice}

\courssub{Je m'entraîne}

\begin{exercice}[Appariement : les types de famille]
Associe chaque type de famille (colonne A) à sa définition (colonne B), puis dessine un arbre généalogique de ta famille en espagnol avec au moins six membres (\esp{abuelo, madre, tío, primo, hermana, yo}\ldots).

\begin{center}
\begin{tabularx}{\linewidth}{|X|X|}
\hline
\rowcolor{popblueL} \textbf{A — Tipo de familia} & \textbf{B — Definición} \\
\hline
1. Familia nuclear & a. Un solo progenitor vive con sus hijos. \\
\hline
2. Familia monoparental & b. Padres divorciados forman nuevas parejas con hijos. \\
\hline
3. Familia extensa & c. Padre, madre e hijos viven juntos. \\
\hline
4. Familia reconstituida & d. Viven juntos abuelos, padres, hijos, tíos y primos. \\
\hline
\end{tabularx}
\end{center}
\end{exercice}

\begin{textebox}[La brecha generacional]
En muchas familias de Yaoundé, los jóvenes y los abuelos no se entienden bien \esp{\ldots\ldots} tienen valores diferentes. \esp{\ldots\ldots} la urbanización, muchos jóvenes viven lejos del pueblo y olvidan las tradiciones. \esp{\ldots\ldots}, los ancianos se sienten solos. \esp{\ldots\ldots} las nuevas tecnologías, los nietos pueden llamar a sus abuelos por vídeo. Si quieres mejorar la relación con tu familia, \esp{\ldots\ldots} escuchar más a los mayores. En conclusión, \esp{\ldots\ldots} dialogar con respeto para construir una familia unida.
\end{textebox}

\begin{exercice}[Texte à trous : la brecha generacional]
Complète le texte ci-dessus avec les connecteurs de la liste : \esp{porque} — \esp{por eso} — \esp{gracias a} — \esp{a causa de} — \esp{deberías} — \esp{hay que}. Chaque connecteur n'est utilisé qu'une seule fois.
\end{exercice}

\begin{exercice}[Transformation de phrases]
Transforme chaque phrase selon la consigne entre parenthèses.
\begin{enumerate}
\item \esp{Los jóvenes usan mucho el móvil. Los padres se enfadan.} (relie les deux phrases avec \esp{por eso})
\item \esp{Hay conflictos porque no hay diálogo.} (remplace \esp{porque} par \esp{a causa de} en transformant la phrase)
\item \esp{Escucha a tus abuelos.} (transforme l'impératif en conseil avec \esp{deberías})
\item \esp{Los padres trabajan mucho. No pasan tiempo con sus hijos.} (relie les deux phrases avec \esp{como})
\item \esp{Es necesario respetar a los mayores.} (réécris la phrase avec \esp{hay que})
\end{enumerate}
\end{exercice}

\begin{exercice}[Traduction : thème]
Traduis en espagnol les phrases suivantes.
\begin{enumerate}
\item Les jeunes discutent avec leurs parents parce qu'ils passent trop de temps sur le téléphone.
\item C'est pourquoi il faut dialoguer en famille.
\item Tu devrais écouter tes grands-parents : ils ont beaucoup d'expérience.
\item À cause des réseaux sociaux, les enfants parlent moins avec leurs parents.
\end{enumerate}
\end{exercice}

\courssub{Je me prépare au Bac}

\begin{textebox}[La familia cambia]
La familia africana y española ha cambiado mucho en los últimos años. Antes, la familia extensa era la norma : abuelos, padres, hijos, tíos y primos vivían juntos en la misma casa o en el mismo barrio. Hoy, sobre todo en las ciudades como Douala o Madrid, la familia nuclear es más frecuente : solo los padres y sus hijos comparten la vivienda.

Este cambio tiene consecuencias importantes. Los abuelos, que antes cuidaban a los nietos y transmitían las tradiciones, viven ahora lejos, a menudo en el pueblo. A causa de la urbanización y del trabajo, los padres pasan poco tiempo con sus hijos. Como los abuelos viven lejos, los nietos pierden el contacto con las tradiciones. Además, las nuevas tecnologías crean una brecha generacional : los jóvenes dominan el móvil y las redes sociales, mientras que los mayores no entienden ese mundo digital.

Los conflictos entre generaciones son frecuentes : los padres quieren autoridad y respeto, y los hijos piden más libertad. Por eso, muchos especialistas recomiendan el diálogo familiar. Hay que escuchar a los demás, comprender sus razones y buscar soluciones juntos. La familia cambia, pero sigue siendo el centro de la vida afectiva y social.
\end{textebox}

\textbf{Glossaire} : \esp{la vivienda} = le logement ; \esp{cuidar a} = s'occuper de ; \esp{la brecha} = le fossé ; \esp{mientras que} = tandis que ; \esp{pedir} = demander ; \esp{seguir siendo} = rester, continuer à être.

\begin{exercice}[Compréhension écrite type Bac]
Lis attentivement le texte \esp{La familia cambia} ci-dessus, puis réponds aux questions en espagnol avec des phrases complètes.

\textbf{Questions de compréhension}
\begin{enumerate}
\item Indique si la phrase suivante est vraie ou fausse et justifie ta réponse avec une citation du texte : \esp{Hoy, la familia extensa es la más frecuente en las ciudades.}
\item ¿Qué dos causas explican que los padres pasen poco tiempo con sus hijos ?
\item ¿Qué es la \esp{brecha generacional} según el texto ?
\item ¿Qué solución proponen los especialistas para los conflictos familiares ?
\item ¿Cuál es la idea principal del último párrafo ?
\end{enumerate}

\textbf{Questions de langue}
\begin{enumerate}
\item Trouve dans le texte deux connecteurs qui expriment la cause et un connecteur qui exprime la conséquence.
\item Réécris la phrase \esp{Los jóvenes dominan el móvil} au passé composé (\esp{pretérito perfecto}).
\end{enumerate}
\end{exercice}

\begin{exercice}[Thème et version type Bac]
\textbf{A. Thème.} Traduis en espagnol le texte suivant.

\begin{quote}
Dans ma famille, il y a souvent des conflits entre les générations. Mon grand-père pense que les jeunes ne respectent plus les traditions parce qu'ils passent tout leur temps sur Internet. C'est pourquoi il est triste. Je crois qu'il faut organiser des réunions de famille pour dialoguer. Tu devrais parler avec lui : il a beaucoup de choses à raconter.
\end{quote}

\textbf{B. Version.} Traduis en français le texte \esp{El diálogo en familia} ci-dessous.
\end{exercice}

\begin{textebox}[El diálogo en familia]
En muchas familias hay tensiones porque los padres y los hijos no se comprenden. A causa del trabajo y del estrés, los padres llegan cansados a casa y no hablan con sus hijos. Por eso, los jóvenes se refugian en el móvil y en las redes sociales. Para mejorar la situación, hay que dedicar tiempo a la familia : comer juntos, conversar y escuchar con paciencia. El diálogo es la clave de la armonía familiar.
\end{textebox}

\begin{exercice}[Expression écrite type Bac]
Rédige en espagnol un texte de \textbf{80 à 100 mots} sur le sujet suivant :

\begin{center}
\esp{En tu familia, ¿hay conflictos entre generaciones? ¿Por qué? Propón dos soluciones.}
\end{center}

\textbf{Consignes :}
\begin{enumerate}
\item Présente brièvement ta famille et la situation (2 ou 3 phrases).
\item Explique au moins une cause des conflits en utilisant \esp{porque} ou \esp{a causa de}.
\item Propose deux solutions en utilisant \esp{hay que}, \esp{deberíamos} ou \esp{por eso}.
\item Soigne l'orthographe, les accents et la conjugaison au présent.
\end{enumerate}
\end{exercice}



\newpage

\coursec{Corrigés des exercices}

\begin{corrige}[Exercice 1 — QCM : la famille et les générations]
\begin{enumerate}
\item \textbf{b. nuclear} : la famille nucléaire est formée uniquement du père, de la mère et de leurs enfants. Les autres options : \esp{extensa} = famille élargie (plusieurs générations) ; \esp{reconstituida} = famille recomposée ; \esp{monoparental} = un seul parent.
\item \textbf{b.} le fossé entre les générations : \esp{brecha} = « fossé, écart », \esp{generacional} = « entre les générations ».
\item \textbf{c.} la cause : \esp{porque} introduit la raison de la dispute (« parce qu'ils passent beaucoup de temps sur le portable »).
\item \textbf{a.} \esp{Deberías escuchar a tus padres.} « Tu devrais écouter tes parents. » Le conditionnel \esp{deberías} (verbe \esp{deber}) exprime le conseil ; les autres phrases sont au présent ou au passé et n'expriment ni conseil ni recommandation.
\end{enumerate}
\end{corrige}

\begin{corrige}[Exercice 2 — Vrai ou faux ?]
\begin{enumerate}
\item \textbf{V} : \esp{Una familia monoparental tiene un solo progenitor : el padre o la madre.} (« Une famille monoparentale n'a qu'un seul parent : le père ou la mère. »)
\item \textbf{F} : \esp{La urbanización influye mucho : los jóvenes viven lejos de sus abuelos.} (« L'urbanisation influence beaucoup les relations : les jeunes vivent loin de leurs grands-parents. »)
\item \textbf{V} : \esp{Los abuelos transmiten las tradiciones, los cuentos y los valores a los nietos.} (« Les grands-parents transmettent les traditions, les contes et les valeurs aux petits-enfants. »)
\item \textbf{V} : \esp{El diálogo permite comprender las razones de los demás y buscar soluciones.} (« Le dialogue permet de comprendre les raisons des autres et de chercher des solutions. »)
\end{enumerate}
Point de vigilance : \esp{transmitir} est régulier au présent (\esp{transmiten}) ; ne pas confondre \esp{los demás} (« les autres ») et \esp{demasiado} (« trop »).
\end{corrige}

\begin{corrige}[Exercice 3 — Vocabulaire : les membres de la famille]
\begin{enumerate}
\item Los \esp{abuelos} de mis \esp{padres} son mis bisabuelos. (« Les grands-parents de mes parents sont mes arrière-grands-parents. »)
\item Mis \esp{padres} son los hijos de mis abuelos. (« Mes parents sont les enfants de mes grands-parents. »)
\item Los hijos de mis \esp{hermanos} son mis sobrinos. (« Les enfants de mes frères et sœurs sont mes neveux et nièces. »)
\item Los \esp{hijos} de mis padres son mis hermanos y yo. (« Les enfants de mes parents sont mes frères et sœurs et moi. »)
\end{enumerate}
Le mot non utilisé est \esp{nietos} (« petits-enfants »). Attention au faux ami : \esp{los padres} = « les parents » (père et mère), pas seulement « les pères ».
\end{corrige}

\begin{corrige}[Exercice 4 — Conjugaison à compléter]
\begin{enumerate}
\item \esp{pasan} : verbe régulier en \esp{-ar}, 3\up{e} personne du pluriel (\esp{pasar} : paso, pasas, pasa, pasamos, pasáis, pasan).
\item \esp{quieren} : verbe à changement radical e $\rightarrow$ ie (\esp{quiero, quieres, quiere, queremos, queréis, quieren}). Attention : après \esp{querer que}, on emploie le subjonctif : \esp{quieren que yo estudie}.
\item \esp{debes} : verbe régulier en \esp{-er}, 2\up{e} personne du singulier (\esp{debo, debes, debe\ldots}).
\item \esp{tenemos} : verbe irrégulier \esp{tener} (\esp{tengo, tienes, tiene, tenemos, tenéis, tienen}) ; \esp{tener que} + infinitif = « devoir, il faut ».
\item \esp{son} : verbe irrégulier \esp{ser}, 3\up{e} personne du pluriel (\esp{soy, eres, es, somos, sois, son}).
\end{enumerate}
\end{corrige}

\begin{corrige}[Exercice 5 — Appariement : les types de famille]
1 — c : \esp{familia nuclear} = père, mère et enfants vivant ensemble. 2 — a : \esp{familia monoparental} = un seul parent vit avec ses enfants. 3 — d : \esp{familia extensa} = grands-parents, parents, enfants, oncles, tantes et cousins sous le même toit. 4 — b : \esp{familia reconstituida} = parents divorcés qui forment de nouveaux couples avec des enfants.

Pour l'arbre généalogique, vérifier : l'accord du possessif (\esp{mi madre, mis abuelos}), le vocabulaire (\esp{abuelo/abuela, tío/tía, primo/prima, hermano/hermana}) et la présentation (\esp{Este es mi abuelo ; se llama\ldots}).
\end{corrige}

\begin{corrige}[Exercice 6 — Texte à trous : la brecha generacional]
\begin{enumerate}
\item \esp{porque} : exprime la cause (« parce qu'ils ont des valeurs différentes »), suivi d'un verbe conjugué.
\item \esp{A causa de} : exprime une cause défavorable (l'urbanisation éloigne les jeunes du village), suivi d'un nom.
\item \esp{Por eso} : exprime la conséquence (« c'est pourquoi les anciens se sentent seuls »).
\item \esp{Gracias a} : exprime une cause favorable (« grâce aux nouvelles technologies »).
\item \esp{deberías} : conseil personnel à la 2\up{e} personne du singulier (« tu devrais écouter davantage les aînés »).
\item \esp{hay que} : obligation générale et impersonnelle (« il faut dialoguer avec respect »).
\end{enumerate}
Astuce à retenir : \esp{gracias a} introduit une cause positive, \esp{a causa de} une cause négative ; \esp{porque} est suivi d'une phrase complète, \esp{a causa de} et \esp{gracias a} d'un groupe nominal.
\end{corrige}

\begin{corrige}[Exercice 7 — Transformation de phrases]
\begin{enumerate}
\item \esp{Los jóvenes usan mucho el móvil, por eso los padres se enfadan.} (« Les jeunes utilisent beaucoup le portable, c'est pourquoi les parents se fâchent. ») \esp{Por eso} relie la cause à la conséquence.
\item \esp{Hay conflictos a causa de la falta de diálogo.} Point de vigilance : \esp{a causa de} est suivi d'un \textbf{nom}, pas d'une phrase verbale ; \esp{no hay diálogo} devient donc \esp{la falta de diálogo} (« le manque de dialogue »).
\item \esp{Deberías escuchar a tus abuelos.} (« Tu devrais écouter tes grands-parents. ») Le conditionnel \esp{deberías} adoucit l'impératif.
\item \esp{Como los padres trabajan mucho, no pasan tiempo con sus hijos.} \esp{Como} = « comme, puisque » ; il se place en tête de phrase. Ne jamais le mettre entre les deux propositions dans ce sens.
\item \esp{Hay que respetar a los mayores.} (« Il faut respecter les aînés. ») \esp{Hay que} + infinitif = obligation impersonnelle.
\end{enumerate}
\end{corrige}

\begin{corrige}[Exercice 8 — Traduction : thème]
\begin{enumerate}
\item \esp{Los jóvenes discuten con sus padres porque pasan demasiado tiempo en el teléfono.} Point de vigilance : « trop de » se traduit \esp{demasiado}, jamais \esp{muy} (« très »).
\item \esp{Por eso hay que dialogar en familia.} « C'est pourquoi » = \esp{por eso} ; « il faut » = \esp{hay que} + infinitif.
\item \esp{Deberías escuchar a tus abuelos : tienen mucha experiencia.} Point de vigilance : la préposition \esp{a} personnelle est obligatoire devant la personne : \esp{escuchar a alguien}. \esp{Experiencia} est féminin : \esp{mucha experiencia}.
\item \esp{A causa de las redes sociales, los niños hablan menos con sus padres.} « À cause de » = \esp{a causa de} + nom ; « les réseaux sociaux » = \esp{las redes sociales}.
\end{enumerate}
\end{corrige}

\begin{corrige}[Exercice 9 — Compréhension écrite type Bac]
\textbf{Questions de compréhension}
\begin{enumerate}
\item \textbf{Falsa.} Justification par une citation du texte : \esp{« Hoy, sobre todo en las ciudades como Douala o Madrid, la familia nuclear es más frecuente »}. C'est la famille nucléaire qui domine en ville, pas la famille élargie.
\item \esp{Los padres pasan poco tiempo con sus hijos a causa de la urbanización y del trabajo.} (« À cause de l'urbanisation et du travail. »)
\item \esp{La brecha generacional es la separación entre los jóvenes, que dominan el móvil y las redes sociales, y los mayores, que no entienden ese mundo digital.} (« Le fossé entre les jeunes, qui maîtrisent le portable et les réseaux sociaux, et les aînés, qui ne comprennent pas ce monde numérique. »)
\item \esp{Los especialistas recomiendan el diálogo familiar : escuchar a los demás, comprender sus razones y buscar soluciones juntos.}
\item \esp{La idea principal es que el diálogo es la solución a los conflictos y que la familia, aunque cambia, sigue siendo el centro de la vida afectiva y social.}
\end{enumerate}
\textbf{Questions de langue}
\begin{enumerate}
\item Connecteurs de cause : \esp{a causa de} (\esp{a causa de la urbanización y del trabajo}) et \esp{como} (\esp{Como los abuelos viven lejos\ldots}). Connecteur de conséquence : \esp{por eso}.
\item \esp{Los jóvenes han dominado el móvil.} Le \esp{pretérito perfecto} = auxiliaire \esp{haber} au présent + participe passé : \esp{han dominado} (« ont dominé »).
\end{enumerate}
\textbf{Barème indicatif (sur 10)} : compréhension 6 points (question 1 : 2 points — 1 pour V/F, 1 pour la citation exacte ; questions 2 à 5 : 1 point chacune) ; langue 4 points (2 points pour les trois connecteurs, 2 points pour le passé composé correct : auxiliaire + participe). Toute réponse doit être une phrase complète en espagnol ; une réponse copiée sans justification ne vaut que la moitié des points.
\end{corrige}

\begin{corrige}[Exercice 10 — Thème et version type Bac]
\textbf{A. Thème (proposition de traduction)}

\esp{En mi familia hay a menudo conflictos entre generaciones. Mi abuelo piensa que los jóvenes ya no respetan las tradiciones porque pasan todo su tiempo en Internet. Por eso está triste. Creo que hay que organizar reuniones de familia para dialogar. Deberías hablar con él : tiene muchas cosas que contar.}

Points de vigilance : « ne\ldots plus » = \esp{ya no} ; « c'est pourquoi » = \esp{por eso} ; « il faut » = \esp{hay que} + infinitif ; « tu devrais » = \esp{deberías} + infinitif ; « à raconter » = \esp{que contar} ; « avec lui » = \esp{con él} (accent obligatoire sur \esp{él} = « lui », sans accent \esp{el} = article).

\textbf{B. Version (proposition de traduction)}

« Dans beaucoup de familles, il y a des tensions parce que les parents et les enfants ne se comprennent pas. À cause du travail et du stress, les parents rentrent fatigués à la maison et ne parlent pas avec leurs enfants. C'est pourquoi les jeunes se réfugient dans le téléphone portable et les réseaux sociaux. Pour améliorer la situation, il faut consacrer du temps à la famille : manger ensemble, converser et écouter avec patience. Le dialogue est la clé de l'harmonie familiale. »

Points de vigilance : \esp{se refugian} = « se réfugient » (verbe pronominal, ne pas oublier le pronom réfléchi en français) ; \esp{la clave} = « la clé » (nom féminin) ; \esp{para mejorar} = « pour améliorer » (expression du but) ; \esp{llegan cansados} = « rentrent fatigués » (accord de l'adjectif au pluriel masculin).

\textbf{Barème indicatif (sur 10)} : thème 5 points (1 point par phrase, dont 0,5 pour le connecteur correctement traduit) ; version 5 points (fidélité du sens 3 points, qualité du français 2 points). Toute faute d'accent modifiant le sens (\esp{el/él}, \esp{si/sí}) est pénalisée.
\end{corrige}

\begin{corrige}[Exercice 11 — Expression écrite type Bac]
\textbf{Exemple de production (environ 95 mots) :}

\esp{En mi familia somos seis personas : mis abuelos, mis padres, mi hermana y yo. A veces hay conflictos entre generaciones porque los jóvenes pasan mucho tiempo en el móvil y los mayores no comprenden las nuevas tecnologías. Además, a causa de la urbanización, vivimos lejos del pueblo y olvidamos las tradiciones. Para mejorar la situación, hay que dialogar con respeto y escuchar a los demás. Deberíamos organizar comidas de familia cada domingo para conversar juntos. Por eso, creo que el diálogo es la mejor solución : la familia cambia, pero sigue siendo lo más importante.}

Traduction : « Dans ma famille, nous sommes six personnes : mes grands-parents, mes parents, ma sœur et moi. Parfois, il y a des conflits entre les générations parce que les jeunes passent beaucoup de temps sur le portable et que les aînés ne comprennent pas les nouvelles technologies. De plus, à cause de l'urbanisation, nous vivons loin du village et nous oublions les traditions. Pour améliorer la situation, il faut dialoguer avec respect et écouter les autres. Nous devrions organiser des repas de famille chaque dimanche pour discuter ensemble. C'est pourquoi je crois que le dialogue est la meilleure solution : la famille change, mais elle reste la chose la plus importante. »

\textbf{Barème indicatif (sur 10)} : respect de la tâche et du nombre de mots (80--100) : 2 points ; au moins un connecteur de cause (\esp{porque}, \esp{a causa de}, \esp{como}) : 2 points ; deux expressions du conseil ou de la conséquence (\esp{hay que}, \esp{deberíamos}, \esp{por eso}) : 2 points ; correction grammaticale (présent de l'indicatif, accords) : 2 points ; orthographe et accents (\esp{tecnologías}, \esp{solución}, \esp{domingo}) : 2 points.

Points de vigilance : ne pas écrire \esp{muy tiempo} pour « beaucoup de temps » (\esp{mucho tiempo}) ; accorder \esp{nuevas tecnologías} ; ne pas oublier \esp{a} devant les personnes (\esp{escuchar a los demás}).
\end{corrige}

\newpage

\coursec{Fiche bilan}

\begin{popfigure}
\begin{tikzpicture}[mindmap, grow cyclic, every node/.style={concept, text=white, font=\small},
  root concept/.append style={concept color=popblue, font=\bfseries},
  level 1 concept/.append style={level distance=3.4cm, sibling angle=51, font=\scriptsize},
  level 2 concept/.append style={level distance=2.2cm, font=\tiny}]
\node[root concept]{La familia: evolución y conflictos}
  child[concept color=popgreen]{ node {Tipos de familia}
    child { node {nuclear} }
    child { node {monoparental} }
    child { node {recompuesta} }
  }
  child[concept color=poporange]{ node {Lexique}
    child { node {abuelos, padres, hijos} }
    child { node {brecha generacional} }
  }
  child[concept color=poppurple]{ node {Causes des conflits}
    child { node {NTIC, valeurs} }
    child { node {urbanización} }
  }
  child[concept color=popgold]{ node {La cause}
    child { node {porque, ya que, puesto que, a causa de} }
  }
  child[concept color=poppink]{ node {La conséquence}
    child { node {por eso, así que, entonces, por lo tanto} }
  }
  child[concept color=popblue!70]{ node {Le conseil}
    child { node {deberías + inf.} }
    child { node {hay que + inf.} }
  }
  child[concept color=popgreen!70!popblue]{ node {Méthode}
    child { node {commentaire de texte} }
  };
\end{tikzpicture}
\legende{Carte mentale de la leçon : la famille, ses évolutions et les conflits entre générations.}
\end{popfigure}

\begin{bilan}
\end{bilan}

\courssub{Lexique clé}
\begin{itemize}
  \item \esp{la familia nuclear / monoparental / extensa / recompuesta} : la famille nucléaire / monoparentale / élargie / recomposée.
  \item \esp{los abuelos, los padres, los hijos, los nietos} : les grands-parents, les parents, les enfants, les petits-enfants.
  \item \esp{la brecha generacional} : le fossé des générations.
  \item \esp{la autoridad, el respeto, el conflicto, el diálogo} : l'autorité, le respect, le conflit, le dialogue.
  \item \esp{las nuevas tecnologías (NTIC), los valores, la urbanización} : les nouvelles technologies, les valeurs, l'urbanisation.
  \item \esp{entenderse, escuchar, dialogar, respetar} : se comprendre, écouter, dialoguer, respecter.
\end{itemize}

\courssub{Grammaire à connaître par cœur}
\begin{center}
\begin{tabularx}{\linewidth}{|l|X|X|}
\hline
\rowcolor{popblueL} Notion & Structures & Exemple \\
\hline
La cause & \esp{porque, ya que, puesto que, como, a causa de + nom} & \esp{Hay conflictos porque los jóvenes usan mucho el móvil.} \\
\hline
La conséquence & \esp{por eso, así que, entonces, por lo tanto, de modo que} & \esp{No se entienden, por eso discuten.} \\
\hline
Le conseil & \esp{deberías / debería + infinitif ; hay que + infinitif ; te aconsejo que + subjonctif} & \esp{Deberías escuchar a tus padres. Hay que dialogar.} \\
\hline
\end{tabularx}
\end{center}

\courssub{Méthodes express}
\begin{itemize}
  \item \textbf{Comprendre un texte} : lire le titre, repérer le thème, souligner le lexique de la famille, relever causes et solutions.
  \item \textbf{Commenter} : introduction (présenter le texte et le problème), développement (causes, conséquences, solutions), conclusion (avis personnel).
  \item \textbf{Traduire} : attention aux faux amis et aux connecteurs ; \esp{porque} répond à \esp{¿por qué?} ; \esp{por eso} introduit la conséquence.
  \item \textbf{Rédiger} : proposer le dialogue familial avec \esp{deberías}, \esp{hay que}, \esp{es importante que + subjonctif}.
\end{itemize}


\begin{aretenir}[Checklist avant le Bac]
\begin{itemize}
  \item $\square$ Je connais le vocabulaire des types de famille et des relations entre générations.
  \item $\square$ Je sais définir \esp{la brecha generacional} et donner des exemples.
  \item $\square$ Je cite au moins trois causes de conflits intergénérationnels (NTIC, valeurs, urbanisation).
  \item $\square$ J'emploie correctement \esp{porque}, \esp{ya que}, \esp{puesto que} et \esp{a causa de} pour la cause.
  \item $\square$ Je distingue \esp{porque} (cause) et \esp{por eso} (conséquence) sans les confondre.
  \item $\square$ J'utilise \esp{así que}, \esp{entonces}, \esp{por lo tanto} pour exprimer la conséquence.
  \item $\square$ Je conjugue \esp{deber} : \esp{deberías + infinitif} pour conseiller.
  \item $\square$ J'emploie \esp{hay que + infinitif} pour un conseil général et impersonnel.
  \item $\square$ Je sais que \esp{te aconsejo que} est suivi du subjonctif.
  \item $\square$ Je respecte les accents espagnols : \esp{¿por qué?}, \esp{porque}, \esp{atención}, \esp{diálogo}.
  \item $\square$ Je sais structurer un commentaire de texte en trois parties.
  \item $\square$ Je peux rédiger cinq phrases proposant le dialogue familial.
\end{itemize}
\end{aretenir}

\findecours

\end{document}
